\documentclass[11pt,a4paper,openany]{book}
\usepackage[utf8]{inputenc}
\usepackage[T1]{fontenc}
\usepackage[french]{babel}
\usepackage{lmodern}
\usepackage{microtype}
\usepackage[a4paper,margin=2.35cm,headheight=15pt]{geometry}
\usepackage{amsmath,amssymb,amsthm,mathtools,mathrsfs,bm}
\usepackage{booktabs,longtable,tabularx,array,multirow}
\usepackage{tikz,tikz-cd,float}
\usetikzlibrary{arrows.meta,positioning,calc,fit,shapes.geometric}
\usepackage{enumitem}
\usepackage{hyperref}
\usepackage{url}
\usepackage{fancyvrb}
\usepackage{fancyhdr}
\usepackage{xspace}
\DeclareUnicodeCharacter{2013}{--}
\DeclareUnicodeCharacter{2014}{---}
\DeclareUnicodeCharacter{2192}{\ensuremath{\to}}
\DeclareUnicodeCharacter{21D2}{\ensuremath{\Rightarrow}}
\DeclareUnicodeCharacter{2264}{\ensuremath{\leq}}
\DeclareUnicodeCharacter{2265}{\ensuremath{\geq}}
\DeclareUnicodeCharacter{227C}{\ensuremath{\preceq}}
\DeclareUnicodeCharacter{2286}{\ensuremath{\subseteq}}
\DeclareUnicodeCharacter{2260}{\ensuremath{\neq}}
\DeclareUnicodeCharacter{2261}{\ensuremath{\equiv}}
\DeclareUnicodeCharacter{21A6}{\ensuremath{\mapsto}}
\DeclareUnicodeCharacter{03BB}{\ensuremath{\lambda}}
\DeclareUnicodeCharacter{00D7}{\ensuremath{\times}}
\setlength{\parindent}{0pt}
\setlength{\parskip}{6pt}
\setlength{\headheight}{25.5pt}
\setlength{\emergencystretch}{2em}
\sloppy
\setlist{itemsep=3pt,topsep=4pt}
\pagestyle{fancy}
\fancyhf{}
\fancyhead[LE,RO]{\thepage}
\fancyhead[RE]{\nouppercase{\leftmark}}
\fancyhead[LO]{\nouppercase{\rightmark}}
\hypersetup{hypertexnames=false,colorlinks=true,linkcolor=black,urlcolor=black,citecolor=black,pdftitle={Théorie des empreintes -- reconstruction order 5, livre consolidé},pdfauthor={Bensalem Nassim}}
\newtheorem{definition}{Définition}[chapter]
\newtheorem{axiom}[definition]{Axiome / postulat}
\newtheorem{hypothesis}[definition]{Hypothèse}
\newtheorem{proposition}[definition]{Proposition}
\newtheorem{lemma}[definition]{Lemme}
\newtheorem{theorem}[definition]{Théorème}
\newtheorem{corollary}[definition]{Corollaire}
\newtheorem{conjecture}[definition]{Conjecture}
\newtheorem{remark}[definition]{Remarque}
\newtheorem{warning}[definition]{Avertissement}
\newtheorem{noGo}[definition]{No-go}
\DeclareMathOperator{\Aff}{Aff}
\newcommand{\R}{\mathbb{R}}
\newcommand{\Rp}{\mathbb{R}_{+}}
\newcommand{\X}{\mathcal X}
\newcommand{\Pset}{\mathcal P}
\newcommand{\Fset}{\mathcal F}
\newcommand{\finger}[2]{I_{#1}^{#2}}
\newcommand{\cl}{\operatorname{Cl}}
\newcommand{\Adm}{\operatorname{Adm}}
\newcommand{\Def}{\operatorname{Def}}
\newcommand{\Reach}{\operatorname{Reach}}
\newcommand{\supp}{\operatorname{Supp}}
\newcommand{\Powfin}{\mathcal P_{\!fin}}
\newcommand{\status}[1]{\textbf{#1}}
\newcolumntype{Y}{>{\raggedright\arraybackslash}X}
\newcolumntype{L}[1]{>{\raggedright\arraybackslash}p{#1}}
\begin{document}
\begin{titlepage}
\centering
{\Large Projet de reconstruction scientifique}\\[1.0cm]
{\Huge\bfseries Théorie des empreintes informationnelles}\\[0.8cm]
{\Large Reconstruction critique, re-développement et livre consolidé — order 5}\\[1.2cm]
{\Large Bensalem Nassim}\\[0.6cm]
\vfill
\textit{Corpus source fondamental : sources\_th\_empreintes\_v5\_1.zip}\par
\textit{La structure du livre suit la logique interne de la théorie ; la traçabilité documentaire est reportée en annexes et registres.}\par
\today
\end{titlepage}
\frontmatter
\tableofcontents
\mainmatter
\part{Fondations de l'empreinte}

\chapter{Objet, portée et méthode scientifique}
\section{Objet}
La théorie reconstruite ici étudie une idée simple mais structurante : un objet n'est pas décrit par une essence supposée, mais par les réponses obtenues lorsqu'un système d'observation lui applique des sondes. Une \emph{empreinte} est donc une représentation induite par un protocole d'observation. Les distinctions entre objets sont celles que le protocole sait effectivement voir.

Cette reconstruction ne prend pas les versions tardives comme automatiquement vraies. La chronologie sert à suivre les modifications ; la validité est décidée par les définitions, les preuves, les contre-exemples, les vérifications et les limites explicites. La branche DD\_Dev\_extension est ainsi intégrée comme programme de recherche, et non comme autorité supérieure aux fondements.

\section{Statuts utilisés}
Tout énoncé est classé selon l'un des statuts suivants : définition ; axiome/postulat ; hypothèse ; conjecture ; proposition/lemme/théorème ; démonstration ; corollaire ; résultat expérimental ; observation ; interprétation ; piste de recherche ; contre-exemple ; limitation. Le document principal n'utilise comme résultats théoriques que les énoncés qui ont reçu un argument suffisant dans le cadre déclaré. Le document d'essai conserve les tentatives non abouties.

\section{Corpus et règle de fidélité}
Le document fondateur est le fichier de socle des seuils et de l'empreinte, référencé dans la traçabilité. Les D\_Base sont reconstruits comme développements de ce socle. Les D\_Dev, DD\_Dev et DD\_Dev\_extension sont traités comme des extensions successives, tandis que les rapports et expériences sont des éléments de preuve expérimentale ou d'historique du développement. Le présent document distingue en outre le corpus courant v5.1 des anciennes archives utilisées par certaines versions expérimentales.

\section{Architecture logique}
La chaîne générale retenue est
\[
\begin{gathered}
\text{espace d'objets}\to\text{observation}\to\text{empreinte}\to\text{fibres/quotient}\\
\text{fibres/quotient}\to\text{géométrie des seuils}\to\text{dynamique interne}\\
\text{dynamique interne}\to\text{apprentissage / recherche}\to\text{exécution / transfert}.
\end{gathered}
\]
Cette chaîne ne doit pas être lue comme une hiérarchie de « puissance » : chaque flèche correspond à une construction déterminée, avec ses hypothèses propres.

\chapter{Observation, indiscernabilité, fibres et quotient}
\section{Espaces, observables et sondes}
Soient \(\Pset\) un ensemble de contextes, \(\X\) un ensemble d'objets observables et \(\Fset\) une famille de sondes. Une application d'observation est donnée par
\[
\Phi:\Pset\times\X\times\Fset\longrightarrow\R.
\]
Pour un contexte \(P\), un objet \(x\) et une famille finie de sondes \(\Fset_K=\{f_1,\dots,f_K\}\), l'empreinte informationnelle est
\[
I_\Phi^P(x)=
\bigl(\Phi(P,x,f_1),\dots,\Phi(P,x,f_K)\bigr).
\]
Pour une famille quelconque \(\mathcal G\subseteq\Fset\), on écrit plus généralement
\[
I_{\Phi,\mathcal G}^{P}:\X\to\R^{\mathcal G},\qquad
x\mapsto\bigl(f\mapsto\Phi(P,x,f)\bigr).
\]

\section{Échelle d'observation}
Une échelle est un couple \((\Lambda,\preceq)\) muni d'une famille croissante de sondes \((\Fset_\lambda)_{\lambda\in\Lambda}\), telle que
\[
\lambda\preceq\mu\Longrightarrow
\Fset_\lambda\subseteq\Fset_\mu.
\]
Une intensité monotone \(\iota:\Lambda\to\Rp\) paramètre la résolution. La capacité brute est
\[
C(\lambda)=|\Fset_\lambda|
\]
dans le cas fini. L'empreinte à l'échelle \(\lambda\) est
\[
I_{\Phi,\lambda}^{P}(x)=I_{\Phi,\Fset_\lambda}^{P}(x).
\]

\section{Indiscernabilité, fibres et quotient}
On définit
\[
x\sim_{\Phi,P,\lambda} y
\iff
I_{\Phi,\lambda}^{P}(x)=I_{\Phi,\lambda}^{P}(y).
\]
Il s'agit d'une relation d'équivalence. Les classes sont les fibres observationnelles. L'application
\[
q_{P,\lambda}:\X\to\X/{\sim_{\Phi,P,\lambda}}
\]
est le quotient canonique, tandis que l'image observable
\[
\mathcal L_{\lambda,P}=I_{\Phi,\lambda}^{P}(\X)
\]
représente les valeurs effectivement observées. L'empreinte factorise canoniquement : l'information pertinente est la classe observée, non l'identité brute de l'objet.

\begin{proposition}[Raffinement monotone]
Si \(\lambda\preceq\mu\), alors \(\sim_{\Phi,P,\mu}\) est plus fine que \(\sim_{\Phi,P,\lambda}\). Autrement dit,
\[
I_{\Phi,\mu}^{P}(x)=I_{\Phi,\mu}^{P}(y)
\Longrightarrow
I_{\Phi,\lambda}^{P}(x)=I_{\Phi,\lambda}^{P}(y).
\]
\end{proposition}
\begin{proof}
Les sondes de \(\Fset_\lambda\) sont incluses dans \(\Fset_\mu\). L'égalité de toutes les composantes de la seconde empreinte implique donc l'égalité des composantes de la première.\qed
\end{proof}

\chapter{Réduction des sondes, fermeture et factorisation}
\section{Équivalence observationnelle des sondes}
À contexte et échelle fixés, deux sondes sont équivalentes lorsque
\[
f\equiv_{P,\lambda}g
\iff
\forall x\in\X,\quad
\Phi(P,x,f)=\Phi(P,x,g).
\]
Cette relation permet de retirer les duplications de probes sans changer les fibres.

\section{Fermeture observationnelle}
Pour \(\mathcal G\subseteq\Fset\), on définit la fermeture
\[
\cl_{P,\lambda}(\mathcal G)
=
\left\{f\in\Fset_\lambda:\exists h:\mathcal L_{\mathcal G}\to\R,
\ \Phi(P,x,f)=h\bigl(I_{\Phi,\mathcal G}^{P}(x)\bigr)\ \forall x\right\}.
\]
Cette famille contient exactement les sondes dont la réponse est déterminée par l'empreinte déjà produite.

\begin{theorem}[Propriétés de la fermeture]
La fermeture observationnelle est extensive, monotone et idempotente :
\[
\mathcal G\subseteq\cl(\mathcal G),\qquad
\mathcal G\subseteq\mathcal H\Rightarrow\cl(\mathcal G)\subseteq\cl(\mathcal H),
\qquad
\cl(\cl(\mathcal G))=\cl(\mathcal G).
\]
De plus, elle possède les mêmes fibres que \(\mathcal G\) et est maximale parmi les familles de sondes conservant ces fibres.
\end{theorem}
\begin{proof}
L'extensivité découle du fait qu'une coordonnée est déterminée par elle-même. La monotonie suit de l'inclusion des informations disponibles. Pour l'idempotence, si une sonde est déjà factorisée à travers l'empreinte de \(\mathcal G\), l'ajouter ne crée aucune nouvelle distinction ; toutes les nouvelles coordonnées sont elles-mêmes fonctions des anciennes. La propriété des fibres suit directement de cette factorisation.\qed
\end{proof}

\section{Théorème de factorisation de l'empreinte}
Soit \(I=I_{\Phi,\mathcal G}^{P}\) et \(\mathcal L=I(\X)\). Pour une application \(A\), les conditions suivantes sont équivalentes : \(A\) est constante sur les fibres de \(I\) ; il existe une application \(\bar A:\mathcal L\to Y\) telle que \(A=\bar A\circ I\) ; il existe une application \(\hat A:\X/{\sim_I}\to Y\) telle que \(A=\hat A\circ q\). Toute deux représentations ayant les mêmes fibres sont canoniquement bijectives sur leur image.

\begin{proof}
Si \(A\) est constante sur les fibres, on peut définir \(\bar A(I(x)):=A(x)\). La constance rend la définition indépendante du représentant. Réciproquement, une composition avec \(I\) est constante sur ses fibres. Le passage au quotient est identique avec \(q(x)=[x]\). Si deux représentations ont les mêmes fibres, chacune est constante sur les fibres de l'autre et induit deux applications inverses entre leurs images.\qed\end{proof}

\part{Seuils et géométrie de résolution}

\chapter{Seuils, hiérarchies et géométrie observationnelle}
\section{Seuil de séparation d'une paire}
Pour deux objets, on définit le seuil d'indiscernabilité
\[
\rho_P(x,y)
=\sup\{\iota(\lambda):x\sim_{\Phi,P,\lambda}y\},
\]
avec les conventions usuelles nécessaires si l'ensemble est vide ou non borné. Pour séparer tout le domaine à une échelle donnée, on introduit
\[
\rho_P^*(\Phi)=
\inf\left\{\iota(\lambda):I_{\Phi,\lambda}^{P}\text{ est injective}\right\},
\]
\emph{lorsque l'ensemble considéré est non vide}. Cette clause d'existence est indispensable.

\section{Fonction de seuil complète}
Pour une paire non ordonnée \(\{x,y\}\), on pose
\[
\Theta_P(\{x,y\})=\rho_P(x,y).
\]
Le seuil critique \(\Theta_P\) est la borne supérieure du segment initial d'indiscernabilité d'une paire. Il constitue un \emph{squelette de seuil} : il détermine les coupures strictement sous et strictement au-dessus du seuil, mais le statut exact au point critique requiert une donnée de bord lorsque le supremum n'est pas nécessairement atteint. Sous une hypothèse d'atteinte aux niveaux critiques, on retrouve la convention fermee
\[
x\equiv_t^K y \iff K(\{x,y\})\ge t.
\]

\begin{theorem}[Loi du minimum double]
Pour tous \(x,y,z\) deux à deux distincts,
\[
\Theta_P(\{x,z\})\ge
\min\left(\Theta_P(\{x,y\}),\Theta_P(\{y,z\})\right).
\]
\end{theorem}
\begin{proof}
Si \(x\) et \(y\) sont indiscernables jusqu'au niveau \(a\), et \(y\) et \(z\) jusqu'au niveau \(b\), alors jusqu'au niveau \(\min(a,b)\) la relation d'indiscernabilité, étant transitive, identifie aussi \(x\) et \(z\). En prenant les supremums des niveaux communs on obtient l'inégalité.\qed\end{proof}

\begin{theorem}[Caractérisation du squelette hiérarchique]
Soit \(K:\binom{\mathcal X}{2}\to\mathbb R_+\) symétrique. Si \(K\) satisfait la loi du minimum double, les relations
\[
x\equiv_t^K y\iff x=y\ \text{ou}\ K(\{x,y\})\ge t
\]
forment une filtration de partitions totalement ordonnée par raffinement et une structure laminaire finie lorsque \(\mathcal X\) est fini. Réciproquement, une telle filtration munie d'une convention de bord fermee aux seuils critiques fournit un \(K\) par borne supérieure des niveaux d'indiscernabilité. Sans convention de bord, la borne supérieure ne détermine pas à elle seule l'appartenance au niveau critique.
\end{theorem}
\begin{proof}
La transitivité découle de la loi du minimum double et la laminarité de l'emboîtement des partitions. Pour la réciproque, la borne supérieure définit le squelette des seuils ; la convention de bord fermee garantit la correspondance exacte avec les coupes \(K\ge t\). Sans cette convention, deux segments initiaux ayant la même borne supérieure mais des bords différents donnent le même \(K\) et des partitions différentes au point critique.\qed\end{proof}

\section{Ultramétrique dérivée : correction formelle}
La représentation métrique n'est pas primitive. Lorsque \(K>0\), on définit correctement
\[
d_K(x,x)=0,
\qquad
 d_K(x,y)=\exp\bigl(-K(\{x,y\})\bigr)\quad(x\ne y).
\]
\begin{proposition}[Ultramétrique dérivée]
Sous la loi du minimum double, \(d_K\) satisfait
\[
d_K(x,z)\le\max\{d_K(x,y),d_K(y,z)\}.
\]
\end{proposition}
\begin{proof}
Pour \(x,y,z\) distincts, la décroissance de \(u\mapsto e^{-u}\) transforme
\(K(x,z)\ge\min(K(x,y),K(y,z))\) en l'inégalité forte triangulaire. Les cas diagonaux sont immédiats puisque \(d_K(x,x)=0\).\qed\end{proof}

\section{Mesure empirique des seuils}
Pour \(2\leq |\X|<\infty\),
\[
\nu_P=
\frac{1}{\binom{|\X|}{2}}
\sum_{\{x,y\}\subseteq\X}
\delta_{\Theta_P(\{x,y\})}.
\]
La résolution cumulée est
\[
\widetilde{\mathcal R}_P(t)=
\frac{1}{\binom{|\X|}{2}}
\sum_{\{x,y\}}
\mathbf 1_{\{t>\Theta_P(\{x,y\})\}}.
\]
Elle est une fonction en escalier dont les sauts correspondent aux atomes de \(\nu_P\). La reconstruction empirique de la géométrie des seuils est donc équivalente à la connaissance de cette mesure atomique.

\section{Régimes de représentation}
On définit deux bornes critiques avec attention aux conventions de saut :
\[
\alpha_-:=\inf\{t:\widetilde{\mathcal R}_P(t)>0\},
\qquad
\alpha_+:=\sup\{t:\widetilde{\mathcal R}_P(t)<1\}.
\]
Avec la convention stricte \(\mathbf 1_{\{t>\Theta\}}\), les assertions sûres sont \(t<\alpha_-\Rightarrow R(t)=0\) et \(t>\alpha_+\Rightarrow R(t)=1\), tandis que les points exactement égaux aux bornes dépendent des atomes. L'intervalle strict \(\alpha_-<t<\alpha_+\) correspond au régime intermédiaire lorsque celui-ci existe.

\begin{remark}[Lecture en trois régimes]
Le vocabulaire « insuffisant / généralisation observationnelle / terminal » désigne trois zones d'une même filtration. Il ne s'agit pas de trois mécanismes indépendants ni d'un score universel.
\end{remark}

\chapter{Dynamique interne, trajectoires et ordre des observations}
\section{Transitions compatibles}
Une transition \(T:P\leadsto P'\) est dite compatible lorsqu'elle respecte les fibres de l'empreinte selon le protocole déclaré. L'objectif n'est pas d'introduire un temps physique externe, mais de définir un ordre interne à partir de transformations observationnelles compatibles.

Pour deux états d'observation, on dit que \(I\preceq J\) lorsque toutes les distinctions visibles dans \(I\) restent visibles dans \(J\). Les transitions compatibles composent une catégorie de transports observationnels ; dans les versions de base, l'anti-symétrisation du préordre mène à un ordre partiel des régimes.

\begin{proposition}[Composition]
Si \(T_1\) transporte \(I_0\) vers \(I_1\) et \(T_2\) transporte \(I_1\) vers \(I_2\), alors la composition \(T_2\circ T_1\) transporte \(I_0\) vers \(I_2\) avec la même notion de compatibilité.
\end{proposition}
\begin{proof}
La préservation des fibres se compose : si deux points sont indiscernables dans l'état initial, ils restent indiscernables après le premier transport, puis après le second.\qed\end{proof}

\section{Trajectoires et profils}
Une trajectoire est une suite \((I_k)_{k\ge0}\) munie de transports compatibles \(r_{k+1\to k}\). On peut associer un profil de stabilisation, qui enregistre les niveaux où les classes deviennent invariantes. Les résultats des D\_Base permettent alors de raisonner sur des chaînes finies ou infinies, sous réserve de distinguer l'existence formelle d'une chaîne de sa calculabilité effective.

\begin{warning}
Une affirmation de convergence ou de stabilisation sur une chaîne infinie exige des hypothèses explicites. Un calcul fini, même répété sur de nombreux cas, ne constitue pas une preuve d'une propriété de toutes les chaînes.
\end{warning}

\chapter{Identité, tour et réflexion interne}
\section{Identité comme famille cohérente}
Pour un état interne \([\kappa]\), l'identité n'est pas ajoutée comme primitive : elle est la famille cohérente des empreintes visibles le long du passé interne,
\[
\operatorname{Id}_{[\kappa]}(I_\kappa)=
(I_\lambda)_{\lambda\trianglelefteq\kappa},
\]
les composantes devant être compatibles par restriction.

\begin{theorem}[Unicité de l'identité interne]
Une famille cohérente de restrictions est entièrement déterminée par son niveau courant et ses transports de restriction ; deux familles ayant les mêmes données courantes et les mêmes restrictions coïncident.
\end{theorem}
\begin{proof}
Chaque composante passée est la restriction canonique de la représentation courante. Il n'existe donc aucun degré de liberté supplémentaire compatible avec les mêmes restrictions.\qed\end{proof}

\section{Tour et réflexion}
Le \emph{tour} rassemble les descriptions internes déjà présentes en suivant les restrictions cohérentes. La réflexion reconstruit ensuite une trajectoire à partir de cette identité et du tour. Les versions D\_Base établissent une idempotence structurelle,
\[
T(T(I))=T(I),\qquad R(R(\tau))=R(\tau).
\]
Cela signifie qu'une seconde réflexion ne crée pas, par simple répétition, une nouvelle structure canonique si la fermeture observationnelle est inchangée.

\begin{remark}[Portée]
Cette idempotence est un résultat de structure interne. Elle ne constitue pas une preuve que la réflexion produit une capacité comportementale nouvelle. Les extensions expérimentales ultérieures confirment au contraire qu'une nouveauté comportementale nécessite une modification explicite de la frontière de ressources ou du codomaine sémantique.
\end{remark}

\part{Théorie générale de la circulation}

\chapter{Apprentissage, généralisation et méta-observation}
\section{Apprentissage canonique}
Les sources D\_Dev décrivent l'apprentissage comme une évolution de familles latentes compatibles : l'expérience n'est pas une simple accumulation de symboles, mais une transformation de la structure observée, avec un ordre de raffinement et des conditions de stabilisation. Les notions de frontière de compatibilité, points fixes et profils de stabilisation sont des prolongements des propriétés de transport des D\_Base.

\section{Généralisation observationnelle}
La généralisation est ici définie géométriquement : une représentation intermédiaire correspond à une coupe non terminale de la filtration des seuils. Elle n'est pas identifiée à une métrique de test statistique. Pour qu'une généralisation soit globale, il faut néanmoins préciser le support couvert, les invariances et la relation entre sondes utilisées à l'apprentissage et sondes utilisées à l'évaluation.

\section{Sondes apprises}
Une sonde apprise peut être ajoutée à la famille d'observation comme opération expérimentale. Elle ne devient une « vérité sémantique » que si son interprétation est certifiée par le protocole. Une sonde qui modifie seulement la représentation sans modifier les fibres ne peut pas être appelée, par cette seule raison, source d'information nouvelle.

\section{Méta-observation}
Observer une représentation ou observer l'observateur ne crée pas automatiquement une nouvelle discrimination. Toute amélioration doit être mesurée comme un changement des fibres, de la fermeture ou du support de sondes. La méta-observation peut révéler une information déjà factorisée ; elle ne suffit pas à démontrer une rupture de fermeture.

\chapter{Langage, sémantique, prédiction et génération}
\section{Séparation des niveaux}
Les D\_Dev consacrés au langage, aux empreintes neuronales, à la prédiction et au pipeline complet imposent une séparation stricte : donnée brute, représentation, sémantique, prédiction, génération et action ne sont pas une seule opération. Une sortie neuronale est une observation ou une réalisation de représentation tant que son interprétation sémantique n'a pas été certifiée.

\section{Empreinte multilingue}
Le traitement multilingue distingue la structure sémantique supposée commune, les réalisations linguistiques et les transformations de surface. La théorie n'autorise pas à déclarer une équivalence de sens par simple ressemblance de forme ; l'équivalence doit être définie relativement aux observables du protocole.

\section{Prédiction et action}
Une prédiction est un objet produit par une règle de transition sur l'information disponible. L'action ajoute une interaction avec un environnement ou un état d'exécution. Cette distinction est conservée dans le pipeline complet : memoire, inférence, prédiction, attention, génération et action sont des modules typés et auditables.

\chapter{Déficit, témoins et transformations certifiées}
\section{Déficit structurel}
Les extensions mathématiques introduisent progressivement un déficit qui représente ce qui manque pour passer de l'état courant à un état cible. Il est crucial de séparer : déficit d'observation \(D^{obs}\), déficit de capacité \(D^{cap}\), déficit de ressource/applicabilité \(D^{res}\), difficulté de recherche \(D^{search}\), déficit de transfert \(D^{transfer}\) et déficit de contrainte \(D^{constraint}\). Ces composantes sont des mesures opérationnelles ; elles ne sont pas des impossibilités mathématiques.

On écrit dans O51
\[
D=(D^{obs},D^{cap},D^{res},D^{search},D^{transfer},D^{constraint}).
\]
Pour deux états,
\[
\Delta D=D_{avant}-D_{apres}
\]
mesure une progression dans le protocole. Par elle-même, une réduction positive ne constitue pas un certificat d'émergence.

\section{Témoins et sondes de séparation}
Une sonde peut être choisie pour distinguer une fibre. Une transformation \(a\) sépare une fibre \(F\) lorsqu'au moins deux éléments de \(F\) reçoivent des sorties différentes sous le support d'observation considéré. La notion est donc relative à l'espace de sortie et à l'observabilité choisie.

\section{Transformation certifiée}
Une transformation utile doit être exécutable, légale, traçable et validable. Le langage expérimental peut être abstrait sous la forme
\[
T::=\operatorname{root}(r)\mid\operatorname{use}(P)\mid\operatorname{seq}(T,T),
\]
les racines \(r\) appartenant à un ledger explicitement déclaré. Un hash peut fixer une identité de programme ; il n'est pas lui-même l'opération sémantique.

\part{Frontières de fermeture et recherche}

\chapter{Fermeture interne et critère de rupture}
\section{Fermeture observationnelle et exécutable}
Lorsque toutes les transformations sont constantes sur les fibres déjà observées ou appartiennent à la clôture des générateurs existants, la composition n'ajoute aucune nouvelle distinction comportementale. Cela explique les résultats négatifs répétés des ordres O27 à O34 : réflexion, imagination syntaxique, quotation/evaluation et certificats peuvent changer la représentation sans sortir de la fermeture.

\begin{theorem}[Conservativité de la promotion interne]
Soit \(\mathcal B'\) obtenu de \(\mathcal B\) uniquement par promotion de programmes dont les symboles sémantiques appartiennent déjà à \(\mathcal B\). Alors
\[
\cl_{exec}(\mathcal B')\subseteq \cl_{exec}(\mathcal B).
\]
\end{theorem}
\begin{proof}
Tout programme promu est interprété par les transitions de \(\mathcal B\). Son comportement appartient donc à la clôture exécutable de \(\mathcal B\). La composition de programmes ayant cette propriété reste dans la même clôture.\qed\end{proof}

\section{Critère de rupture de niveau}
Les extensions O26 et suivantes formulent une rupture possible à partir de niveaux typés \(W_\alpha\), \(W_\beta\). Une capacité candidate \(C^*\) manifeste une extension stricte du comportement seulement si
\[
C^*\in\Reach(W_\beta),
\qquad
C^*\notin\cl_\alpha(G_\alpha),
\]
et si le problème visé est solvable à \(\beta\) mais pas avec les ressources \(\alpha\). Cette condition est un critère relatif au ledger et au codomaine ; elle n'est pas une preuve de création « ex nihilo » au sens absolu.

\section{Contre-exemple général à la monotonie naïve}
Une contrainte plus restrictive peut supprimer une branche utile. Ainsi une diminution de l'ensemble admissible n'implique pas nécessairement une amélioration de tous les composantes du déficit. Le système doit donc conserver les campagnes négatives et distinguer progrès local, réduction de recherche et résolution effective.

\chapter{Recherche récursive et imagination formalisée}
\section{O27--O34 : résultats de fermeture}
Les ordres O27 à O34 ont servi de laboratoires de falsification interne. Ils montrent successivement que : la réflexion peut être conservative ; la redescente ne crée pas une opération absente ; une hiérarchie de représentations ne crée pas à elle seule une hiérarchie comportementale ; l'imagination peut produire des objets de langage sans ajouter une primitive ; la certification peut être complète dans le monde déclaré sans devenir une certification indépendante de la sémantique elle-même.

\section{O35--O39 : simulation, interprétation et synthèse}
O35 sépare simulation et exécution. O36 produit un langage génératif dérivé de primitives. O38 dérive un évaluateur à partir de primitives déclarées et vérifie que celui-ci reste dans la fermeture. O39 effectue une synthèse de tri pour de petits \(n\), avec résultats bornés et ablations. Ces résultats sont des démonstrations de construction dans des mondes finis ; ils ne constituent pas une capacité universelle.

\section{O40 : planification et deux domaines}
O40 remplace une recherche gloutonne par A* et applique le même noyau à deux environnements finis. Sur ces domaines, l'optimalité est démontrable lorsque les hypothèses d'A* sont satisfaites et l'espace est suffisamment borné. Le résultat important pour la théorie est méthodologique : le vérificateur fait partie des ressources racines. On ne peut donc pas compter la vérification comme une ressource gratuite découverte par le système.

\section{O41--O45 : preuve comme objet et frontière ouverte}
O41 à O43 formalisent des états de preuve, les statuts d'arrêt et la séparation entre même moteur de replay et vérification indépendante. O44 déroule dix micro-contextes puis atteint une cible Riemann pour laquelle aucune ressource applicable n'est présente ; il ne s'agit ni d'une preuve ni d'une réfutation. O45 ouvre la recherche récursive sans borne interne, mais l'exécution historique a été arrêtée par une limite physique de ressources. Le statut correct est un arrêt par ressources, non un résultat mathématique négatif.

\section{O46--O47 : Global/Internal et promotion}
La boucle Global/Internal sépare la sélection de capacité et la réobservation de l'état. O47 ajoute le déficit explicite, la promotion de principes exécutables et la validation hors construction. Les limites restent néanmoins nettes : primitives mathématiques déclarées, contextes construits et espace de recherche restreint.

\chapter{Théorie des contraintes, promotion et transfert}

Le déficit rend explicite ce qui manque à une tâche relativement à un état, une observation et un protocole. La théorie des contraintes introduite dans les extensions ne remplace pas cette structure : elle agit sur l'espace des transitions admissibles et doit rester reliée à une provenance observable.

\section{Contrainte comme objet de la théorie}
\begin{definition}[Objet contrainte]
Une contrainte est un tuple
\[
C=(\mathrm{Dom},\varphi,\mathrm{scope},\ell,\mathrm{deps},\mathrm{cost},\mathrm{prov}),
\]
où le domaine, le prédicat, le périmètre, le niveau, les dépendances, le coût et la provenance sont déclarés. Pour un espace de recherche $S$, l'espace admissible est
\[
S_C=\{s\in S: \varphi_C(s)\}.
\]
Pour une famille $\mathcal C$ de contraintes,
\[
S_{\mathcal C}=\bigcap_{C\in\mathcal C}S_C.
\]
\end{definition}

\begin{definition}[Raffinement de contrainte]
On écrit $C'\preceq C$ lorsque $S_{C'}\subseteq S_C$. Le raffinement est strict lorsque l'inclusion l'est.
\end{definition}

\begin{proposition}[Monotonie de l'admissibilité]
Si $C'\preceq C$, toute transition admissible sous $C'$ est admissible sous $C$.
\end{proposition}
\begin{proof}
L'assertion est exactement l'inclusion $S_{C'}\subseteq S_C$.
\end{proof}

\section{Sécurité, fuite d'information et provenance}
Une contrainte est sûre relativement à un protocole lorsque ses dépendances sont dérivées uniquement de la base d'information publique autorisée, des observations, du déficit, des objectifs publics et des ressources déclarées. Une contrainte qui code directement ou indirectement une solution cachée, un témoin optimal ou une réponse interdite est fuyarde. Une analyse de mots-clés ne suffit pas à établir cette absence de fuite : le contrôle porte sur la dépendance sémantique, la frontière des modules et les artefacts du ledger.

\begin{warning}[Portée]
La sûreté est une propriété relative au protocole inspecté. Elle ne constitue pas un théorème de non-interférence pour tout programme possible.
\end{warning}

\section{Génération et recherche sous contraintes}
Le générateur de contraintes peut être formalisé par une application partielle
\[
\operatorname{Ref}: (I,D,L,R,\mathcal C)\longmapsto\text{contraintes candidates},
\]
où seuls l'empreinte $I$, le déficit $D$, la bibliothèque $L$, les ressources $R$ et les contraintes publiques courantes sont accessibles. La théorie exige que la provenance de chaque contrainte produite soit traçable.

Une contrainte active n'est pas une annotation : elle modifie le domaine des transitions testées. Une recherche sous contraintes peut alors être représentée par
\[
\operatorname{Search}(S,\mathcal C,R,D),
\]
avec une frontière dont chaque transition satisfait toutes les contraintes actives. Un budget de profondeur, d'états ou d'expansions borne seulement le préfixe effectivement parcouru.

\begin{theorem}[Complétude conditionnelle de la recherche par approfondissement]
Supposons qu'une solution existe à profondeur finie $d^*$, que le branchement soit fini à chaque profondeur, que l'exécution et le test de suffisance soient décidables sur le domaine considéré et que toutes les frontières jusqu'à $d^*$ soient effectivement parcourues. Alors l'approfondissement itératif rencontre une solution à une profondeur au plus $d^*$.
\end{theorem}
\begin{proof}
La solution finie appartient à une frontière de profondeur au plus $d^*$. Sous l'hypothèse de branchement fini, cette frontière est finie et l'énumération complète finit par rencontrer le chemin solution ; le vérificateur le reconnaît.
\end{proof}

\section{Déficit de capacité : reconstruction après correction}
Dans les sources O50, la quantité $D^{cap}=\sum_i w_i\min\{1,r_i(L)\}$ mélange une distance générique $r_i$ et une lecture binaire utilisée par le prototype. Pour éviter cette ambiguïté, le livre retient une définition opérationnelle explicite :
\[
D^{cap}_{0/1}(L)=\sum_i w_i\mathbf 1_{\{\text{aucun chemin certifié admissible pour }i\text{ dans }L\}}.
\]
Une version métrique peut être définie séparément lorsqu'une notion de distance certifiée $r_i$ et son échelle sont fixées. Les deux quantités ne doivent pas être confondues.

\section{Utilité : statut expérimental}
La formule
\[
U(C'\mid\mathcal C)=\frac{\Delta D+\lambda\Delta I}{1+\operatorname{cost}(C')}
\]
est conservée comme score expérimental sous trois déclarations supplémentaires : les directions de $\Delta D$ et $\Delta I$, l'échelle de leurs unités et la valeur de $\lambda$ doivent être fixées par le protocole. Il ne s'agit pas d'un invariant de la théorie ni d'un théorème d'optimalité.

\section{Bibliothèque append-only et promotion}
Si la bibliothèque est append-only, une promotion ajoute une capacité sans invalider les capacités précédentes. L'idempotence de promotion demande une clé d'identification et un critère de doublon explicitement déclarés. Sous validation par replay, certificat et graphe acyclique, la promotion conserve les capacités antérieures et ajoute un objet exécutable déterminé sur le support certifié.

\section{Transfert et émergence relative}
Une capacité est dite nouvelle seulement relativement à une bibliothèque et à un support de validation : il faut distinguer différence syntaxique, différence de représentation et différence comportementale. Un transfert vers un support $H'$ n'est acquis que si le corps exécutable est rejoué sur $H'$ avec les préconditions déclarées et atteint l'objectif public du support. L'énoncé d'émergence
\[
\operatorname{Emerges}(P\mid R,\mathcal C,H)
\]
est donc une propriété de protocole combinant ressources, contraintes, exécution, déficit et support de validation ; il ne constitue pas une affirmation d'apparition ex nihilo.

\section{Résultats expérimentaux O50--O51 : statut}
Les textes O50--O51 décrivent une architecture expérimentale où les contraintes sont réellement utilisées par la recherche, où les capacités promues sont exécutables et où des contrôles de récursion, transfert, ablation et anti-oracle sont déclarés. Le corpus source actuel ne contient toutefois aucun fichier source Python, notebook ou package correspondant aux noms de modules cités dans O51. Les commandes et chiffres de campagne sont donc conservés comme résultats rapportés par les documents, et non comme une reproduction indépendante effectuée par l'ORDER 4.

\begin{warning}[Résultat non reproduit indépendamment]
L'absence du code exécutable dans le corpus actuel empêche de transformer les nombres de campagne O51 en validation indépendante. Cette limite porte sur la reproductibilité documentaire, pas sur la cohérence interne des définitions mathématiques du langage de contraintes.
\end{warning}


\part{Reconstruction, extensions et niveaux de preuve}

\chapter{Circulation interne : profils, trajectoires, identité et recollement}
\label{chap:restauration-dbase}

Cette partie est ajoutée à l'ordre 1 parce que l'audit de couverture a montré que la version précédente décrivait les D\_Base surtout par synthèse alors que ces fichiers contiennent des résultats structurants qui méritent une exposition autonome. Les résultats ci-dessous reprennent les objets et énoncés effectivement présents dans les trois fichiers D\_Base. Les démonstrations sont reformulées uniquement lorsque cela est nécessaire pour rendre explicites les prémisses.

\section{Dynamique observationnelle et circulation des empreintes}

\begin{definition}[Régime observationnel et empreinte]
Un régime est un couple \(\rho=(P,\mathcal G)\), où \(P\) est un contexte d'observation et \(\mathcal G\) une famille finie de sondes. L'empreinte induite est une application
\[
I_\rho:X\longrightarrow L_\rho.
\]
La relation d'indiscernabilité est
\[
x\sim_\rho y\quad\Longleftrightarrow\quad I_\rho(x)=I_\rho(y),
\]
et le quotient est \(\Pi_\rho=X/{\sim_\rho}\), avec projection canonique \(q_\rho\).
\end{definition}

\begin{definition}[Compatibilité de transition]
On écrit \(\rho\rightsquigarrow\sigma\) lorsque
\[
I_\rho(x)=I_\rho(y)\Longrightarrow I_\sigma(x)=I_\sigma(y)
\qquad\forall x,y\in X.
\]
La compatibilité signifie donc que toute indiscernabilité conservée par \(\rho\) le reste par \(\sigma\). La réflexivité et la transitivité en découlent immédiatement ; la compatibilité forme ainsi un préordre.
\end{definition}

\begin{theorem}[Théorème fondamental de circulation des empreintes]
Pour deux régimes \(\rho,\sigma\), les propriétés suivantes sont équivalentes :
\begin{enumerate}
\item \(\rho\rightsquigarrow\sigma\) ;
\item il existe une unique application \(T_{\rho\to\sigma}:L_\rho\to L_\sigma\) telle que \(T_{\rho\to\sigma}\circ I_\rho=I_\sigma\) ;
\item il existe une unique application \(\overline T_{\rho\to\sigma}:\Pi_\rho\to\Pi_\sigma\) telle que \(\overline T_{\rho\to\sigma}\circ q_\rho=q_\sigma\).
\end{enumerate}
\end{theorem}

\begin{proof}
L'implication \(1\Rightarrow2\) est la propriété de factorisation : si deux éléments de \(L_\rho\) sont produits par deux objets \(x,y\) avec la même valeur, la compatibilité impose que leurs images par \(I_\sigma\) coïncident ; on peut donc poser \(T_{\rho\to\sigma}(I_\rho(x)):=I_\sigma(x)\). La définition est indépendante du représentant et l'unicité suit de la surjectivité de \(I_\rho\) sur son image réduite. L'implication \(2\Rightarrow1\) s'obtient en composant l'égalité de factorisation avec l'hypothèse \(I_\rho(x)=I_\rho(y)\). Le raisonnement sur les classes d'équivalence donne de même \(1\Leftrightarrow3\).
\end{proof}

\begin{theorem}[Composition des transitions]
Si \(\rho\rightsquigarrow\sigma\rightsquigarrow\tau\), alors
\[
T_{\rho\to\tau}=T_{\sigma\to\tau}\circ T_{\rho\to\sigma},
\qquad
\overline T_{\rho\to\tau}=\overline T_{\sigma\to\tau}\circ\overline T_{\rho\to\sigma}.
\]
En particulier \(T_{\rho\to\rho}=\mathrm{Id}\) et \(\overline T_{\rho\to\rho}=\mathrm{Id}\).
\end{theorem}

\begin{proof}
Les deux membres composés prennent, sur toute empreinte \(I_\rho(x)\), la valeur \(I_\tau(x)\). L'unicité du transport canonique impose donc l'égalité. Le cas identité est le même argument avec \(\tau=\rho\).
\end{proof}

\section{Ordre interne, profils et représentation canonique}

Les D\_Base ne réduisent pas la circulation à une simple suite de fonctions. Ils construisent un ordre interne des régimes et un profil d'antériorité.

\begin{definition}[Ordre de transport]
Sur les classes de régimes, on définit \([\rho]\trianglelefteq[\sigma]\) lorsque la classe de \(\rho\) est antérieure à celle de \(\sigma\) dans le préordre de compatibilité quotienté par équivalence mutuelle. Les sources montrent que, sur cet ensemble quotienté, \(\trianglelefteq\) est réflexif, transitif et antisymétrique : c'est donc un ordre partiel interne.
\end{definition}

\begin{definition}[Profil d'antériorité]
Pour un instant observationnel \([\kappa]\), le profil d'antériorité est l'idéal principal
\[
\operatorname{Prof}([\kappa])=\{[\lambda]:[\lambda]\trianglelefteq[\kappa]\}.
\]
Il représente exactement les niveaux antérieurs accessibles par l'ordre interne déclaré.
\end{definition}

\begin{proposition}[Monotonie et séparation des profils]
Si \([\lambda]\trianglelefteq[\kappa]\), alors \(\operatorname{Prof}([\lambda])\subseteq\operatorname{Prof}([\kappa])\). Si deux classes ont des profils réciproquement inclus et que l'ordre interne est antisymétrique, elles coïncident.
\end{proposition}

\begin{proof}
La première inclusion suit par transitivité. La seconde est une application directe de l'antisymétrie : les deux inclusions donnent \([\lambda]\trianglelefteq[\kappa]\) et \([\kappa]\trianglelefteq[\lambda]\), donc les classes sont égales.
\end{proof}

\begin{theorem}[Représentation canonique de l'ordre interne]
L'application
\[
\mathcal A:[\kappa]\longmapsto\operatorname{Prof}([\kappa])
\]
représente fidèlement l'ordre interne :
\[
[\lambda]\trianglelefteq[\kappa]
\Longleftrightarrow
\mathcal A([\lambda])\subseteq\mathcal A([\kappa]).
\]
\end{theorem}

\begin{proof}
Le sens direct est la monotonie précédente. Pour la réciproque, l'inclusion des profils implique en particulier \([\lambda]\in\operatorname{Prof}([\kappa])\), donc \([\lambda]\trianglelefteq[\kappa]\). L'application est donc un plongement fidèle de l'ordre dans les idéaux principaux.
\end{proof}

\begin{warning}[Portée de la minimalité]
Les sources appellent cette représentation « minimale » ou « initiale ». Cette propriété doit être lue dans la catégorie des représentations qui préservent et reflètent effectivement \(\trianglelefteq\), avec les conditions de morphisme déclarées dans le document source. Elle ne constitue pas une minimalité absolue parmi toutes les structures mathématiques possibles.
\end{warning}

\section{Trajectoires, chaînes et bifurcations}

\begin{definition}[Trajectoire observationnelle]
Une trajectoire est une famille \((\rho_t)_{t\in J}\) indexée par un ensemble ordonné \(J\), telle que \(s\le t\Rightarrow\rho_s\rightsquigarrow\rho_t\). Pour \(x\in X\), elle produit la trajectoire d'empreinte \(E_x(t)=I_{\rho_t}(x)\).
\end{definition}

\begin{proposition}[Propagation de l'indiscernabilité]
Si \(E_x(t_0)=E_y(t_0)\), alors \(E_x(t)=E_y(t)\) pour tout \(t\ge t_0\).
\end{proposition}

\begin{proof}
La compatibilité fournit, pour chaque \(t\ge t_0\), un transport canonique \(T_{\rho_{t_0}\to\rho_t}\). En appliquant ce transport à l'égalité \(E_x(t_0)=E_y(t_0)\), on obtient \(E_x(t)=E_y(t)\).
\end{proof}

\begin{definition}[Concaténation associative partielle]
Deux trajectoires sont composables lorsque la cible de la première coïncide avec la source de la seconde. La concaténation \(\Gamma\star\Delta\) est alors la trajectoire obtenue en juxtaposant leurs chaînes compatibles. La composition est associative dès que les trois concaténations concernées sont effectivement définies.
\end{definition}

\begin{definition}[Bifurcation]
Une bifurcation apparaît lorsqu'un préfixe commun admet au moins deux prolongements compatibles distincts. Cette notion est relative à l'ordre et à l'ensemble des régimes considérés ; elle n'implique pas à elle seule une multiplicité ontologique du système sous-jacent.
\end{definition}

\section{Identité, tour et réflexion : restauration D\_Base}

\begin{definition}[Famille cohérente d'empreintes visibles stabilisées]
Sur un support \(S\) stable par antériorité, une famille \((\mathcal I_\kappa)_{[\kappa]\in S}\) est cohérente lorsque les restrictions imposées par l'ordre interne coïncident avec les transports canoniques prévus par la théorie. Une telle famille encode simultanément les vues visibles de l'empreinte à tous les niveaux antérieurs du support.
\end{definition}

\begin{definition}[Identité principale]
Pour une empreinte au niveau \([\kappa]\), son identité principale est la famille cohérente de ses restrictions sur
\[
S_{[\kappa]}=\{[\lambda]:[\lambda]\trianglelefteq[\kappa]\}.
\]
Elle n'ajoute pas une information externe : elle organise les restrictions déjà forcées par la structure observationnelle.
\end{definition}

\begin{theorem}[Existence, unicité et auto-décomposition]
Sous la cohérence et les règles de restriction des D\_Base, toute empreinte visible stabilisée admet une identité principale unique. La restriction de l'identité principale à un niveau antérieur coïncide avec l'identité principale de la composante considérée.
\end{theorem}

\begin{proof}
L'existence est donnée par la famille des restrictions canoniques de l'empreinte. La cohérence garantit qu'une même composante obtenue par deux chemins compatibles possède la même valeur. Pour l'unicité, toute famille cohérente ayant la même composante maximale doit, par les conditions de restriction, prendre exactement ces valeurs sur chaque niveau antérieur. L'auto-décomposition est alors immédiate par restriction.
\end{proof}

\begin{definition}[Sous-système cofinal et recollement]
Un sous-système \(T\subseteq S\) est cofinal si tout élément de \(S\) est antérieur à un élément de \(T\). Une famille cohérente sur \(T\) est recolleable lorsque les accès concurrents à une même composante donnent les mêmes restrictions.
\end{definition}

\begin{theorem}[Recollement canonique]
Une famille cohérente recolleable sur un sous-système cofinal détermine une unique famille cohérente sur tout le support.
\end{theorem}

\begin{proof}
Pour chaque niveau \([\lambda]\in S\), choisir un niveau cofinal \([\kappa]\in T\) au-dessus de lui et prendre la restriction canonique de \(\mathcal J_\kappa\) vers \([\lambda]\). La recollabilité assure l'indépendance du choix. La cohérence résulte de la composition des transports et l'unicité de la détermination sur les niveaux cofinaux.
\end{proof}

\begin{definition}[Construction canonique interne admissible]
Il s'agit d'une construction obtenue par composition finie de restrictions, transports canoniques, profils, filtrations, quotients, factorisations, représentations minimales et comparaisons de familles cohérentes, sans consultation d'une donnée extérieure à la structure ainsi définie.
\end{definition}

\begin{theorem}[Clôture interne : portée exacte]
Toute lecture qui utilise uniquement les objets précédents est déterminée par l'identité principale correspondante. Le passage à l'identité est donc sans perte pour ce langage interne.
\end{theorem}

\begin{warning}[Lecture de la réflexion]
L'idempotence des tours/réflexions des D\_Base est un résultat de fermeture interne. Elle ne suffit pas, à elle seule, à établir une création de capacité comportementale. Une nouveauté comportementale exige une nouvelle ressource, un nouveau codomaine, un nouvel accès observationnel ou un autre changement de frontière explicitement certifié.
\end{warning}

\section{D\_Base : invariants, chaînes infinies et morphismes}

Le troisième D\_Base ajoute une couche d'invariants qui a été insuffisamment visible dans la version précédente.

\begin{definition}[Signature de chaîne et invariant de trajectoire]
Une signature de chaîne encode les niveaux observés et leurs données de transport selon une normalisation fixée. Un invariant de trajectoire est une quantité qui reste inchangée sous les transports compatibles admissibles et les transformations de représentation autorisées.
\end{definition}

\begin{proposition}[Rigidité par signature]
Deux trajectoires possédant la même signature normalisée coïncident pour les invariants qui sont définis comme fonctions de cette signature. Cette conclusion est relative à la classe d'invariants déclarée ; elle ne constitue pas une identité de toutes les données non observées.
\end{proposition}

\begin{definition}[Chaîne infinie et stabilisation observationnelle]
Une chaîne infinie est une suite ordonnée de régimes compatibles. Une stabilisation des profils signifie qu'au-delà d'un rang, les profils observés deviennent constants ou stabilisés au sens précis de l'ordre déclaré.
\end{definition}

\begin{theorem}[Fermeture observationnelle sous stabilisation]
Sous les hypothèses de cofinalité/stabilisation utilisées dans le D\_Base, toute donnée qui dépend seulement des profils stabilisés est déterminée par le niveau stabilisé correspondant.
\end{theorem}

\begin{definition}[Morphism observatoire]
Un morphisme observationnel est une application entre structures observationnelles qui respecte les empreintes, les transports et les relations d'indiscernabilité selon les conditions annoncées.
\end{definition}

\begin{proposition}[Stabilité sous morphisme]
Sous les conditions de morphisme, l'image d'une chaîne compatible est une chaîne compatible et les profils correspondants sont préservés au sens du morphisme.
\end{proposition}

\section{Lecture consolidée des D\_Base}

Les D\_Base établissent donc une continuité théorique nette :
\[
\begin{gathered}
\text{empreinte}\to\text{fibre/quotient}\to\text{compatibilité}\to\text{transport},\\
\text{transport}\to\text{ordre interne}\to\text{profil}\to\text{identité},\\
\text{identité}\to\text{tour/réflexion}.
\end{gathered}
\]
La dynamique interne, l'identité et les chaînes ne sont pas des ajouts indépendants de la théorie initiale : elles réutilisent systématiquement les mêmes factorisations observationnelles. Les développements ultérieurs doivent être raccordés à cette chaîne avant d'être qualifiés de nouvelles capacités.

\chapter{Généralisation observationnelle, apprentissage et recherche}
\section{Généralisation observationnelle globale}

Le corpus D\_Dev contient une ligne de développement qui mérite d'être restaurée explicitement : la généralisation observationnelle globale et la notion de sens doivent être définies sur le quotient effectivement vu par la famille admissible de sondes.

\begin{definition}[Quotient observationnel global]
Soit \(\mathcal R_{\mathrm{adm}}\) une famille de régimes admissibles. On pose
\[
x\equiv_{\mathrm{glob}}y
\quad\Longleftrightarrow\quad
I_\rho(x)=I_\rho(y)\quad\forall\rho\in\mathcal R_{\mathrm{adm}}.
\]
Le quotient global est \(Q_{\mathrm{glob}}=X/{\equiv_{\mathrm{glob}}}\).
\end{definition}

\begin{theorem}[Facteurs globaux]
Toute quantité \(C:X\to Y\) qui factorise par chaque quotient admissible et qui est donc constante sur les classes de \(\equiv_{\mathrm{glob}}\) induit une application unique \(\overline C:Q_{\mathrm{glob}}\to Y\) telle que \(C=\overline C\circ q_{\mathrm{glob}}\).
\end{theorem}

\begin{warning}[Portée]
La généralisation est exacte sur le support effectivement couvert par les régimes admissibles. Elle ne donne pas automatiquement accès à des distinctions absentes du quotient global.
\end{warning}

\section{Prédiction dérivée et limites de l'action}

Les D\_Dev consacrés à la prédiction dérivent des fréquences conditionnelles à partir des branches effectivement observées. Le statut correct est probabiliste et conditionnel : pour une famille de branches de poids positifs, la règle
\[
\Pr(B\mid A)=\frac{N(A,B)}{N(A)}
\]
est une définition empirique/conditionnelle sur les observations disponibles, et non une loi métaphysique du monde.

\begin{proposition}[Indépendance par réordonnancement des observations]
Lorsque les comptes utilisés dans la loi restent inchangés, un réordonnancement des observations ne change pas les probabilités empiriques dérivées. Cette invariance est une propriété du protocole de comptage, pas une preuve de causalité.
\end{proposition}

\begin{warning}[Causalité]
Une prédiction conditionnelle ne devient pas une causalité par simple changement de vocabulaire. Une action causalement interprétée demande un mécanisme d'intervention, une transformation certifiée ou une hypothèse supplémentaire reliant l'intervention aux observations.
\end{warning}

\section{Formalisme DD\_Dev : déficit, générateurs et recherche}

Les versions DD\_Dev 9--15, ainsi que \texttt{book\_suite\_th\_emp\_2.tex}, développent une théorie de la complétude générative et de la recherche sous contrainte. L'audit indépendant distingue ici les énoncés mathématiques conditionnels des ambitions universelles.

\begin{definition}[Déficit structuré]
Pour une cible \(S\), le déficit \(D(S)\) est l'ensemble des obligations observables non satisfaites par l'état courant, relativement à la représentation et au protocole de satisfaction déclarés. Les sources distinguent notamment déficits d'observation, de capacité, de ressource/applicabilité, de recherche, de transfert et de contrainte.
\end{definition}

\begin{proposition}[Réduction exacte par témoins]
Lorsque chaque obligation du déficit peut être séparée par un ensemble fini de témoins et que la satisfaction d'un témoin est décidable sur le domaine considéré, la recherche d'un sous-ensemble de témoins satisfaisant toutes les obligations se reformule exactement comme un problème de couverture/hitting-set.
\end{proposition}

\begin{warning}[Condition de finitude]
L'équivalence avec un hitting-set calculable ne donne pas à elle seule une procédure universelle sur un espace infini ou indécidable. Les résultats algorithmiques du corpus supposent explicitement une frontière finie ou une procédure de décision/certification suffisante.
\end{warning}

\begin{theorem}[Complétude générative conditionnelle]
Soit \(G\) un ensemble de générateurs dont la clôture \(\operatorname{Cl}(G)\) est effectivement définie. Si toute cible d'une classe \(T\) admet une dérivation finie par les générateurs et si le vérificateur de suffisance est correct et décidable sur le domaine considéré, une énumération exhaustive des dérivations finies finit par rencontrer une dérivation suffisante.
\end{theorem}

\begin{proof}
\small
Une dérivation suffisante existe par hypothèse et est finie. Une énumération diagonale complète des séquences finies de générateurs finit par l'atteindre, puis le vérificateur correct l'accepte. L'argument ne porte ni sur les cibles sans dérivation ni sur un espace dont l'énumération ou la certification n'est pas effectivement disponible.
\normalsize
\end{proof}

\begin{proposition}[Terminaison sous réduction stricte du déficit]
Sur un ordre bien fondé de déficits, si chaque étape acceptée réduit strictement le déficit et si aucune infinité descendante n'existe, une recherche guidée termine après un nombre fini d'étapes.
\end{proposition}

\begin{warning}[Sélection gloutonne et EXHAUSTED]
Un algorithme glouton qui ne trouve plus de transition admissible est seulement \emph{GREEDY\_STUCK}. Le statut \emph{EXHAUSTED} exige au contraire une preuve que toute la frontière de recherche déclarée a été épuisée. Cette distinction, déjà signalée dans O49/O50, est conservée comme exigence de provenance.
\end{warning}

\section{Bisimulation, pruning et optimisation}

Le corpus ultérieur distingue utilement deux types de pruning : une réduction démontrée par bisimulation/commutation, et une suppression heuristique. Seule la première conserve une portée théorique de préservation.

\begin{proposition}[Pruning conservatif]
Si une relation \(\approx\) est une bisimulation pour le système de dérivation et que les deux états reliés ont les mêmes possibilités de satisfaction jusqu'à la cible considérée, fusionner ces états ne modifie pas l'existence d'une solution dans le quotient.
\end{proposition}

\begin{warning}[A*]
Les conclusions d'optimalité d'A* sont conditionnelles aux hypothèses usuelles de complétude, de coût borné/non négatif et, selon le résultat visé, d'admissibilité/cohérence de l'heuristique. Elles ne doivent pas être étendues aux campagnes expérimentales où ces conditions n'ont pas été vérifiées.
\end{warning}

\section{Architecture IA et universalités : statut exact}

Les DD\_Dev formalisent plusieurs universalités distinctes : représentation, génération, recherche, transfert, causalité active et « intelligence ». L'audit conserve cette pluralité et refuse de les confondre.

\begin{definition}[Universalité effective]
Une procédure est effectivement universelle pour une classe de tâches si, pour chaque tâche de cette classe, l'existence d'une solution, son espace de candidats et la procédure de vérification sont tous spécifiés de manière calculable/exécutable dans le domaine déclaré.
\end{definition}

\begin{warning}[Universalité de l'intelligence]
Le corpus ne fournit pas une preuve générale d'une universalité absolue identifiant le modèle des empreintes aux capacités d'un LLM dans toutes leurs dimensions. Ce résultat reste ouvert. Les éléments démontrés concernent des sous-problèmes précisément spécifiés et des domaines explicitement bornés.
\end{warning}

\chapter{Imagination formalisée, auto-observation et redescente}

La partie \emph{imagination} n'est pas introduite ici comme une faculté mystérieuse. Elle désigne une construction interne de scénarios, d'actions typées, de critères, de révisions et de certificats qui peut être examinée comme un objet de la théorie. La question scientifique porte alors sur le passage éventuel entre cet objet et une réalisation opératoire.

\section{Objet imaginaire}
\begin{definition}[Plan imaginaire]
Un plan imaginaire est un objet structuré
\[
\Pi=(s_0,a_1,\ldots,a_n,g,\mathsf{test},\mathsf{rev}),
\]
où $s_0$ est l'état observé, les $a_i$ sont des actions typées déclarées, $g$ est l'objectif public, $\mathsf{test}$ est le dispositif de vérification et $\mathsf{rev}$ décrit les règles de révision. Le plan appartient au laboratoire tant qu'aucun pont de réalisation vers un domaine externe n'est attesté.
\end{definition}

\section{Chaîne observation -- imagination -- révision}
La chaîne structurante peut être lue sur deux niveaux :
\[
\begin{array}{c}
\text{état observé}\to\text{plan imaginaire}\to\text{simulation}\to\text{certificat}
\cr
\text{révision}\to\text{auto-observation}\to\text{redescente}
\end{array}
\]
Chaque flèche est soumise à une condition différente. Une simulation n'est pas une exécution externe ; un certificat de cohérence n'est pas une preuve d'atteinte d'un objectif hors du domaine vérifié ; une redescente syntaxique n'est pas une nouvelle capacité comportementale.

\section{Conservativité de l'imagination interne}
Lorsque toutes les primitives d'un plan appartiennent à une fermeture exécutable $\cl(G)$ et que la révision ne fait qu'en composer des éléments, la construction imaginaire reste dans la fermeture. Cette proposition explique les contrôles négatifs répétés des laboratoires finis O27--O40. Elle ne dit pas que toute imagination possible est conservative : elle dit que la construction déclarée l'est tant que son vocabulaire de réalisation reste fermé.

\begin{noGo}[Pas de pont magique]
Aucune transformation interne de représentation ne peut, à elle seule, constituer une preuve d'une modification d'un monde extérieur qui n'appartient pas au domaine de réalisation déclaré. Pour obtenir une telle conclusion, il faut introduire et auditer séparément un pont de réalisation, ses ressources, ses préconditions et sa vérification.
\end{noGo}

\section{Apprentissage canonique : plusieurs propriétés à séparer}
Les sources et O21 imposent de distinguer : existence d'un minimum, unicité d'un minimiseur, choix canonique, et minimisation d'une énergie. Une construction peut être minimale par inclusion sans être unique ; un sélecteur déterministe peut être canonique relativement à un ordre lexicographique choisi sans être la canonicité intrinsèque de la théorie.

\section{Recherche de nouveautés}
La notion de nouveauté est relativisée à un support de validation. Une construction peut être nouvelle comme procédure, tout en étant dérivable des primitives mathématiques déclarées. Inversement, sortir d'une fermeture syntaxique ne suffit pas à prouver une nouveauté sémantique. Cette distinction permet de conserver les résultats expérimentaux sans leur attribuer une portée universelle.

\section{Statut scientifique}
Les pistes d'imagination qui ne sont pas démontrées sont conservées dans le document d'essai. Lorsqu'une piste devient une construction mathématique correctement typée, elle peut rejoindre le corps du livre sous son hypothèse exacte. Les résultats O21 relatifs à l'imagination sont ainsi intégrés comme constructions et no-go conditionnels, non comme affirmations psychologiques ou physiques.


\chapter{Reconstruction hiérarchique, seuils et information conservée}

La reconstruction hiérarchique doit distinguer la donnée indexée par paire, la filtration de partitions, l'arbre étiqueté et les résumés scalaires. Cette distinction est la conséquence directe des contre-exemples de perte d'affectation.

\section{Donnée de seuil complétée}
Pour une paire distincte $\{x,y\}$, poser
\[
S_{x,y}=\{\lambda:x\sim_{\Phi,P,\lambda}y\},
\qquad
\Theta_P(\{x,y\})=\sup\{\iota(\lambda):\lambda\in S_{x,y}\}.
\]
La donnée de bord
\[
b_P(\{x,y\})=1\iff\exists\lambda\in S_{x,y}:\iota(\lambda)=\Theta_P(\{x,y\})
\]
complète le squelette. Avec $H_P=(\Theta_P,b_P)$, la lecture exacte devient
\[
x\sim_{\Phi,P,\lambda}y
\iff
\iota(\lambda)<\Theta_P(\{x,y\})
\ \text{ ou }\
(\iota(\lambda)=\Theta_P(\{x,y\})\land b_P(\{x,y\})=1).
\]

\section{Loi du minimum double}
Pour $x,y,z$ deux à deux distincts, la cohérence des fibres entraîne
\[
\Theta_P(\{x,z\})\geq
\min\bigl(\Theta_P(\{x,y\}),\Theta_P(\{y,z\})\bigr).
\]
La restriction aux triplets de points distincts est essentielle : $\Theta_P$ est définie sur $\binom{X}{2}$. La même propriété produit des relations superniveau réflexives et une filtration de partitions.

\section{Représentation quotient et canonicité}
Une filtration finie possède une réalisation par quotients
$X/{\sim_t}$. Cette réalisation est canonique au sens précis de la propriété universelle du quotient. Une représentation arbitraire partageant les mêmes partitions peut cependant porter des données de codage supplémentaires ; la filtration ne les détermine donc pas automatiquement.

\section{Ce que conserve la mesure}
Pour $|X|\ge2$, la mesure empirique
\[
\mu_P=\frac{1}{\binom{|X|}{2}}
\sum_{\{x,y\}\subseteq X}\delta_{\Theta_P(\{x,y\})}
\]
conserve la distribution des seuils avec multiplicité. Elle oublie, en général, quelle paire porte quel seuil. De même, pour
\[
F_{\mu_P}(t)=\mu_P((-\infty,t)),
\]
on a
\[
F_{\mu_P}(b)-F_{\mu_P}(a)=\mu_P([a,b)).
\]
Toute formule utilisant cette convention doit conserver l'orientation des frontières. Une reconstruction de branche nécessitant l'affectation des seuils doit donc prendre comme entrée $\Theta_P$ (ou la hiérarchie équivalente), et non $\mu_P$ seul.

\section{Probabilités de branche : reconstruction conservatrice}
Si une classe $A$ est munie d'une règle d'échantillonnage uniforme sur ses éléments et si $B$ est un enfant non vide, on peut définir
\[
\mathbb P(B\mid A)=\frac{|B|}{|A|}.
\]
Cette probabilité dépend du protocole d'échantillonnage ; elle n'est pas une conséquence de la mesure scalaire globale $\mu_P$. Pour une chaîne de classes, le produit des probabilités conditionnelles décrit la probabilité de trajectoire sous cette règle.



\chapter{Ordre 5 : modèles, algèbre des contraintes et résolution mathématique}
Le passage à l'ordre 5 ne change pas les fondations ; il ajoute une couche de développement qui rend explicites trois interfaces nouvelles : l'interface entre contraintes et empreintes, l'interface entre modèles computationnels et théorie, et l'interface entre les seuils critiques et leurs données de frontière.

\section{Audit des nouveaux modèles O37--O51}
Le corpus complémentaire contient des archives de modèles des ordres O37, O38, O39, O40, O47, O48, O49, O50 et O51. Les archives sont traitées comme des objets expérimentaux. Pour O51, le package exécutable fourni a été rejoué deux fois avec la graine $5151$, quatre contextes, profondeur maximale $6$ et $12000$ expansions. Les sorties sont identiques au niveau du hachage. Cette reproductibilité est un résultat du package sur son support fini ; elle ne transforme ni la campagne en théorème de convergence ni le statut de ses primitives en preuve d'universalité.

\section{Correction de la sémantique des contraintes}
Une contrainte $C$ est interprétée par son ensemble de transitions admissibles $\Adm(C)$. Le raffinement est
\[
C_1\preceq C_2 \iff \Adm(C_1)\subseteq\Adm(C_2).
\]
La conjonction réalise l'intersection et la disjonction réalise l'union. Une relaxation mathématique cohérente peut donc être définie comme une disjonction contrôlée. Le code O51 fourni porte toutefois une opération nommée `relax` qui évalue seulement son premier argument ; l'artefact est conservé, mais il n'est pas utilisé comme preuve d'une sémantique de relaxation.

\section{Empreinte sous contraintes}
Pour une empreinte $F_H$ et une famille de prédicats $(C_j)$, l'ordre 5 introduit
\[
F_{H,\mathcal C}(x)=\bigl(F_H(x),C_1(x),\ldots,C_m(x)\bigr).
\]
Cette construction raffine les fibres de $F_H$. Si toutes les contraintes se factorisent par $F_H$, elle n'apporte aucune séparation supplémentaire. Elle fournit ainsi une interface précise entre la théorie des contraintes et la théorie des empreintes.

\section{Monotonie du déficit}
Pour une cible $A$, le déficit par collisions est
\[
\Def_F(A)=\{(x,y):F(x)=F(y),\ A(x)\ne A(y)\}.
\]
Si $F'$ raffine $F$, alors $\Def_{F'}(A)\subseteq\Def_F(A)$. Sur un domaine fini, le nombre de collisions incompatibles ne peut donc qu'être conservé ou diminuer sous raffinement de l'empreinte. Cette monotonie porte sur la définition précise du déficit ; elle ne permet pas d'affirmer qu'une stratégie de recherche computationnelle est globalement monotone.

\section{Seuils : le point d'égalité}
Le squelette $\Theta_P$ ne suffit pas à décider le comportement au seuil lorsque le supremum n'est pas connu comme atteint. L'ordre 5 conserve la donnée $b_P$ et utilise le critère exact :
\[
x\sim_{\Phi,P,\lambda}y\iff \iota(\lambda)<\Theta_P(\{x,y\})\ \text{ou}\ \bigl(\iota(\lambda)=\Theta_P(\{x,y\})\land b_P(\{x,y\})=1\bigr).
\]
Cette correction est propagée jusqu'aux partitions critiques et aux interprétations des mesures de seuils.

\section{Nouveauté finie et abstraction procédurale}
Les modèles O48--O51 utilisent des signatures comportementales sur des supports de validation finis. Une différence de signature prouve une différence locale, pas une différence globale. De même, un schéma procédural extrait de fragments communs doit satisfaire des conditions de compatibilité domaine/codomaine avant de devenir un programme exécutable. Cette distinction empêche de confondre une régularité observée avec une loi universelle.

\section{Imagination comme branche de résolution}
L'imagination reste traitée comme une construction interne : état observé, plan hypothétique, simulation, certificat, révision. Les contraintes peuvent filtrer les transitions de la simulation, et les sorties simulées peuvent elles-mêmes alimenter une empreinte raffinée. Les réalisations externes restent séparées par un pont de certification.

\section{Théorie des contraintes : statut final de l'ordre 5}
L'ordre 5 conserve les objets de contrainte, le raffinement, la promotion et le transfert relatifs aux protocoles O50--O51. Il corrige cependant le point faible du mot `relax` et distingue l'algèbre mathématique de l'implémentation historique. Les scores d'utilité restent heuristiques, les budgets restent des bornes d'exécution, et `SOLVED` reste un statut expérimental du protocole.

\section{Nouveaux développements mathématiques}
Les démonstrations complètes de la couche mathématique développée à l'ordre 5 sont données dans les dossiers ordre 52. Le corps du livre en conserve les résultats structurants ; les variantes et conjectures restent dans le laboratoire.

\chapter{Expériences, résultats bornés et critères de preuve}
\section{Résultats théoriques consolidés}
Les éléments les plus stables de la reconstruction sont : la définition observationnelle de l'empreinte ; l'équivalence par fibres ; la factorisation canonique ; la fermeture des sondes ; le raffinement monotone ; la fonction de seuil ; la loi du minimum double ; la reconstruction hiérarchique par coupes ; l'ultramétrique dérivée avec diagonale correctement définie ; la lecture en mesure atomique ; la composition des transports compatibles ; l'ordre interne ; l'idempotence des opérations de tour/réflexion dans la fermeture considérée ; les théorèmes de conservativité lorsqu'aucune ressource sémantique nouvelle n'est introduite.

\section{Résultats expérimentaux bornés}
Les ordres O31, O35, O38, O39, O40, O44, O45, O46, O47, O49, O50 et O51 fournissent des démonstrations computationnelles dans des mondes déclarés et finis, des tests de falsification, des audits et des mesures. Ils établissent ce qui a été exécuté dans leurs cadres précis ; ils n'établissent pas une intelligence générale ni la découverte autonome de nouvelles mathématiques.

\section{Questions ouvertes}
Restent ouvertes, entre autres : une construction de rupture de fermeture qui ne masque pas sa ressource supplémentaire ; des conditions générales de complétude des générateurs de témoins/contraintes ; des résultats de convergence sur des espaces récursifs non finis ; une théorie générale du transfert entre supports ; une caractérisation plus complète de la nouveauté sémantique sous changement de codomaine ; et toute universalité au-delà des univers explicitement déclarés.

\section{Hypothèse de progression conditionnelle}
Une conjecture récurrente des extensions est qu'une progression locale puisse être garantie sous génération complète des contraintes, capacité admissible et validation exhaustive du support. Cette proposition reste conditionnelle et conjecturale tant que ces hypothèses n'ont pas reçu une preuve générale dans un cadre précisément spécifié.


\chapter{État de certitude et frontières de l'ordre 5}
Les résultats de l'ordre 5 sont répartis en quatre niveaux. Les définitions et propositions démontrées utilisent uniquement les hypothèses déclarées. Les résultats conditionnels sont accompagnés de leurs conditions. Les campagnes O48--O51 sont expérimentales et bornées. Les conjectures, notamment celles relatives à la complétude de la génération de contraintes, restent ouvertes.

Aucune branche entière n'a été supprimée à la suite d'une correction. Les modifications principales sont des reconstructions : algèbre des contraintes, seuil critique complété, déficit sous raffinement, nouveauté finie et exécutabilité typée des abstractions. Le registre de traçabilité permet de revenir aux sources et aux artefacts de modèle.

\chapter{Frontières, contre-exemples, conjectures et programme de recherche}

Les limites font partie de la théorie consolidée : elles indiquent précisément où une conclusion s'arrête. Les contre-exemples ne sont pas des échecs de présentation ; ils empêchent des identifications trop fortes.

\section{Contre-exemples structurants}
Le premier contre-exemple est le bord du supremum : $S_1=(0,1)$ et $S_2=(0,1]$ ont le même supremum mais des comportements différents au point $1$. Le second est la conservation de distribution sans conservation d'affectation : deux fonctions de seuil sur un ensemble de quatre objets peuvent avoir le même multiensemble de valeurs mais permuter les paires porteuses du seuil élevé. Le troisième est la reparamétrisation monotone non injective, qui peut fusionner des niveaux.

\section{Domaines et normalisations}
Toute utilisation de $\binom{|X|}{2}$ dans une normalisation exige $|X|\ge2$. Toute division conditionnelle $\omega(B)/\omega(A)$ exige $\omega(A)>0$. Les cas vide, singleton et dégénéré doivent être traités séparément plutôt que recouverts implicitement par une convention.

\section{Claims d'universalité}
Les mots « universel », « canonique », « unique », « complet », « sans perte » sont interprétés relativement à leur domaine explicitement déclaré. La canonicité du quotient est une propriété de construction ; la canonicité d'un sélecteur ou d'un latent exige des choix supplémentaires. Les résultats d'existence et d'unicité ne sont jamais confondus.

\section{Questions ouvertes}
Restent notamment ouvertes : une théorie générale de la rupture de fermeture avec ressources nouvelles explicitement certifiées ; une complétude générale des générateurs de contraintes ; des résultats de convergence sur des espaces non finis ; une caractérisation générale du transfert entre supports ; et une formalisation catégorique complète du « latent universel » évoqué dans certaines sources.

\begin{conjecture}[Progression conditionnelle]
Sous une génération complète des contraintes, l'existence d'une capacité admissible, un vérificateur correct et une exploration exhaustive du support fini, une stratégie de raffinement peut conduire à une suite de déficits strictement décroissants jusqu'à satisfaction de la cible. Sans ces hypothèses, aucune convergence générale n'est affirmée.
\end{conjecture}


\appendix

\chapter{Dossier formel détaillé du socle DD\_Base}
Le présent dossier reprend, sous une forme compilable et traçable, le noyau formel du manuscrit fondamental \path{DD_Base_theorie_empreintes_informationnelles_seuils_corrigee_v7.tex} : définitions, propositions, théorèmes, démonstrations, exemples et développements théoriques ultérieurs présents dans ce manuscrit. Il est annexé à la reconstruction afin de conserver les arguments détaillés plutôt que de les réduire à un simple résumé. Les corrections explicitement enregistrées dans le journal d'erreurs sont appliquées lorsque leur portée est locale et vérifiable ; les passages exploratoires restent signalés comme tels.

\begin{warning}
Ce dossier est une transcription structurée du socle formel, et non une validation automatique de chaque phrase historique. Les statuts scientifiques consolidés du présent ouvrage priment lorsqu'une formulation ancienne a été corrigée ou restreinte par l'audit inter-source.
\end{warning}

\section{Partie I. Fondations}

\subsection{Cadre abstrait}

Soient
\[
\mathcal P,\qquad \mathcal X,\qquad \mathcal F
\]
trois ensembles non vides, où \(\mathcal P\) désigne l'ensemble des contextes d'observation, \(\mathcal X\) l'ensemble des objets étudiés, et \(\mathcal F\) l'ensemble des sondes, opérateurs d'interrogation ou contextes auxiliaires.

On fixe une famille finie de sondes
\[
\mathcal F_K=\{f_1,\ldots,f_K\}\subset \mathcal F,
\qquad K\in\mathbb N^*.
\]
Plus généralement, toute famille \(\mathcal G\subseteq\mathcal F\), finie ou non, peut être utilisée.

\begin{definition}[Fonctionnel d'observation]
Un \emph{fonctionnel d'observation} est une application
\[
\Phi:\mathcal P\times\mathcal X\times\mathcal F\longrightarrow \mathbb R.
\]
Pour tout triplet \((P,x,f)\), \(\Phi(P,x,f)\) est la réponse scalaire produite par l'observation de \(x\) dans le contexte \(P\) au moyen de la sonde \(f\).
\end{definition}

\begin{definition}[Empreinte informationnelle]
Pour chaque contexte \(P\in\mathcal P\), on définit
\[
I_\Phi^P:\mathcal X\longrightarrow \mathbb R^K,
\qquad
I_\Phi^P(x)=
\begin{bmatrix}
\Phi(P,x,f_1)\\
\vdots\\
\Phi(P,x,f_K)
\end{bmatrix}.
\]
Pour toute famille \(\mathcal G\subseteq\mathcal F\), on définit aussi
\[
I_{\Phi,\mathcal G}^P:\mathcal X\longrightarrow \mathbb R^{\mathcal G},
\qquad
I_{\Phi,\mathcal G}^P(x)(g)=\Phi(P,x,g).
\]
\end{definition}

\begin{remark}
Aucune structure interne sur \(\mathcal X\) n'est requise. La discrimination et les régimes d'indiscernabilité sur \(\mathcal X\) sont induits a posteriori par le choix de \(\Phi\) et des familles de sondes.
\end{remark}

\subsection{Théorie mathématique des échelles observationnelles}

L'échelle ne désigne pas ici une taille spatiale. Elle désigne la quantité d'information observationnelle accessible, donc la contrainte imposée par le régime d'observation.

\begin{definition}[Structure d'échelle observationnelle]
Une \emph{structure d'échelle} est la donnée d'un ensemble partiellement ordonné \((\Lambda,\preceq)\) et d'une famille croissante de sous-familles
\[
\bigl(\mathcal F_\lambda\bigr)_{\lambda\in\Lambda},
\qquad
\mathcal F_\lambda\subseteq \mathcal F,
\]
telles que
\[
\lambda\preceq\mu \Longrightarrow \mathcal F_\lambda\subseteq \mathcal F_\mu.
\]
On associe à cette structure une fonction d'intensité
\[
\iota:\Lambda\longrightarrow \mathbb R_+
\]
monotone pour l'ordre \(\preceq\), c'est-à-dire
\[
\lambda\preceq\mu \Longrightarrow \iota(\lambda)\le \iota(\mu).
\]
\end{definition}

\begin{definition}[Information accessible brute]
On note
\[
C(\lambda):=|\mathcal F_\lambda|
\]
la quantité brute d'information accessible à l'échelle \(\lambda\). Lorsque \(\mathcal F_\lambda\) est finie, \(C(\lambda)\in\mathbb N\); lorsque \(\mathcal F_\lambda\) est infinie, \(C(\lambda)\) est une cardinalité.
\end{definition}

\begin{definition}[Empreinte à l'échelle \(\lambda\)]
Pour tout \(P\in\mathcal P\) et tout \(\lambda\in\Lambda\), on définit
\[
I_{\Phi,\lambda}^P:\mathcal X\longrightarrow \mathbb R^{\mathcal F_\lambda},
\qquad
I_{\Phi,\lambda}^P(x)(f)=\Phi(P,x,f)\quad(f\in\mathcal F_\lambda).
\]
Lorsque \(\mathcal F_\lambda\) est finie de cardinal \(K_\lambda\), on identifie \(\mathbb R^{\mathcal F_\lambda}\simeq \mathbb R^{K_\lambda}\).
\end{definition}

\begin{definition}[Indiscernabilité à l'échelle \(\lambda\)]
On définit, pour tout \(P\in\mathcal P\) et \(\lambda\in\Lambda\), une relation \(\sim_{\Phi,P,\lambda}\) sur \(\mathcal X\) par
\[
x\sim_{\Phi,P,\lambda}y
\quad\Longleftrightarrow\quad
I_{\Phi,\lambda}^P(x)=I_{\Phi,\lambda}^P(y).
\]
\end{definition}

\begin{definition}[Projection quotient]
On note
\[
\pi_\lambda:\mathcal X\longrightarrow \mathcal X/{\sim_{\Phi,P,\lambda}}
\]
la projection quotient.
\end{definition}

\begin{definition}[Image observationnelle réduite]
On pose
\[
\mathcal L_{\lambda,P}:=I_{\Phi,\lambda}^P(\mathcal X)\subseteq \mathbb R^{\mathcal F_\lambda}.
\]
Cette image est la représentation réduite des objets à l'échelle \(\lambda\).
\end{definition}

\begin{proposition}[Raffinement par intensité]
Si \(\lambda\preceq\mu\), alors
\[
x\sim_{\Phi,P,\mu}y \Longrightarrow x\sim_{\Phi,P,\lambda}y.
\]
En particulier, l'augmentation de l'intensité observationnelle ne peut qu'affiner les classes d'indiscernabilité.
\end{proposition}

\begin{proof}
L'égalité des empreintes sur \(\mathcal F_\mu\) implique, par restriction, l'égalité sur la sous-famille \(\mathcal F_\lambda\subseteq\mathcal F_\mu\).
\end{proof}

\begin{corollary}[Segments initiaux de non-séparation]
Supposons \(\Lambda\) totalement ordonné. Pour tout \(x,y\in\mathcal X\), l'ensemble
\[
S_{x,y}:=\{\lambda\in\Lambda:\ x\sim_{\Phi,P,\lambda}y\}
\]
est un segment initial de \(\Lambda\).
\end{corollary}

\begin{proof}
C'est une reformulation immédiate de la proposition précédente.
\end{proof}

\begin{definition}[Seuil de résolution]
Supposons \(\Lambda\) totalement ordonné et \(\iota\) injective. Pour tout couple \((x,y)\in\mathcal X^2\), on définit le seuil de résolution par
\[
\rho_P(x,y):=\sup\{\iota(\lambda):\ \lambda\in S_{x,y}\},
\]
lorsque cette borne supérieure existe. On appelle \emph{seuil séparant} de \((\mathcal X,\Phi,P)\) la quantité
\[
\rho_P^\ast(\Phi):=\inf\{\iota(\lambda):\ I_{\Phi,\lambda}^P \text{ est injective}\},
\]
lorsqu'elle existe.
\end{definition}

\begin{proposition}[Lecture par seuil et borne de résolution]
Sous les hypothèses de la définition précédente, la relation exacte est
\[
x\sim_{\Phi,P,\lambda}y
\quad\Longleftrightarrow\quad
\lambda\in S_{x,y}.
\]
Si
\[
\rho_P(x,y)=\sup\{\iota(\lambda):\lambda\in S_{x,y}\},
\]
alors, pour \(S_{x,y}\) non vide,
\[
\iota(\lambda)<\rho_P(x,y)\Longrightarrow x\sim_{\Phi,P,\lambda}y,
\qquad
\iota(\lambda)>\rho_P(x,y)\Longrightarrow x\not\sim_{\Phi,P,\lambda}y.
\]
Au niveau critique \(\iota(\lambda)=\rho_P(x,y)\), l'appartenance dépend de l'atteinte du supremum et ne peut pas être déduite du seul seuil sans hypothèse supplémentaire sur le bord.
\end{proposition}

\begin{proof}
L'ensemble \(S_{x,y}\) est un segment initial par monotonie. Si \(\iota(\lambda)<\sup\iota(S_{x,y})\), il existe un élément de \(\iota(S_{x,y})\) strictement supérieur à \(\iota(\lambda)\), et l'initialité impose alors \(\iota(\lambda)\in\iota(S_{x,y})\). Si \(\iota(\lambda)>\sup\iota(S_{x,y})\), l'appartenance est impossible. Le cas d'égalité est précisément le cas de bord et requiert de savoir si le supremum appartient à \(\iota(S_{x,y})\).
\end{proof}

\begin{remark}
La distinction micro/macro n'est donc pas fixée par la taille de l'objet. Elle est fixée par l'intensité observationnelle disponible : deux objets de taille très différente peuvent être indistinguables à basse intensité, et un même objet peut devenir richement différencié à haute intensité.
\end{remark}

\subsection{Réduction observationnelle des sondes}

Le résultat suivant formalise l'élimination exacte des sondes redondantes sans changement de fibre.

\begin{definition}[Équivalence observationnelle des sondes]
Fixons \(P\in\mathcal P\) et \(\lambda\in\Lambda\). Pour \(f,g\in\mathcal F_\lambda\), on écrit
\[
f\equiv_{P,\lambda} g
\quad\Longleftrightarrow\quad
\Phi(P,x,f)=\Phi(P,x,g)\ \text{ pour tout }x\in\mathcal X.
\]
\end{definition}

\begin{definition}[Fermeture observationnelle]
Pour toute famille \(\mathcal G\subseteq\mathcal F_\lambda\), on définit
\[
\operatorname{Cl}_{P,\lambda}(\mathcal G)
=
\Bigl\{
f\in\mathcal F_\lambda:\ \exists\,\psi_f:\mathbb R^{\mathcal G}\to\mathbb R,\ 
\Phi(P,x,f)=\psi_f\bigl(I_{\Phi,\mathcal G}^P(x)\bigr)\ \forall x\in\mathcal X
\Bigr\}.
\]
\end{definition}

\begin{remark}[Intuition]
Deux sondes sont redondantes lorsqu'elles ne distinguent jamais des fibres différentes. La fermeture observationnelle décrit l'ensemble maximal des sondes qui n'ajoutent aucune information classificatoire.
\end{remark}

\begin{definition}[Quotient des sondes et information effective]
On pose
\[
\mathcal F_{\lambda,P}^{\sharp}:=\mathcal F_\lambda/\!\equiv_{P,\lambda}.
\]
La fonction induite \(\bar\Phi\) est définie par
\[
\bar\Phi(P,x,[f])=\Phi(P,x,f).
\]
La quantité
\[
C_{\mathrm{eff},P}(\lambda):=\bigl|\mathcal F_{\lambda,P}^{\sharp}\bigr|
\]
est l'information accessible \emph{effective} à l'échelle \(\lambda\).
\end{definition}

\begin{theorem}[Réduction observationnelle des sondes]
Fixons \(P\in\mathcal P\) et \(\lambda\in\Lambda\). Alors :
\begin{enumerate}[label=\arabic*.]
\item \(\equiv_{P,\lambda}\) est une relation d'équivalence sur \(\mathcal F_\lambda\).
\item La fermeture \(\operatorname{Cl}_{P,\lambda}\) est extensive, monotone et idempotente.
\item Pour toute famille \(\mathcal G\subseteq\mathcal F_\lambda\), les empreintes \(I_{\Phi,\mathcal G}^P\) et \(I_{\Phi,\operatorname{Cl}_{P,\lambda}(\mathcal G)}^P\) ont exactement les mêmes fibres.
\item Si \(\mathcal H\subseteq\mathcal F_\lambda\) a les mêmes fibres que \(\mathcal G\), alors \(\mathcal H\subseteq \operatorname{Cl}_{P,\lambda}(\mathcal G)\). Autrement dit, \(\operatorname{Cl}_{P,\lambda}(\mathcal G)\) est la plus grande famille de sondes qui ne change pas la structure de fibres de \(\mathcal G\).
\item Le quotient \(\mathcal F_{\lambda,P}^{\sharp}\) porte une fonction induite \(\bar\Phi\) bien définie, et la cardinalité \(C_{\mathrm{eff},P}(\lambda)\) mesure le nombre de directions de réponse effectivement distinctes à l'échelle \(\lambda\).
\end{enumerate}
\end{theorem}

\begin{proof}
(1) La réflexivité, la symétrie et la transitivité sont immédiates par égalité des réponses.

(2) Si \(\mathcal G\subseteq\mathcal H\), toute factorisation via \(I_{\Phi,\mathcal G}^P\) se compose avec la projection \(\mathbb R^{\mathcal H}\to\mathbb R^{\mathcal G}\), donc \(\operatorname{Cl}_{P,\lambda}(\mathcal G)\subseteq \operatorname{Cl}_{P,\lambda}(\mathcal H)\). L'extensivité est le cas \(\mathcal G\subseteq \operatorname{Cl}_{P,\lambda}(\mathcal G)\). Pour l'idempotence, si \(f\in \operatorname{Cl}_{P,\lambda}(\operatorname{Cl}_{P,\lambda}(\mathcal G))\), alors chaque coordonnée de \(I_{\Phi,\operatorname{Cl}_{P,\lambda}(\mathcal G)}^P\) factorise déjà par \(I_{\Phi,\mathcal G}^P\); la composition donne donc une factorisation de \(f\) par \(I_{\Phi,\mathcal G}^P\).

(3) Si \(I_{\Phi,\mathcal G}^P(x)=I_{\Phi,\mathcal G}^P(y)\), alors pour tout \(f\in\operatorname{Cl}_{P,\lambda}(\mathcal G)\) on a \(\Phi(P,x,f)=\Phi(P,y,f)\) par définition de la fermeture. Ainsi les fibres de \(I_{\Phi,\mathcal G}^P\) sont contenues dans celles de \(I_{\Phi,\operatorname{Cl}_{P,\lambda}(\mathcal G)}^P\). L'inclusion réciproque vient de \(\mathcal G\subseteq\operatorname{Cl}_{P,\lambda}(\mathcal G)\).

(4) Si \(\mathcal H\) a les mêmes fibres que \(\mathcal G\), alors chaque \(h\in\mathcal H\) est constant sur les fibres de \(I_{\Phi,\mathcal G}^P\), donc \(h\in \operatorname{Cl}_{P,\lambda}(\mathcal G)\).

(5) Si \(f\equiv_{P,\lambda}g\), alors leurs profils de réponses coïncident, donc \([f]\mapsto \Phi(P,x,f)\) est bien définie. La quantité \(C_{\mathrm{eff},P}(\lambda)\) compte les directions réellement distinctes de réponse ; elle est indépendante de toute redondance interne.
\end{proof}

\begin{definition}[Dimension effective d'une famille]
Soit \(\mathcal G\subseteq\mathcal F\) une famille finie. On définit
\[
d_{\mathrm{eff},P}(\mathcal G):=\dim \Aff\!\bigl(I_{\Phi,\mathcal G}^P(\mathcal X)\bigr),
\]
lorsque cette dimension est finie. En particulier,
\[
d_{\mathrm{eff},P}(\lambda):=d_{\mathrm{eff},P}(\mathcal F_\lambda).
\]
\end{definition}

\begin{proposition}[Monotonie et invariance de la dimension effective]
Soient \(\mathcal G_1,\mathcal G_2\subseteq\mathcal F\) deux familles finies.
\begin{enumerate}[label=\arabic*.]
\item Si \(\mathcal G_1\subseteq \mathcal G_2\), alors
\[
d_{\mathrm{eff},P}(\mathcal G_1)\le d_{\mathrm{eff},P}(\mathcal G_2).
\]
\item Si \(\sigma:\mathcal G_1\to\mathcal G_2\) est une bijection qui ne fait que re-indexer les sondes sans changer les réponses coordonnées, alors
\[
d_{\mathrm{eff},P}(\mathcal G_1)=d_{\mathrm{eff},P}(\mathcal G_2).
\]
\item Si \(\mathcal H\) s'obtient à partir de \(\mathcal G\) en ajoutant seulement des sondes qui sont des combinaisons affines des coordonnées de \(I_{\Phi,\mathcal G}^P\), alors
\[
d_{\mathrm{eff},P}(\mathcal H)=d_{\mathrm{eff},P}(\mathcal G).
\]
\end{enumerate}
\end{proposition}

\begin{proof}
(1) La projection \(\mathbb R^{\mathcal G_2}\to \mathbb R^{\mathcal G_1}\) envoie \(I_{\Phi,\mathcal G_2}^P(\mathcal X)\) sur \(I_{\Phi,\mathcal G_1}^P(\mathcal X)\), donc la dimension affine ne peut pas décroître quand on passe à un plus grand nombre de coordonnées.

(2) Un re-indexage des sondes correspond à une permutation des coordonnées de l'espace ambiant, donc à un automorphisme affine. La dimension affine de l'image est inchangée.

(3) Si les coordonnées ajoutées sont des combinaisons affines des précédentes, l'image de \(\mathcal H\) est contenue dans une enveloppe affine de même dimension que celle de \(\mathcal G\), et réciproquement \(\mathcal G\subseteq\mathcal H\). Les deux dimensions coïncident.
\end{proof}

\begin{proposition}[Borne dimensionnelle intrinsèque]
Si \(C_{\mathrm{eff},P}(\lambda)<\infty\), alors
\[
d_{\mathrm{eff},P}(\lambda)\le C_{\mathrm{eff},P}(\lambda).
\]
\end{proposition}

\begin{proof}
Après réduction observationnelle, \(\mathcal L_{\lambda,P}=I_{\Phi,\lambda}^P(\mathcal X)\) s'inscrit dans un espace de coordonnées de dimension \(C_{\mathrm{eff},P}(\lambda)\). L'enveloppe affine d'un sous-ensemble d'un espace de dimension \(m\) ne peut dépasser cette dimension.
\end{proof}

\begin{proposition}[Géométrie induite]
Fixons \(P\in\mathcal P\) et \(\lambda\in\Lambda\). Soit \(\|\cdot\|\) une norme sur \(\mathbb R^{\mathcal F_\lambda}\) lorsque \(\mathcal F_\lambda\) est finie. On définit
\[
d_{\Phi,P,\lambda}(x,y)=\|I_{\Phi,\lambda}^P(x)-I_{\Phi,\lambda}^P(y)\|,
\qquad x,y\in\mathcal X.
\]
Alors \(d_{\Phi,P,\lambda}\) est une pseudométrique sur \(\mathcal X\). De plus,
\[
I_{\Phi,\lambda}^P:(\mathcal X,d_{\Phi,P,\lambda})\longrightarrow \mathcal L_{\lambda,P}
\]
est une isométrie surjective. Si \(\mathcal F_\lambda\) est séparante en \(P\), alors \(d_{\Phi,P,\lambda}\) est une métrique et \(I_{\Phi,\lambda}^P\) est une bijection isométrique entre \((\mathcal X,d_{\Phi,P,\lambda})\) et \(\mathcal L_{\lambda,P}\).
\end{proposition}

\begin{proof}
Les axiomes d'une norme donnent la positivité, la symétrie et l'inégalité triangulaire. La surjectivité sur l'image est tautologique. L'injectivité suit de la proposition précédente lorsque \(\mathcal F_\lambda\) est séparante.
\end{proof}

\subsection{Théorème de factorisation}

\begin{theorem}[Factorisation de l'empreinte]
Fixons \(P\in\mathcal P\) et \(\lambda\in\Lambda\). Soient \(Y\) un ensemble quelconque et \(T:\mathcal X\to Y\) une application. Les assertions suivantes sont équivalentes :
\begin{enumerate}[label=\textbf{(\arabic*)}]
\item \(T\) est constante sur les fibres de \(I_{\Phi,\lambda}^P\), c'est-à-dire
\[
I_{\Phi,\lambda}^P(x)=I_{\Phi,\lambda}^P(y)\Longrightarrow T(x)=T(y).
\]
\item Il existe une unique application \(\widetilde T:\mathcal L_{\lambda,P}\to Y\) telle que
\[
T=\widetilde T\circ I_{\Phi,\lambda}^P.
\]
\item Il existe une unique application \(\widehat T:\mathcal X/{\sim_{\Phi,P,\lambda}}\to Y\) telle que
\[
T=\widehat T\circ \pi_\lambda.
\]
\item Si \(J:\mathcal X\to Z\) est une application ayant exactement les mêmes fibres que \(I_{\Phi,\lambda}^P\), alors il existe une unique bijection \(\beta:\mathcal L_{\lambda,P}\to J(\mathcal X)\) telle que
\[
J=\beta\circ I_{\Phi,\lambda}^P.
\]
\end{enumerate}

En particulier, le triplet \(\bigl(I_{\Phi,\lambda}^P,\pi_\lambda,\mathcal L_{\lambda,P}\bigr)\) décrit une même donnée observationnelle à travers trois incarnations naturellement équivalentes. Toute représentation minimale fidèle de l'information à l'échelle \(\lambda\) est unique à bijection près.
\end{theorem}

\begin{proof}
L'implication \((1)\Rightarrow(2)\) est immédiate en définissant \(\widetilde T\bigl(I_{\Phi,\lambda}^P(x)\bigr):=T(x)\). L'hypothèse garantit que cette définition est indépendante du représentant. L'implication \((2)\Rightarrow(3)\) suit du quotient par les fibres de \(I_{\Phi,\lambda}^P\). L'implication \((1)\Rightarrow(4)\) est le principe de minimalité fibrée. Les réciproques sont immédiates. L'unicité dans chaque cas est imposée par la définition de la factorisation sur une image ou un quotient.
\end{proof}

\begin{remark}[Pont vers la Partie II]
Toute structure construite uniquement à partir des fibres observables se factorise par les empreintes et les quotients associés. La Partie II montrera que, lorsque l'on organise ces fibres en niveaux de seuil, la fonction de seuil complète suffit à reconstruire toute la structure observationnelle.
\end{remark}

\subsection{Grand exemple développé : classification observée d'une base de données finie}

On considère six objets
\[
\mathcal X=\{x_1,x_2,x_3,x_4,x_5,x_6\},
\]
et quatre sondes \(f_1,f_2,f_3,f_4\). Les réponses sont données par le tableau suivant :
\[
\begin{array}{c|cccc}
& f_1 & f_2 & f_3 & f_4\\ \hline
x_1 & 0 & 0 & 0 & 0\\
x_2 & 0 & 0 & 1 & 0\\
x_3 & 1 & 0 & 0 & 0\\
x_4 & 1 & 1 & 0 & 0\\
x_5 & 2 & 0 & 0 & 0\\
x_6 & 2 & 0 & 0 & 1
\end{array}
\]

On fixe une échelle discrète
\[
\Lambda=\{\lambda_1\prec\lambda_2\prec\lambda_3\prec\lambda_4\},
\qquad
\iota(\lambda_k)=k.
\]
On prend
\[
\mathcal F_{\lambda_1}=\{f_1\},\quad
\mathcal F_{\lambda_2}=\{f_1,f_2\},\quad
\mathcal F_{\lambda_3}=\{f_1,f_2,f_3\},\quad
\mathcal F_{\lambda_4}=\{f_1,f_2,f_3,f_4\}.
\]

Les partitions obtenues sont alors :
\[
\mathcal X/{\sim_{\Phi,P,\lambda_1}}=
\bigl\{\{x_1,x_2\},\{x_3,x_4\},\{x_5,x_6\}\bigr\},
\]
\[
\mathcal X/{\sim_{\Phi,P,\lambda_2}}=
\bigl\{\{x_1,x_2\},\{x_3\},\{x_4\},\{x_5,x_6\}\bigr\},
\]
\[
\mathcal X/{\sim_{\Phi,P,\lambda_3}}=
\bigl\{\{x_1\},\{x_2\},\{x_3\},\{x_4\},\{x_5,x_6\}\bigr\},
\]
\[
\mathcal X/{\sim_{\Phi,P,\lambda_4}}=
\bigl\{\{x_1\},\{x_2\},\{x_3\},\{x_4\},\{x_5\},\{x_6\}\bigr\}.
\]

La hiérarchie est donc
\[
\begin{aligned}
\{x_1,\ldots,x_6\}
&\succ \{x_1,x_2\}\,|\,\{x_3,x_4\}\,|\,\{x_5,x_6\}\\
&\succ \{x_1,x_2\}\,|\,\{x_3\}\,|\,\{x_4\}\,|\,\{x_5,x_6\}\\
&\succ \{x_1\}\,|\,\{x_2\}\,|\,\{x_3\}\,|\,\{x_4\}\,|\,\{x_5,x_6\}\\
&\succ \text{singletons.}
\end{aligned}
\]

Une fonction de seuil complète compatible avec cette hiérarchie est donnée par
\[
\Theta_P(\{x_1,x_2\})=2.5,\qquad
\Theta_P(\{x_3,x_4\})=1.5,\qquad
\Theta_P(\{x_5,x_6\})=3.5,
\]
et pour toute paire de points appartenant à des blocs distincts dès l'échelle \(\lambda_1\),
\[
\Theta_P(\{x_i,x_j\})=0.5.
\]
Cette définition satisfait immédiatement la loi du minimum double : dans tout triplet, le plus petit des seuils est \(0.5\) ou, si deux points appartiennent au même bloc de niveau supérieur, le minimum est atteint au moins deux fois par les deux paires qui croisent ce bloc.

Le spectre critique est alors
\[
\Sigma_P=\{0.5,\,1.5,\,2.5,\,3.5\},
\]
et la mesure empirique vaut
\[
\nu_P
=
\frac{12}{15}\,\delta_{0.5}
+\frac{1}{15}\,\delta_{1.5}
+\frac{1}{15}\,\delta_{2.5}
+\frac{1}{15}\,\delta_{3.5}.
\]
En particulier, la courbe de résolution atomique est
\[
\widetilde{\mathcal R}_P(t)=
\begin{cases}
0, & t\le 0.5,\\[0.3em]
\frac{12}{15}, & 0.5<t\le 1.5,\\[0.3em]
\frac{13}{15}, & 1.5<t\le 2.5,\\[0.3em]
\frac{14}{15}, & 2.5<t\le 3.5,\\[0.3em]
1, & t>3.5.
\end{cases}
\]
Les valeurs de résolution aux échelles discrètes sont donc
\[
\mathcal R_P(\lambda_1)=\frac{12}{15},\qquad
\mathcal R_P(\lambda_2)=\frac{13}{15},\qquad
\mathcal R_P(\lambda_3)=\frac{14}{15},\qquad
\mathcal R_P(\lambda_4)=1.
\]
Ce calcul montre simultanément trois phénomènes :
\begin{itemize}[leftmargin=1.4em]
\item l'information brute augmente à chaque niveau ;
\item dans le tableau choisi, l'information effective augmente jusqu'au dernier niveau ; si l'une des sondes finales est redondante, le dernier niveau peut au contraire devenir un plateau ;
\item la résolution n'augmente que par sauts atomiques visibles sur la mesure des seuils.
\end{itemize}
L'exemple fournit ainsi une lecture complète de la théorie : empreinte, quotient, seuil, arbre, mesure et résolution y sont calculés explicitement.

\begin{definition}[Donnée de bord des seuils]
Pour une paire distincte \(\{x,y\}\), on définit
\[
b_P(\{x,y\})=1
\quad\Longleftrightarrow\quad
\exists\lambda\in S_{x,y}:\ \iota(\lambda)=\Theta_P(\{x,y\}),
\]
et \(b_P(\{x,y\})=0\) sinon. La donnée de seuil complétée est
\[
H_P:=(\Theta_P,b_P).
\]
Elle distingue le seuil-squelette de son statut de bord. Pour \(x\ne y\), on a alors l'identité exacte
\[
x\sim_{\Phi,P,\lambda}y
\Longleftrightarrow
\iota(\lambda)<\Theta_P(\{x,y\})
\ \text{ ou }\ 
\bigl(\iota(\lambda)=\Theta_P(\{x,y\})\land b_P(\{x,y\})=1\bigr).
\]
\end{definition}

\begin{remark}[Figure source]
Le manuscrit source contient ici une figure illustrative ; sa présence est documentée dans le corpus, mais l'image externe n'est pas embarquée dans cette reconstruction.
\end{remark}

\subsection{Exemples complémentaires}

\begin{proposition}[Exemple élémentaire de séparation]
Soit \(\mathcal X=\mathbb R^2\), \(\mathcal F=\{f_1,f_2,f_3\}\), et définissons
\[
\Phi(P,(u,v),f_1)=u,\qquad
\Phi(P,(u,v),f_2)=v,\qquad
\Phi(P,(u,v),f_3)=u+v.
\]
Alors la famille \(\{f_1\}\) n'est pas séparante, la famille \(\{f_1,f_2\}\) est séparante, et l'ajout de \(f_3\) est redondant.
\end{proposition}

\begin{proof}
Avec \(\{f_1\}\), tous les points de même première coordonnée sont indiscernables. Avec \(\{f_1,f_2\}\), l'empreinte est l'identité sur \(\mathbb R^2\), donc injective. Enfin \(f_3\) est déterminé par \(f_1\) et \(f_2\), donc il ne raffine rien.
\end{proof}

\begin{proposition}[Graphes finis]
Soit \(\mathcal X\) l'ensemble des graphes simples sur l'ensemble des sommets \(\{1,2,3,4\}\). Pour chaque arête \(e\), on définit une sonde
\[
f_e(G)=\mathbf 1_{\{e\in E(G)\}},
\]
et l'on ajoute une sonde
\[
g(G)=\#\{\text{triangles de }G\}.
\]
Alors la famille des sondes \(f_e\) est séparante, et la sonde \(g\) appartient à la fermeture observationnelle de cette famille.
\end{proposition}

\begin{proof}
La donnée de toutes les sondes \(f_e\) reconstitue exactement l'ensemble des arêtes du graphe, donc le graphe lui-même ; la famille est donc séparante. La quantité \(g(G)\) s'exprime comme une fonction polynomiale des variables \(f_e(G)\), puisqu'un triangle est déterminé par ses trois arêtes ; elle dépend donc uniquement des valeurs de la famille séparante.
\end{proof}

\begin{proposition}[Signaux et images]
Soit \(\mathcal X\) la famille des signaux
\[
u_{a,b,c}(t)=a+b\cos(2\pi t)+c\sin(2\pi t),
\qquad (a,b,c)\in\mathbb R^3.
\]
On considère les trois sondes d'évaluation
\[
f_1(u)=u(0),\qquad f_2(u)=u\!\left(\tfrac14\right),\qquad f_3(u)=u\!\left(\tfrac12\right).
\]
Alors l'empreinte
\[
u\longmapsto \bigl(u(0),u(\tfrac14),u(\tfrac12)\bigr)
\]
est injective. En particulier, la théorie s'applique à des familles de fonctions, de signaux et d'images numériques sur une grille finie.
\end{proposition}

\begin{proof}
On a
\[
\begin{pmatrix}
u(0)\\[0.2em]
u\!\left(\tfrac14\right)\\[0.2em]
u\!\left(\tfrac12\right)
\end{pmatrix}
=
\begin{pmatrix}
1&1&0\\
1&0&1\\
1&-1&0
\end{pmatrix}
\begin{pmatrix}
a\\b\\c
\end{pmatrix},
\]
et la matrice affichée est inversible. Les coefficients \((a,b,c)\) sont donc déterminés de manière unique par les trois valeurs d'empreinte.
\end{proof}

\section{Partie II. Géométrie des seuils}

\subsection{Spectre critique, courbe de résolution et lecture par seuil}

Fixons \(P\in\mathcal P\). Supposons \(\mathcal X\) fini, \(\Lambda\) totalement ordonné et \(\iota:\Lambda\to\mathbb R_+\) injective.

\begin{definition}[Fonction de seuil complète]
On note
\[
\Theta_P:\binom{\mathcal X}{2}\longrightarrow \mathbb R_+,
\qquad
\Theta_P(\{x,y\})=\rho_P(x,y).
\]
Cette application encode, pour chaque paire non ordonnée \(\{x,y\}\), l'intensité critique à partir de laquelle la paire devient séparée.
\end{definition}

\begin{definition}[Spectre critique d'un contexte]
On appelle \emph{spectre critique} de \(P\) l'ensemble
\[
\Sigma_P:=\{\Theta_P(\{x,y\}):\ x,y\in\mathcal X,\ x\neq y\}\subseteq \mathbb R_+.
\]
\end{definition}

\begin{definition}[Courbe de résolution observée]
On définit pour tout \(\lambda\in\Lambda\) la courbe de résolution par
\[
\mathcal R_P(\lambda)
:=
\frac{1}{\binom{|\mathcal X|}{2}}
\sum_{\{x,y\}\subseteq\mathcal X}
\mathbf 1_{\{\iota(\lambda)>\Theta_P(\{x,y\})\}}.
\]
Elle mesure la proportion de paires d'objets déjà distinguées à l'échelle \(\lambda\).
\end{definition}

\begin{remark}
Le spectre critique \(\Sigma_P\) n'est que l'image de \(\Theta_P\), sans multiplicité ni localisation par paire. Il ne suffit donc pas, à lui seul, à reconstruire la hiérarchie observationnelle ; la donnée complète est la fonction \(\Theta_P\).
\end{remark}

\begin{proposition}[Lecture par seuil avec traitement du bord]
Sous les hypothèses précédentes, la relation exacte est
\[
x\sim_{\Phi,P,\lambda}y
\quad\Longleftrightarrow\quad
\lambda\in S_{x,y},
\qquad
S_{x,y}:=\{\lambda\in\Lambda:\ x\sim_{\Phi,P,\lambda}y\}.
\]
En posant
\[
\Theta_P(\{x,y\})=\sup\{\iota(\lambda):\lambda\in S_{x,y}\},
\]
on obtient, pour toute paire dont \(S_{x,y}\) est non vide,
\[
\iota(\lambda)<\Theta_P(\{x,y\})\ \Longrightarrow\ x\sim_{\Phi,P,\lambda}y,
\]
\[
\iota(\lambda)>\Theta_P(\{x,y\})\ \Longrightarrow\ x\not\sim_{\Phi,P,\lambda}y.
\]
Au niveau critique \(\iota(\lambda)=\Theta_P(\{x,y\})\), l'appartenance de \(\lambda\) à \(S_{x,y}\) dépend de l'atteinte ou non de la borne supérieure ; elle ne peut donc pas être déduite du seul supremum sans hypothèse supplémentaire d'atteinte (par exemple fermeture appropriée de l'ensemble admissible).
\end{proposition}

\begin{proof}
Par monotonie, \(S_{x,y}\) est un segment initial de \(\Lambda\), donc son image \(\iota(S_{x,y})\) est un segment initial de l'image ordonnée \(\iota(\Lambda)\). Si \(\iota(\lambda)<\sup\iota(S_{x,y})\), il existe un élément de \(\iota(S_{x,y})\) strictement supérieur à \(\iota(\lambda)\), et la propriété de segment initial impose alors \(\iota(\lambda)\in\iota(S_{x,y})\). Inversement, si \(\iota(\lambda)>\sup\iota(S_{x,y})\), alors \(\iota(\lambda)\notin\iota(S_{x,y})\). L'égalité est précisément le cas de bord : elle équivaut à l'appartenance du point critique à \(\iota(S_{x,y})\), ce que le supremum seul ne garantit pas.
\end{proof}

\begin{remark}[Figure source]
Le manuscrit source contient ici une figure illustrative ; sa présence est documentée dans le corpus, mais l'image externe n'est pas embarquée dans cette reconstruction.
\end{remark}

\subsection{Pourquoi la loi du minimum double généralise l'ultramétrie}

\begin{proposition}[Cohérence des seuils]
Soit \(K:\binom{\mathcal X}{2}\to\mathbb R_+\) une application symétrique. Les assertions suivantes sont équivalentes :
\begin{enumerate}[label=\arabic*.]
\item Pour tout \(t\in\mathbb R_+\), la relation
\[
x\equiv_t^K y \Longleftrightarrow x=y \text{ ou } K(\{x,y\})\ge t
\]
est transitive.
\item \(K\) satisfait la loi du minimum double
\[
K(\{x,z\})\ge \min\bigl(K(\{x,y\}),K(\{y,z\})\bigr)
\qquad
\text{pour tous }x,y,z\in\mathcal X\text{ deux à deux distincts}.
\]
\end{enumerate}
\end{proposition}

\begin{proof}
Supposons (1). Fixons \(x,y,z\). En prenant \(t=\min(K(\{x,y\}),K(\{y,z\}))\), on a \(x\equiv_t^K y\) et \(y\equiv_t^K z\); la transitivité donne \(x\equiv_t^K z\), donc \(K(\{x,z\})\ge t\), ce qui est la loi du minimum double.

Réciproquement, supposons la loi du minimum double. Si \(x\equiv_t^K y\) et \(y\equiv_t^K z\), alors \(K(\{x,y\})\ge t\) et \(K(\{y,z\})\ge t\); la loi du minimum double entraîne \(K(\{x,z\})\ge t\), donc \(x\equiv_t^K z\). La transitivité est donc vraie pour tout \(t\).
\end{proof}

\begin{theorem}[Théorème structural des seuils]
Sous les hypothèses précédentes, la fonction de seuil complète \(\Theta_P\) possède les propriétés suivantes.
\begin{enumerate}[label=\textbf{(\Alph*)}]
\item \textbf{Loi du minimum double.} Pour tous \(x,y,z\) deux à deux distincts,
\[
\Theta_P(\{x,z\})\ge \min\bigl(\Theta_P(\{x,y\}),\Theta_P(\{y,z\})\bigr).
\]
Équivalemment, pour toute triple \(\{x,y,z\}\), le minimum des trois seuils
\[
\Theta_P(\{x,y\}),\quad \Theta_P(\{y,z\}),\quad \Theta_P(\{z,x\})
\]
est atteint au moins deux fois.

\item \textbf{Reconstruction de la hiérarchie des partitions hors du bord, et rôle du bord.} Pour tout \(\lambda\in\Lambda\) et tout \(x,y\in\mathcal X\), la relation exacte est
\[
x\sim_{\Phi,P,\lambda}y
\quad\Longleftrightarrow\quad
\lambda\in S_{x,y}.
\]
La fonction \(\Theta_P\) détermine la classification pour les niveaux strictement sous ou strictement au-dessus du seuil. Lorsque \(\iota(\lambda)=\Theta_P(\{x,y\})\), il faut en plus savoir si la borne est atteinte. Ainsi, \(\Theta_P\) seul ne détermine la filtration complète aux points critiques sans hypothèse d'atteinte ou donnée de bord complémentaire. Réciproquement,
\[
\Theta_P(\{x,y\})
=
\sup\{\iota(\lambda):\ \lambda\in S_{x,y}\}.
\]

\item \textbf{Structure arborescente du squelette.} Pour tout \(t\in\mathbb R_+\), la relation superniveau fermée
\[
x\equiv_t
\quad\Longleftrightarrow\quad
x=y\ \text{ ou }\ \Theta_P(\{x,y\})\ge t
\]
est une relation d'équivalence sur \(\mathcal X\). Si \(t_1\le t_2\), alors \(\equiv_{t_2}\) raffine \(\equiv_{t_1}\). La famille des classes
\[
\overline{\mathcal T}_P:=\{[x]_t:\ x\in\mathcal X,\ t\in\mathbb R_+\}
\]
est laminaire : pour tous \(A,B\in\overline{\mathcal T}_P\),
\[
A\cap B\neq\varnothing \Longrightarrow A\subseteq B \text{ ou } B\subseteq A.
\]
En particulier, munie de l'inclusion, \(\overline{\mathcal T}_P\) est un arbre enraciné fini de racine \(\mathcal X\). La famille observée \(\mathcal T_P\) coïncide avec cette construction aux niveaux non critiques ; aux niveaux critiques, la coïncidence exige la condition de bord encodée par \(b_P\). Pour tous \(A,B\in\overline{\mathcal T}_P\), l'ancêtre commun minimal \(A\vee B\) existe et est unique. Si l'on exige que cet ancêtre appartienne à \(\mathcal T_P\), il faut ajouter l'hypothèse que le niveau correspondant est effectivement observé. Si \(A\in\overline{\mathcal T}_P\) n'est pas minimal pour l'inclusion, l'ensemble de ses enfants est non vide, ses enfants sont deux à deux disjoints, et
\[
A=\bigsqcup_{B\in\operatorname{Ch}(A)}B.
\]

\item \textbf{Classification par coupes.} Pour tout \(t\in\mathbb R_+\), les blocs maximaux de la relation superniveau fermée forment la partition du squelette. Pour \(t=\iota(\lambda)\), elle coïncide avec la partition effectivement observée lorsque le seuil critique est atteint selon la convention fermée ; sinon, la différence est localisée au niveau critique. Deux contextes \(P\) et \(Q\) induisent la même hiérarchie de partitions si et seulement si leurs familles de partitions \((\equiv_t)_{t\ge 0}\) sont isomorphes niveau par niveau par une bijection préservant les seuils.
\end{enumerate}

En particulier, toute structure observationnelle construite uniquement à partir des fibres observables à l'échelle est déterminée par \(H_P=(\Theta_P,b_P)\) ; \(\Theta_P\) seul détermine le squelette hors bord.
\end{theorem}

\begin{proof}
(A) Pour \(x,y\in\mathcal X\), posons
\[
S_{x,y}:=\{\lambda\in\Lambda:\ x\sim_{\Phi,P,\lambda}y\}.
\]
Par la définition des relations \(\sim_{\Phi,P,\lambda}\), chaque \(S_{x,y}\) est un segment initial de \(\Lambda\). De plus, pour tout \(\lambda\in\Lambda\),
\[
x\sim_{\Phi,P,\lambda}y\ \text{ et }\ y\sim_{\Phi,P,\lambda}z
\Longrightarrow
x\sim_{\Phi,P,\lambda}z,
\]
donc
\[
S_{x,y}\cap S_{y,z}\subseteq S_{x,z}.
\]
En appliquant \(\iota\) et en prenant les bornes supérieures, on obtient
\[
\Theta_P(\{x,z\})
=
\sup \iota(S_{x,z})
\ge
\sup \iota(S_{x,y}\cap S_{y,z})
=
\min\bigl(\Theta_P(\{x,y\}),\Theta_P(\{y,z\})\bigr),
\]
ce qui est la loi du minimum double. La formulation équivalente sur les triples est immédiate.

(B) La relation exacte est \(x\sim_{\Phi,P,\lambda}y \Longleftrightarrow \lambda\in S_{x,y}\). La quantité \(\Theta_P(\{x,y\})\) est la borne supérieure de \(\iota(S_{x,y})\). Elle détermine la classification aux niveaux strictement sous et strictement au-dessus du seuil. Au niveau critique, l'appartenance dépend de l'atteinte du supremum ; une donnée de bord ou une hypothèse d'atteinte est donc nécessaire pour reconstruire exactement la filtration complète. Réciproquement, la famille des segments \(S_{x,y}\) détermine \(\Theta_P\) par prise du supremum.

(C) La réflexivité de \(\equiv_t\) est imposée par le cas \(x=y\), et la symétrie est évidente. Pour la transitivité, si \(x\equiv_t y\) et \(y\equiv_t z\), alors
\[
\Theta_P(\{x,y\})\ge t
\quad\text{et}\quad
\Theta_P(\{y,z\})\ge t.
\]
Par la loi du minimum double,
\[
\Theta_P(\{x,z\})\ge t,
\]
d'où \(x\equiv_t z\). Ainsi \(\equiv_t\) est une relation d'équivalence. Si \(t_1\le t_2\), alors \(\equiv_{t_2}\) est plus fine que \(\equiv_{t_1}\) car la condition \(\Theta_P(\{x,y\})\ge t_2\) est plus restrictive que \(\Theta_P(\{x,y\})\ge t_1\).

Si \(A=[x]_{t_1}\) et \(B=[y]_{t_2}\) intersectent, choisissons \(z\in A\cap B\). Si \(t_1\le t_2\), alors \(\equiv_{t_2}\) est plus fine que \(\equiv_{t_1}\), donc
\[
[z]_{t_2}\subseteq [z]_{t_1},
\]
d'où \(B\subseteq A\). Le cas \(t_2\le t_1\) est symétrique. La famille \(\overline{\mathcal T}_P\) est donc laminaire.

Comme \(\mathcal X\) est fini, il n'existe qu'un nombre fini de sous-ensembles de \(\mathcal X\), donc \(\overline{\mathcal T}_P\) et \(\mathcal T_P\) sont finies. Le plus petit ancêtre commun de deux éléments existe alors par le même argument de finitude et de laminarité que précédemment : l'ensemble des ancêtres communs est non vide car il contient \(\mathcal X\), et deux minimaux ne peuvent être distincts qu'en se recouvrant, ce qui est interdit par la laminarité. Le join \(A\vee B\) est donc bien défini et unique dans \(\mathcal T_P\).

Enfin, si \(A\in\overline{\mathcal T}_P\) n'est pas minimal, le fait que \(\mathcal X\) soit fini garantit qu'il existe un seuil suffisamment grand pour lequel toutes les classes sont des singletons ; ainsi, pour tout \(x\in A\), l'ensemble
\[
\mathcal S_x(A):=\{C\in\overline{\mathcal T}_P:\ x\in C\subsetneq A\}
\]
est non vide et fini. Il admet donc un élément maximal, que l'on note \(B_x\). Si
\[
B_x\subsetneq C\subsetneq A,
\]
alors \(x\in C\), donc \(C\in\mathcal S_x(A)\), ce qui contredit la maximalité de \(B_x\). Ainsi \(B_x\) est un enfant de \(A\), et les enfants étant deux à deux disjoints par laminarité on obtient
\[
A=\bigsqcup_{B\in\operatorname{Ch}(A)}B.
\]

(D) Les blocs maximaux de la relation \(\equiv_t\) sont précisément ses classes d'équivalence. À un niveau critique \(t=\iota(\lambda)\), cette partition coïncide avec \(\mathcal X/{\sim_{\Phi,P,\lambda}}\) seulement lorsque le bord est atteint selon la convention déclarée. Si deux contextes \(P\) et \(Q\) ont la même hiérarchie de partitions, alors ils ont les mêmes classes à chaque niveau de seuil, donc les mêmes blocs dans leurs arbres complétés. L'identité sur les blocs définit alors un isomorphisme d'arbres étiquetés. Réciproquement, un isomorphisme d'arbres étiquetés préserve les blocs maximaux à chaque seuil, donc les partitions. La hiérarchie est donc entièrement déterminée par l'arbre étiqueté.
\end{proof}

\begin{corollary}[Caractérisation des fonctions de seuil admissibles]
Une application symétrique
\[
K:\binom{\mathcal X}{2}\longrightarrow \mathbb R_+
\]
est la fonction de seuil complète d'une hiérarchie observationnelle si et seulement si elle vérifie la loi du minimum double
\[
K(\{x,z\})\ge \min\bigl(K(\{x,y\}),K(\{y,z\})\bigr)
\qquad
\text{pour tous }x,y,z\in\mathcal X\text{ deux à deux distincts}.
\]
Dans ce cas, les relations superniveau
\[
x\equiv_t^K y
\quad\Longleftrightarrow\quad
x=y\ \text{ ou }\ K(\{x,y\})\ge t
\]
forment une filtration de partitions, et l'arbre étiqueté associé est unique.
\end{corollary}

\begin{proof}
Le sens direct suit du théorème structural des seuils appliqué à \(K=\Theta_P\).

Pour le sens réciproque, supposons que \(K\) satisfasse la loi du minimum double. Pour tout \(t\in\mathbb R_+\), la relation \(\equiv_t^K\) est réflexive et symétrique par définition. Sa transitivité découle de la loi du minimum double : si \(x\equiv_t^K y\) et \(y\equiv_t^K z\), alors
\[
K(\{x,y\})\ge t \quad\text{et}\quad K(\{y,z\})\ge t,
\]
donc
\[
K(\{x,z\})\ge t,
\]
d'où \(x\equiv_t^K z\). Ainsi \(\equiv_t^K\) définit une filtration de partitions.

Réciproquement, pour deux objets \(x,y\), l'ensemble des niveaux \(t\) tels que \(x\equiv_t^K y\) forme un segment initial, dont la borne supérieure est
\[
K(\{x,y\})=\sup\{t\in\mathbb R_+:\ x\equiv_t^K y\}.
\]
Cette formule fixe \(K\) de manière unique à partir de la filtration, et donc fixe également l'arbre étiqueté associé.
\end{proof}

\begin{corollary}[Représentation par treillis de filtrations]
La donnée d'une fonction \(K:\binom{\mathcal X}{2}\to\mathbb R_+\) satisfaisant la loi du minimum double est équivalente à la donnée d'un treillis de partitions totalement ordonné par raffinement, récupérable par les relations \(\equiv_t^K\). En particulier, l'arbre étiqueté construit à partir de \(K\) n'ajoute aucune donnée extra : il représente la même structure de filtrations sous une forme laminaire.
\end{corollary}

\begin{proof}
La filtration \((\equiv_t^K)_{t\in\mathbb R_+}\) est totalement ordonnée par raffinement et le passage de la filtration à l'ensemble des classes maximales fournit un treillis laminaire de partitions. Réciproquement, un tel treillis permet de reconstruire \(K\) par la formule de borne supérieure du théorème précédent.
\end{proof}

\subsection{Représentation atomique de la résolution}

\begin{proposition}[Passage à une ultramétrique]
Supposons que \(K:\binom{\mathcal X}{2}\to(0,+\infty)\) satisfasse la loi du minimum double. Alors
\[
d_K(x,x):=0,\qquad
d_K(x,y):=\exp\!\bigl(-K(\{x,y\})\bigr)\quad(x\neq y)
\]
définit une ultramétrique sur \(\mathcal X\).
\end{proposition}

\begin{proof}
La symétrie est immédiate par symétrie de $K$. La séparation vient du cas diagonal défini par $d_K(x,x)=0$ et de $K(\{x,y\})>0$ pour $x\neq y$. Pour la forte inégalité triangulaire, la loi du minimum double donne
\[
K(\{x,z\})\ge \min\bigl(K(\{x,y\}),K(\{y,z\})\bigr).
\]
En appliquant la fonction décroissante \(u\mapsto e^{-u}\), on obtient
\[
d_K(x,z)\le \max\bigl(d_K(x,y),d_K(y,z)\bigr),
\]
ce qui est précisément l'inégalité ultramétrique.
\end{proof}

\begin{remark}
La loi du minimum double joue donc le même rôle formel que la forte inégalité triangulaire. La différence est conceptuelle : la théorie ne part pas d'une distance, mais d'une fonction de seuil issue des observations. Le dendrogramme n'est pas présupposé ; il est reconstruit à partir des niveaux d'indiscernabilité, dans l'esprit des représentations ultramétriques des hiérarchies observées \textsuperscript{[références bibliographiques du manuscrit source]}.
\end{remark}

\begin{definition}[Mesure empirique des seuils]
Sous l'hypothèse \(|\mathcal X|\ge2\), on définit la mesure empirique des seuils par
\[
\nu_P
:=
\frac{1}{\binom{|\mathcal X|}{2}}
\sum_{\{x,y\}\subseteq\mathcal X}\delta_{\Theta_P(\{x,y\})}.
\]
Son support est exactement \(\Sigma_P\).
\end{definition}

\begin{theorem}[Représentation atomique de la résolution]
Sous l'hypothèse \(|\mathcal X|\ge2\), soit
\[
\widetilde{\mathcal R}_P:\mathbb R_+\longrightarrow [0,1],
\qquad
\widetilde{\mathcal R}_P(t)
:=
\frac{1}{\binom{|\mathcal X|}{2}}
\sum_{\{x,y\}\subseteq\mathcal X}
\mathbf 1_{\{t>\Theta_P(\{x,y\})\}}.
\]
Alors la mesure de Lebesgue--Stieltjes associée à \(\widetilde{\mathcal R}_P\), notée \(\mu_P\), coïncide avec \(\nu_P\), c'est-à-dire
\[
\mu_P([s,t))=\widetilde{\mathcal R}_P(t)-\widetilde{\mathcal R}_P(s)
\qquad\Longleftrightarrow\quad
\mu_P=\nu_P.
\]
De plus, pour toute intensité critique \(\alpha\in\Sigma_P\),
\[
\widetilde{\mathcal R}_P(\alpha^+)-\widetilde{\mathcal R}_P(\alpha^-)
=
\nu_P(\{\alpha\})
=
\frac{1}{\binom{|\mathcal X|}{2}}
\#\bigl\{\{x,y\}\subseteq\mathcal X:\ \Theta_P(\{x,y\})=\alpha\bigr\}.
\]
\end{theorem}

\begin{proof}
Par définition de \(\nu_P\),
\[
\nu_P((-\infty,t))
=
\frac{1}{\binom{|\mathcal X|}{2}}
\sum_{\{x,y\}\subseteq\mathcal X}
\delta_{\Theta_P(\{x,y\})}((-\infty,t))
=
\frac{1}{\binom{|\mathcal X|}{2}}
\sum_{\{x,y\}\subseteq\mathcal X}
\mathbf 1_{\{t>\Theta_P(\{x,y\})\}},
\]
ce qui est exactement \(\widetilde{\mathcal R}_P(t)\). La mesure de Lebesgue--Stieltjes associée à une fonction croissante en escalier est précisément la somme de ses sauts, donc \(\mu_P=\nu_P\).

Le saut en \(\alpha\) ne reçoit une contribution que des paires dont le seuil exact vaut \(\alpha\), d'où la formule annoncée.
\end{proof}

\begin{corollary}[Prédiction de plateau]
Dans toute expérience où l'intensité observationnelle est augmentée de manière monotone, la résolution observée ne peut croître que par paliers. Plus précisément, si
\[
[t_1,t_2]\cap\Sigma_P=\varnothing,
\]
alors \(\widetilde{\mathcal R}_P\) est constante sur \([t_1,t_2]\).

Toute évolution absolument continue de \(\widetilde{\mathcal R}_P\) sur un intervalle contenant un point de \(\Sigma_P\) est incompatible avec la version finie de la théorie.
\end{corollary}

\begin{proof}
Si \([t_1,t_2]\cap\Sigma_P=\varnothing\), alors aucune masse de \(\nu_P\) ne se trouve dans cet intervalle. La fonction de répartition \(\widetilde{\mathcal R}_P\) ne peut donc pas y varier.

Comme \(\mu_P=\nu_P\) est purement atomique, toute variation absolument continue sur un intervalle contenant un point de \(\Sigma_P\) est impossible dans le cadre fini.
\end{proof}

\begin{corollary}[Différence testable avec une description lisse]
Tout modèle concurrent qui prédit une courbe de résolution absolument continue sur un intervalle contenant un point de \(\Sigma_P\), ou qui refuse la loi de saut ci-dessus, est expérimentalement distinguable de la version finie de la théorie dès que l'on peut balayer l'intensité et compter les paires séparées.
\end{corollary}

\begin{proof}
Dans la théorie finie, \(\widetilde{\mathcal R}_P\) est une fonction en escalier et sa mesure associée est purement atomique. Un modèle absolument continu ne peut donc pas reproduire cette structure.
\end{proof}

\begin{proposition}[Calcul expérimental minimal]
Considérons
\[
\mathcal X=\{a,b,c\},\qquad \Lambda=\{\lambda_1\prec\lambda_2\prec\lambda_3\},
\]
avec
\[
\iota(\lambda_1)<\iota(\lambda_2)<\iota(\lambda_3),
\]
et
\[
\mathcal F_{\lambda_1}=\{f_1\},\qquad
\mathcal F_{\lambda_2}=\{f_1,f_2\},\qquad
\mathcal F_{\lambda_3}=\{f_1,f_2,f_3\},
\]
où les réponses sont données par
\[
\begin{array}{c|ccc}
& a & b & c \\ \hline
\Phi(P,\cdot,f_1) & 0 & 0 & 1 \\
\Phi(P,\cdot,f_2) & 0 & 1 & 1 \\
\Phi(P,\cdot,f_3) & 0 & 1 & 1
\end{array}
\]
Un choix compatible de fonction de seuil complète est
\[
\Theta_P(\{a,c\})=\Theta_P(\{b,c\})=\alpha_1,
\qquad
\Theta_P(\{a,b\})=\alpha_2,
\]
avec
\[
\alpha_1<\iota(\lambda_1)\le \alpha_2<\iota(\lambda_2).
\]
On obtient alors
\[
\nu_P=\frac23\,\delta_{\alpha_1}+\frac13\,\delta_{\alpha_2},
\]
et
\[
\widetilde{\mathcal R}_P(t)=
\begin{cases}
0,& t\le \alpha_1,\\[0.2em]
\frac23,& \alpha_1<t\le \alpha_2,\\[0.2em]
1,& t>\alpha_2.
\end{cases}
\]
En particulier,
\[
\mathcal R_P(\lambda_1)=\frac23,\qquad
\mathcal R_P(\lambda_2)=1,\qquad
\mathcal R_P(\lambda_3)=1.
\]
De plus,
\[
C(\lambda_1)=1,\quad C(\lambda_2)=2,\quad C(\lambda_3)=3,
\]
tandis que
\[
C_{\mathrm{eff},P}(\lambda_1)=1,\quad
C_{\mathrm{eff},P}(\lambda_2)=2,\quad
C_{\mathrm{eff},P}(\lambda_3)=2.
\]
Ce calcul montre simultanément un saut réel de résolution entre \(\lambda_1\) et \(\lambda_2\), puis un plateau parfait entre \(\lambda_2\) et \(\lambda_3\), malgré l'augmentation de l'information brute.

\paragraph{Lecture structurelle.}
Dans cet exemple, la hiérarchie de partitions est
\[
\{a,b,c\}
\succ
\{a,b\}\,|\,\{c\}
\succ
\{a\}\,|\,\{b\}\,|\,\{c\},
\]
et la fonction \(\Theta_P\) enregistre exactement le niveau auquel chaque paire se sépare. Le schéma est celui d'une filtration de partitions à deux sauts, dont la mesure empirique des seuils possède deux atomes.
\end{proposition}

\begin{proof}
Le calcul ci-dessus est immédiat à partir des définitions de \(\nu_P\) et \(\widetilde{\mathcal R}_P\).
\end{proof}

\section{Conclusion structurale}

La première partie établit qu'une théorie construite à partir de réponses observationnelles se factorise par l'empreinte, ses fibres, son quotient et ses représentations minimales. La seconde partie montre qu'à échelle finie la hiérarchie complète des partitions observables est encodée par une fonction de seuil qui satisfait une loi de minimum double, dont la structure arborescente est unique. La courbe de résolution n'est alors rien d'autre que la fonction de répartition d'une mesure atomique portée par le spectre critique. Ainsi, la structure observationnelle introduite dans cet article est gouvernée par deux objets : l'empreinte réduite en Partie I, et la fonction de seuil complète en Partie II.

\section{Ce qui est nouveau}

\begin{enumerate}[label=\arabic*.]
\item L'objet primaire n'est pas une distance ni une partition, mais une procédure d'observation. La théorie commence au niveau des réponses et non au niveau d'une géométrie déjà supposée.
\item La réduction observationnelle élimine les sondes redondantes sans perdre les fibres. Elle distingue l'information brute de l'information effectivement nouvelle.
\item La fonction de seuil joue le rôle d'objet de synthèse : elle reconstruit toutes les partitions observables, l'arbre des classes et la mesure de résolution.
\item La loi du minimum double ne sert pas seulement à imiter les ultramétriques ; elle fournit un critère exact pour reconnaître quand une hiérarchie de partitions provient d'une observation filtrée.
\end{enumerate}

\section{Ouverture physico-mathématique}

La théorie précédente demeure entièrement mathématique. Elle suggère cependant un programme supplémentaire : étudier si des structures physiques peuvent être reconstruites à partir de l'espace latent des empreintes et de ses invariants. La prochaine étape serait de formuler, dans un cadre séparé, un théorème physico-mathématique de reconstruction où l'espace, la dynamique, les symétries et certaines lois effectives émergent comme invariants d'un système d'observation itérée.

Une telle étape n'est pas démontrée ici. Elle constitue le prolongement naturel du cadre développé dans cet article, et elle doit être abordée avec des hypothèses physiques explicites, séparées des résultats purement mathématiques.

\section*{Remerciements}
La mise en forme, certaines reformulations et la relecture linguistique ont bénéficié d'une assistance automatisée de rédaction. L'auteur assume l'ensemble des énoncés, des démonstrations et des choix scientifiques.







\section{Résultats techniques complémentaires}

\subsection{Monotonie, redondance et robustesse}

Soient \(\mathcal G_1,\mathcal G_2\subseteq\mathcal F\) deux familles de sondes.

\begin{proposition}[Monotonie par raffinement]
Si \(\mathcal G_1\subseteq\mathcal G_2\), alors, pour tout \(P\in\mathcal P\),
\[
x\sim_{\Phi,P,\mathcal G_2}y \Longrightarrow x\sim_{\Phi,P,\mathcal G_1}y.
\]
En particulier, l'ajout de sondes ne peut qu'affiner les classes d'équivalence.
\end{proposition}

\begin{proof}
L'égalité des empreintes sur \(\mathcal G_2\) implique l'égalité sur la sous-famille \(\mathcal G_1\).
\end{proof}

\begin{corollary}
Si une famille \(\mathcal G_1\) est séparante en \(P\), alors toute famille \(\mathcal G_2\supseteq\mathcal G_1\) est également séparante en \(P\).
\end{corollary}

\begin{proof}
Immédiat par la proposition précédente.
\end{proof}

\begin{proposition}[Sondes redondantes]
Soit \(g\in\mathcal F\). Si, pour tout \(x\in\mathcal X\), la valeur \(\Phi(P,x,g)\) est déterminée par les valeurs \(\{\Phi(P,x,f):f\in\mathcal G\}\) d'une famille \(\mathcal G\subseteq\mathcal F\), alors l'ajout de \(g\) à \(\mathcal G\) ne raffine pas la relation d'indiscernabilité induite par \(\mathcal G\).
\end{proposition}

\begin{proof}
Deux objets égaux sur \(\mathcal G\) ont la même valeur sur \(g\), donc les relations d'indiscernabilité induites par \(\mathcal G\) et \(\mathcal G\cup\{g\}\) coïncident.
\end{proof}

\begin{proposition}[Robustesse uniforme]
Si deux fonctionnels \(\Phi,\Psi:\mathcal P\times\mathcal X\times\mathcal F\to\mathbb R\) satisfont
\[
\sup_{(P,x,f)}|\Phi(P,x,f)-\Psi(P,x,f)|\le \varepsilon,
\]
alors, pour toute famille finie \(\mathcal F_K\) et tout \(P\in\mathcal P\),
\[
\|I_\Phi^P(x)-I_\Psi^P(x)\|_\infty\le \varepsilon
\qquad
\text{pour tout }x\in\mathcal X.
\]
\end{proposition}

\begin{proof}
Par définition,
\[
\|I_\Phi^P(x)-I_\Psi^P(x)\|_\infty
=
\max_{1\le k\le K} |\Phi(P,x,f_k)-\Psi(P,x,f_k)|
\le \varepsilon.
\]
\end{proof}

\subsection{Tours d'empreintes}

\begin{definition}[Tour d'empreintes]
On appelle \emph{tour d'empreintes} une suite d'espaces
\[
E_0=\mathcal X,\qquad E_{n+1}=I_n(E_n)\subseteq \mathbb R^{K_n}
\]
où, pour chaque \(n\ge 0\), \(I_n:E_n\to \mathbb R^{K_n}\) est une empreinte informationnelle construite sur le niveau \(E_n\).
\end{definition}

\begin{proposition}[Stabilité de la hiérarchie]
Supposons que, pour chaque \(n\ge 0\), \(I_n\) soit injective. Alors chacun des niveaux \(E_n\) est identifiable à \(E_0\) comme ensemble. Si, de plus, chaque \(I_n\) est isométrique pour une distance induite, alors tous les niveaux sont isométriques entre eux.
\end{proposition}

\begin{proof}
L'injectivité identifie \(E_n\) à son image. L'itération donne l'identification de chaque niveau à \(E_0\). Si les \(I_n\) sont isométriques, les identifications préservent en outre les distances.
\end{proof}

\begin{remark}
La hiérarchie des empreintes formalise l'idée d'empreinte de l'empreinte. Elle ne crée pas automatiquement de nouvelle information brute si chaque étage ne dépend que du précédent, mais elle peut faire émerger de nouvelles géométries, de nouvelles symétries et de nouveaux invariants. Elle exprime aussi la possibilité d'observer les observables, puis les observables des observables, et ainsi de suite.
\end{remark}

%












%
%%%%%%%%%%%%%%%%%%%%%%%%%%%%%%%%%%%%%%%%%%%%%%%%%%%%%%%%%%%%%%%%%%%%%%%%%%%%%%%%%%%%%%%%%%%%%%%%%%%%%%%%%%%%%%%%%%%%%%%%%%%%%%%%%%%%%%%%%%%%%%%%%%%%%%%%%%%%%%%%%%%%%%%%%%%%%%%%%%%%%%%%%%%%%%%%%Principe universel de descente et de commutation des empreintes


% Théorème général unificateur de la théorie des empreintes
% issu de l'analyse des exemples Naszimger / double résolution / transport observationnel, mais formulé ici comme principe universel

\section{Principe universel de descente et de commutation des empreintes}

\begin{definition}[Structure observationnelle admissible]
Fixons un contexte \(P\in\mathcal P\), une échelle \(\lambda\in\Lambda\) et une famille finie de sondes \(\mathcal G\subseteq\mathcal F_\lambda\).
On pose
\[
I:=I_{\Phi,\mathcal G}^P,\qquad
\mathcal L:=\mathcal L_{\mathcal G,P}=I(\mathcal X),\qquad
\pi:=\pi_{\mathcal G}:\mathcal X\to\mathcal X/{\sim_{\Phi,P,\mathcal G}}.
\]

On appelle \emph{structure observationnelle admissible} toute donnée \(\mathsf S\) obtenue à partir de \(I\) par itération finie des constructions canoniques de cette théorie, telles que :
\begin{enumerate}[label=\textbf{(\arabic*)}]
\item le passage à l'image \(I(\mathcal X)\),
\item le passage au quotient par les fibres de \(I\),
\item la restriction à une sous-famille de sondes,
\item la construction, à partir de données constantes sur les fibres de \(I\), d'une relation, d'une fonction, d'un arbre, d'une donnée de seuil, d'une mesure ou d'une probabilité.
\end{enumerate}

On dit qu'une telle structure \(\mathsf S\) est \emph{fibre-compatible} si toute représentation canonique associée à \(\mathsf S\) est constante sur les fibres de \(I\), c'est-à-dire si, pour la représentation considérée,
\[
I(x)=I(y)\Longrightarrow \mathsf S(x)=\mathsf S(y).
\]
\end{definition}

\begin{remark}
La définition précédente n'introduit aucun objet primitif nouveau : elle regroupe, sous une seule étiquette, des constructions déjà démontrées ailleurs dans la théorie des empreintes.
\end{remark}

\begin{theorem}[Principe de descente et de commutation des empreintes — portée auditée]
Sous les hypothèses de la structure observationnelle admissible et si \(\mathsf S\) est fibre-compatible :
\begin{enumerate}[label=\textbf{(\alph*)}]
\item toute donnée fonctionnelle constante sur les fibres de \(I\) se factorise de manière unique à travers l'image \(\mathcal L\) et le quotient ;
\item la fermeture observationnelle conserve les fibres lorsqu'elle est définie selon le théorème de réduction des sondes ;
\item pour une structure de seuil, \(\Theta_P\) détermine le squelette hors des points critiques et \(H_P=(\Theta_P,b_P)\) détermine la filtration exacte ;
\item \(\widetilde{\mathcal R}_P\) et \(\mu_P\) déterminent la distribution des seuils avec multiplicité, mais pas leur affectation aux paires étiquetées ;
\item une invariance exacte sous reparamétrisation exige un isomorphisme strictement croissant de l'échelle qui conserve l'ordre des niveaux et le statut de bord. Une application monotone non injective peut fusionner des niveaux.
\end{enumerate}
Aucune universalité catégorique supplémentaire n'est déduite sans définition explicite des objets, morphismes et conditions d'existence concernés.
\end{theorem}
\begin{proof}
Les points (a) et (b) suivent de la factorisation par les fibres et de la fermeture observationnelle. Le point (c) suit de la lecture exacte par \(H_P\). Le point (d) suit de la construction de la mesure comme somme de masses indexées par les paires, dont le regroupement oublie l'affectation. Le point (e) suit de la conservation de l'ordre strict et des niveaux critiques.\qed\end{proof}

\begin{corollary}[Naszimger comme cas de quotient observationnel]
La reconstruction du chat de Naszimger n'est rien d'autre qu'un cas particulier du principe précédent : la classe non résolue, sa projection quotient et son raffinement par niveaux de sondes correspondent à la descente d'une structure observationnelle admissible à travers \(I\) et \(\pi\).
\end{corollary}

\begin{proof}
Dans cet exemple, la donnée fondamentale est une classe observationnelle non réduite, puis son raffinement par les sondes à plus haute résolution. C'est exactement le cas \textbf{(a)} du théorème, appliqué à une structure discrète finie.
\end{proof}



\begin{corollary}[Transport observationnel comme cas de factorisation]
Toute opération \(\mathcal G\)-compatible est un cas particulier du cas \textbf{(a)} du théorème : l'application physique ou logique \(F\) descend de manière unique à l'application induite \(\mathcal C_F\) sur l'image observationnelle réduite.
\end{corollary}

\begin{proof}
C'est exactement le théorème de factorisation de l'empreinte appliqué au morphisme \(I\circ F\). La commutation
\[
I\circ F=\mathcal C_F\circ I
\]
est la forme la plus simple du principe universel de descente.
\end{proof}

\begin{remark}
Dans ce cadre, les trois exemples ne sont donc pas trois théories séparées. Ils sont trois incarnations du même schéma :
\[
\mathcal X
\;\xrightarrow{\,I\,}\;
\mathcal L
\cong
\mathcal X/{\sim_{\Phi,P,\mathcal G}}
\Longrightarrow
\begin{cases}
\text{réduction observationnelle des sondes},\\
\text{reconstruction complète des partitions},\\
\text{complétion seuil canonique},\\
\text{arbre étiqueté},\\
\text{mesure atomique de la résolution},\\
\text{transport observationnel}.
\end{cases}
\]
Le théorème précédent formule cette signature comme un principe général.
\end{remark}


%







%
%%%%%%%%%%%%%%%%%%%%%%%%%%%%%%%%%%%%%%%%%%%%%%%%%%%%%%%%%%%%%%%%%%%%%%%%%%%%%%%%%%%%%%%%%%%%%%%%%%%%%%%%%%%%%%%%%%%%%%%%%%%%%%%%%%%%%%%%%%%%%%%%%%%%%%%%%%%%%%%%%%%%%%%%%%%%%%%%%%%%%%%%%%%%%%%%%generalisation


\section{La représentation comme géométrie de résolution observationnelle}

\subsection{Idée directrice}

Dans le langage ordinaire de la représentation, on parle d'\emph{objet d'entrée}, d'\emph{observable}, de \emph{généralisation}, de \emph{sous-résolution} et de \emph{sur-résolution}.  
Dans le langage de la théorie des empreintes, ces mots se relisent de façon plus rigoureuse :

\begin{itemize}[leftmargin=1.4em]
\item l'\emph{objet d'entrée} est un élément \(x\in\mathcal X\) ;
\item l'\emph{observable} est une valeur d'empreinte \(I_{\Phi,\mathcal G}^P(x)\) ;
\item un système de représentation est un choix de sondes et d'échelle qui induit une filtration de fibres ;
\item une représentation est ici \emph{un niveau de cette filtration}, donc une position dans une géométrie de résolution ;
\item la généralisation correspond à un \emph{régime observationnel intermédiaire} de la géométrie des seuils, c'est-à-dire à un niveau où \(0<\widetilde{\mathcal R}_P(t)<1\) ;
\item la sous-résolution et la résolution terminale correspondent respectivement à une résolution insuffisante et à une résolution déjà maximale pour le système d'observation du chapitre.
\end{itemize}

Autrement dit, représenter ne signifie pas d'abord mémoriser, mais \emph{raffiner des fibres observables}.  
Le présent chapitre ne suppose aucune structure extérieure à la théorie des empreintes : il étudie une forme interne de la représentation, entièrement lisible dans les fibres, les quotients et les seuils déjà construits.

\begin{remark}[Interprétation observationnelle]
Dans ce cadre, une représentation n'est ni un vecteur, ni un simple embedding, ni une fonction isolée : c'est une position dans une filtration observationnelle canonique, donc un niveau de résolution.
\end{remark}

\begin{remark}[Cas particulier vectoriel]
Si une représentation est donnée par un vecteur de réponses
\[
E=(E_1,\ldots,E_d):\mathcal X\to\mathbb R^d,
\]
alors chaque coordonnée \(E_i\) joue le rôle d'une sonde. La représentation vectorielle \(E\) se lit donc comme une empreinte à famille finie. Enrichir le nombre de coordonnées ou la famille de sondes revient à raffiner la filtration des fibres. Dans ce sens, une représentation vectorielle est représentable dans la théorie des empreintes.
\end{remark}

\begin{remark}[Terminologie observationnelle]
Nous utilisons ici les termes \emph{généralisation}, \emph{sous-résolution} et \emph{sur-résolution} dans un sens strictement observationnel. Ils désignent des positions relatives dans la géométrie des seuils, et non des notions importées d'une théorie extérieure de l'apprentissage.
\end{remark}

\begin{remark}
Le chapitre ne remplace pas la géométrie observationnelle déjà construite dans la théorie. Il en déploie une conséquence interne : si un système de représentation est modélisé par une suite de familles de sondes croissantes, alors sa dynamique se lit dans la filtration des partitions observables et dans l'arbre canonique associé.
\end{remark}

\subsection{Raffinement observationnel d'une trajectoire}

\begin{definition}[Trajectoire observationnelle de représentation]
Fixons un contexte \(P\in\mathcal P\), un ensemble fini d'objets d'entrée \(\mathcal D\subseteq\mathcal X\), et une famille croissante de sondes
\[
\mathcal G_0\subseteq \mathcal G_1\subseteq \cdots \subseteq \mathcal G_n \subseteq \mathcal F_\lambda
\]
correspondant à une suite finie de niveaux de représentation.

Pour chaque \(k\in\{0,\ldots,n\}\), on pose
\[
I_k:=I_{\Phi,\mathcal G_k}^P,
\qquad
\mathcal L_k:=I_k(\mathcal X),
\qquad
\Pi_k:=\mathcal D/{\sim_{\Phi,P,\mathcal G_k}},
\]
où \(\Pi_k\) désigne la partition induite sur l'ensemble d'étude \(\mathcal D\).

On appelle \emph{trajectoire observationnelle de représentation} la chaîne de restrictions
\[
\mathcal L_n
\xrightarrow{\,r_{n\to n-1}\,}
\mathcal L_{n-1}
\xrightarrow{\,r_{n\to n-2}\,}
\cdots
\xrightarrow{\,r_{1\to 0}\,}
\mathcal L_0,
\]
où \(r_{\ell\to k}\) est la restriction aux coordonnées de \(\mathcal G_k\), de sorte que
\[
I_k=r_{k+1\to k}\circ I_{k+1}
\qquad (0\le k<n).
\]
Cette chaîne décrit un raffinement progressif de l'observation.

On dit qu'une telle trajectoire est \emph{fibre-compatible} si elle est obtenue par simple ajout de sondes, donc si chaque étape est bien une projection de la suivante par oubli de coordonnées.
\end{definition}

\begin{definition}[Compatibilité avec la fonction de seuil complète]
On dit que la trajectoire \((\mathcal G_k)_{0\le k\le n}\) est \emph{compatible avec} la fonction de seuil complète \(\Theta_P\) s'il existe une suite croissante d'intensités
\[
\iota(\lambda_0)\preceq \iota(\lambda_1)\preceq \cdots \preceq \iota(\lambda_n)
\]
telle que, pour tout \(k\),
\[
\Pi_k=\mathcal D/{\sim_{\Phi,P,\lambda_k}}.
\]
Autrement dit, la trajectoire considérée est une sous-filtration de la filtration canonique déterminée par \(\Theta_P\).
\end{definition}

\begin{proposition}[Filtration canonique induite]
Si \(\mathcal G_k\subseteq \mathcal G_\ell\), alors
\[
x\sim_{\Phi,P,\mathcal G_\ell}y \Longrightarrow x\sim_{\Phi,P,\mathcal G_k}y
\qquad \text{pour tous }x,y\in\mathcal X.
\]
En particulier, la famille \(\bigl(\Pi_k\bigr)_{0\le k\le n}\) est une filtration de partitions sur \(\mathcal D\) : chaque \(\Pi_{k+1}\) raffine \(\Pi_k\).
\end{proposition}

\begin{proof}
L'égalité des empreintes sur \(\mathcal G_\ell\) implique, par restriction, l'égalité sur la sous-famille \(\mathcal G_k\). C'est exactement la monotonie du raffinement observationnel. La restriction à \(\mathcal D\) donne alors la filtration des partitions.
\end{proof}

\subsection{Application du Principe universel de descente et de commutation des empreintes}

\begin{remark}[Résultat central]
Le théorème suivant porte le cœur interprétatif du chapitre.
\end{remark}

\begin{theorem}[Descente canonique des trajectoires de représentation]
\label{thm:descente-trajectoire}
Sous les hypothèses précédentes, la famille
\[
\mathsf S^{\mathrm{traj}}
:=
\bigl(I_k,\mathcal L_k,\Pi_k,r_{k+1\to k}\bigr)_{0\le k<n}
\]
induit une structure observationnelle admissible au sens du \emph{Principe universel de descente et de commutation des empreintes}. En particulier, elle admet une descente canonique et un passage canonique au quotient, uniques, et tous les diagrammes élémentaires associés commutent.

Dans le cadre fini de la Partie II, si cette trajectoire est compatible avec la fonction de seuil complète \(\Theta_P\), alors la filtration \((\Pi_k)\) est une sous-filtration de la filtration canonique déterminée par \(\Theta_P\). Le point \textbf{(d)} du \emph{Principe universel de descente et de commutation des empreintes} fournit alors la chaîne canonique
\[
(\Pi_k,\Theta_P,\mathcal T_P)\text{ : hiérarchie indexée},\qquad \widetilde{\mathcal R}_P\leftrightarrow\mu_P\text{ : résumé scalaire}.
\]
\end{theorem}

\begin{proof}
La famille \(\mathsf S^{\mathrm{traj}}\) est construite par itération finie de restrictions, d'images et de quotients à partir des empreintes \(I_k\). Elle relève donc, par définition, des constructions canoniques admissibles du Principe universel.

Par le point \textbf{(a)} du Principe universel, toute représentation canonique associée à une donnée fibre-compatible se factorise de manière unique à travers l'image réduite et le quotient observationnel. Cela donne l'existence et l'unicité des structures descendues et quotientées, ainsi que la commutation des carrés élémentaires.

Dans le cadre fini de la Partie II, la compatibilité avec \(\Theta_P\) signifie que la filtration \((\Pi_k)\) est la restriction d'une filtration canonique. Le point \textbf{(d)} du Principe universel fournit alors la correspondance
\[
(\Pi_k,\Theta_P,\mathcal T_P)\text{ : hiérarchie indexée},\qquad \widetilde{\mathcal R}_P\leftrightarrow\mu_P\text{ : résumé scalaire},
\]
sans redémonstration supplémentaire.
\end{proof}

\begin{proposition}[Caractère canonique de la réalisation quotient]
Toute trajectoire fibre-compatible possède une réalisation canonique par sa filtration de quotients. Cette canonicité porte sur la réalisation quotient et sur les données explicitement descendues des fibres ; elle ne s'étend pas automatiquement à une représentation arbitraire qui partagerait les mêmes partitions.
\end{proposition}
\begin{proof}
À chaque niveau, le quotient par les fibres est déterminé par la relation d'indiscernabilité. Les applications de passage sont imposées par les inclusions de fibres. Une autre réalisation quotient ayant les mêmes fibres est reliée à la première par la bijection induite sur les classes. Une représentation enrichie de données supplémentaires n'est pas couverte par cette unicité.
\qed\end{proof}

\begin{remark}[Indiscernabilité des trajectoires ayant la même filtration canonique]
Deux trajectoires qui induisent la même filtration canonique sont indiscernables du point de vue de toutes les structures observationnelles admissibles. Dans ce sens précis, les sondes sont secondaires par rapport à la géométrie des fibres ; cette seconde phrase est une lecture interprétative du résultat précédent, et non son contenu démontré.
\end{remark}

\begin{remark}[Figure source]
Le manuscrit source contient ici une figure illustrative ; sa présence est documentée dans le corpus, mais l'image externe n'est pas embarquée dans cette reconstruction.
\end{remark}

\subsection{Seuils canoniques de représentation}

\begin{remark}[Courbe de résolution canonique]
On rappelle la courbe de résolution canonique déjà associée à la mesure atomique de résolution \(\mu_P=\nu_P\) :
\[
\widetilde{\mathcal R}_P(t)
:=
\nu_P\bigl((-\infty,t)\bigr)
=
\frac{1}{\binom{|\mathcal D|}{2}}
\sum_{\{x,y\}\subseteq\mathcal D}
\mathbf 1_{\{t>\Theta_P(\{x,y\})\}}.
\]
\end{remark}

\begin{proposition}[Gain de résolution]
Pour deux niveaux \(k<\ell\), on pose
\[
G_P(k,\ell):=
\widetilde{\mathcal R}_P\bigl(\iota(\lambda_\ell)\bigr)
-
\widetilde{\mathcal R}_P\bigl(\iota(\lambda_k)\bigr).
\]
Alors
\[
G_P(k,\ell)=\mu_P\bigl([\iota(\lambda_k),\iota(\lambda_\ell))\bigr),
\]
et ce gain est nul si et seulement si aucun saut atomique de la résolution ne se produit entre les deux niveaux.
\end{proposition}

\begin{proof}
La définition de \(\widetilde{\mathcal R}_P\) donne directement
\[
\widetilde{\mathcal R}_P\bigl(\iota(\lambda_\ell)\bigr)-\widetilde{\mathcal R}_P\bigl(\iota(\lambda_k)\bigr)
=
\nu_P\bigl((\iota(\lambda_k),\iota(\lambda_\ell)]\bigr)
=
\mu_P\bigl([\iota(\lambda_k),\iota(\lambda_\ell))\bigr).
\]
Comme \(\mu_P\) est atomique, l'égalité à zéro équivaut à l'absence d'atome dans l'intervalle considéré.
\end{proof}

\begin{remark}[Seuils critiques canoniques]
On note les seuils critiques déjà lus sur la courbe de résolution :
\[
\alpha_-(P):=\inf\{t\in\mathbb R_+:\ \widetilde{\mathcal R}_P(t)>0\},
\qquad
\alpha_+(P):=\sup\{t\in\mathbb R_+:\ \widetilde{\mathcal R}_P(t)<1\}.
\]
Ces bornes sont entièrement lues sur la mesure de résolution \(\mu_P\).
\end{remark}

\begin{proposition}[Intervalle de résolution intermédiaire]
La courbe \(\widetilde{\mathcal R}_P\) est nulle en dessous du premier saut, puis strictement intermédiaire entre ses premier et dernier sauts, puis égale à \(1\) au-delà du dernier saut. Plus précisément,
\[
t<\alpha_-(P)\ \Longrightarrow\ \widetilde{\mathcal R}_P(t)=0,
\]
\[
\alpha_-(P)<t<\alpha_+(P)\ \Longrightarrow\ 0<\widetilde{\mathcal R}_P(t)<1,
\]
\[
t>\alpha_+(P)\ \Longrightarrow\ \widetilde{\mathcal R}_P(t)=1.
\]
En particulier,
\[
\{t\in\mathbb R_+:0<\widetilde{\mathcal R}_P(t)<1\}=(\alpha_-(P),\alpha_+(P)].
\]
\end{proposition}

\begin{proof}
Comme \(\mu_P=\nu_P\) est purement atomique et que \(\widetilde{\mathcal R}_P\) est sa fonction de répartition, la courbe est une fonction en escalier dont les sauts sont exactement les atomes de \(\mu_P\). Les deux bornes \(\alpha_-(P)\) et \(\alpha_+(P)\) sont donc respectivement le premier et le dernier seuil critique lus sur cette courbe. Il s'ensuit que \(\widetilde{\mathcal R}_P\) est nulle avant le premier saut, strictement entre \(0\) et \(1\) entre le premier et le dernier saut, puis vaut \(1\) après le dernier saut.
\end{proof}

\begin{remark}[Portée de la résolution terminale]
La résolution terminale est ici une notion strictement observationnelle : elle signifie uniquement que toutes les distinctions canoniques ont déjà été réalisées. Dans un contexte d'apprentissage, ce régime pourra être interprété comme une forme de sur-résolution observationnelle.
\end{remark}

\subsection{Tripartition observationnelle des régimes de représentation}

\begin{remark}[Sens de la généralisation observationnelle]
La généralisation ne renvoie ici ni à une erreur de test ni à une performance statistique ; elle désigne uniquement un régime intermédiaire de résolution, caractérisé par \(0<\widetilde{\mathcal R}_P(t)<1\).
\end{remark}

\begin{remark}[Résultat central]
Le résultat suivant constitue l'interprétation observationnelle fondamentale de la notion de représentation.
\end{remark}

\begin{theorem}[Tripartition observationnelle des régimes de représentation]
\label{thm:tripartition-representation}
Sous les hypothèses précédentes, supposons que la trajectoire \((\mathcal G_k)_{0\le k\le n}\) soit fibre-compatible et compatible avec la fonction de seuil complète \(\Theta_P\). Alors les trois régimes suivants sont déterminés par la position de \(\iota(\lambda_n)\) par rapport aux seuils canoniques \(\alpha_-(P)\) et \(\alpha_+(P)\) :

\begin{enumerate}[label=\textbf{(\arabic*)}]
\item \textbf{Résolution insuffisante.} Si \(\iota(\lambda_n)< \alpha_-(P)\), alors l'observation reste insuffisante : la partition induite sur \(\mathcal D\) est la partition triviale, et aucune séparation non triviale n'est encore réalisée.

\item \textbf{Généralisation observationnelle.} Si
\[
\alpha_-(P)< \iota(\lambda_n)< \alpha_+(P),
\]
alors la représentation correspond à une coupe non terminale de l'arbre canonique \(\mathcal T_P\), donc à un nœud non terminal de la géométrie observationnelle.

\item \textbf{Résolution terminale.} Si \(\iota(\lambda_n)> \alpha_+(P)\), alors la résolution est terminale. Au point exact \(\alpha_+(P)\), la valeur dépend de la convention de seuil et des atomes ; elle ne peut être classée sans hypothèse supplémentaire.
\end{enumerate}

De plus, si l'intervalle \((\alpha_-(P),\alpha_+(P))\) est vide, la trajectoire ne possède pas de régime intermédiaire.
\end{theorem}

\begin{proof}
Par le Principe universel, plus précisément son point \textbf{(d)}, la fonction de seuil complète détermine canoniquement la filtration des partitions puis l'arbre laminaire ; les seuils critiques correspondent donc aux niveaux de coupe de cet arbre. Les niveaux \(\iota(\lambda)\) ne font que paramétrer cette géométrie interne, sans ajouter de donnée étrangère.

Si \(\iota(\lambda_n)< \alpha_-(P)\), la courbe de résolution n'a pas encore franchi le premier saut : aucune séparation non triviale n'est disponible, d'où la résolution insuffisante.

Si \(\alpha_-(P)< \iota(\lambda_n)< \alpha_+(P)\), la courbe de résolution a franchi certains sauts mais pas tous. La partition induite est alors une coupe non terminale de l'arbre canonique : elle n'est ni la partition triviale de la racine, ni la partition discrète des feuilles. C'est précisément le régime intermédiaire de représentation, appelé ici généralisation observationnelle.

Si \(\iota(\lambda_n)> \alpha_+(P)\), tous les seuils critiques sont franchis. La représentation est alors au niveau terminal. Le point exactement égal à \(\alpha_+(P)\) reste convention-dépendant.

La dernière assertion est immédiate : si l'intervalle \((\alpha_-(P),\alpha_+(P))\) est vide, il n'existe aucun niveau strictement intermédiaire entre le premier et le dernier seuil critique.
\end{proof}

\begin{remark}[Lecture des trois phases]
La résolution insuffisante, la généralisation observationnelle et la résolution terminale ne sont pas, dans ce cadre, trois mystères séparés : ce sont trois positions d'une même trajectoire dans l'arbre des seuils.
\end{remark}

\subsection{Lecture pour les systèmes de représentation}

\begin{remark}[Interprétation observationnelle de l'entrée et de la sortie]
Dans toute lecture observationnelle d'un système de représentation, l'objet d'entrée est un \(x\in\mathcal X\), tandis que la sortie produite par le système est un observable de la forme
\[
I_{\Phi,\mathcal G}^P(x).
\]
Changer le système d'observation, les sondes ou l'échelle de résolution modifie donc la partition induite sur l'espace des objets d'entrée. Dans ce sens, une représentation est une position dans une filtration observationnelle, et non un objet isolé.
\end{remark}

\begin{corollary}[Représentation arborescente de la filtration]
Une filtration finie d'empreintes observables possède une représentation canonique en arbre de classes. Cet arbre n'ajoute aucune donnée extra : il représente la même structure de filtrations sous une forme laminaire, conformément au point \textbf{(d)} du \emph{Principe universel de descente et de commutation des empreintes}. Les nœuds non terminaux représentent des classes encore non atomisées, tandis que les feuilles terminales représentent les atomes de résolution.
\end{corollary}

\begin{proof}
La fonction de seuil complète détermine la filtration canonique des partitions, l'arbre laminaire, les intervalles canoniques et la mesure atomique de résolution. L'arbre ne rajoute donc aucune donnée extra : il représente la même structure de filtrations sous une forme laminaire. Si un système de représentation construit progressivement des distinctions entre observables, alors ses niveaux se lisent comme des coupes de cet arbre. Les nœuds non terminaux correspondent aux classes qui survivent à plusieurs niveaux de résolution ; les feuilles correspondent aux atomes apparus à haute résolution.
\end{proof}

\begin{corollary}[Gain de résolution et changement de sondes]
Le gain de résolution entre deux niveaux de représentation est mesuré par l'accroissement de la courbe \(\widetilde{\mathcal R}_P\), donc par la masse atomique de résolution franchie entre les deux niveaux. Il n'est pas une propriété de l'objet seul, mais du couple formé par l'objet et le système d'observation.
\end{corollary}

\begin{proof}
Le gain de résolution entre deux niveaux est
\[
G_P(k,\ell)
=
\widetilde{\mathcal R}_P\bigl(\iota(\lambda_\ell)\bigr)
-
\widetilde{\mathcal R}_P\bigl(\iota(\lambda_k)\bigr)
=
\mu_P\bigl([\iota(\lambda_k),\iota(\lambda_\ell))\bigr).
\]
Il dépend donc des seuils franchis par la trajectoire d'observation, et non de l'objet seul.
\end{proof}

\begin{remark}
Dans cette perspective, l'entraînement d'un système de représentation ressemble à la construction progressive d'un arbre de résolution. Cette lecture ouvre naturellement la transition vers les systèmes de représentation concrets, sans introduire de structure supplémentaire à la théorie des empreintes.
\end{remark}


\begin{theorem}[Réalisation canonique d'une filtration observationnelle]
Toute filtration finie de partitions issue des empreintes possède une réalisation canonique par les quotients \(q_k:\mathcal X\to\mathcal X/{\sim_k}\) et les applications de passage induites par inclusion des fibres. Cette réalisation est unique à isomorphisme près parmi les réalisations explicitement définies comme quotients de ces mêmes fibres.

La filtration ne détermine pas, sans hypothèse supplémentaire, toute représentation arbitraire partageant les mêmes partitions : une représentation peut contenir une information de codage ou de structure supplémentaire.
\end{theorem}
\begin{proof}
L'inclusion \(\sim_{k+1}\subseteq\sim_k\) induit une application unique \(r_{k+1\to k}\) telle que \(q_k=r_{k+1\to k}\circ q_{k+1}\). L'unicité à isomorphisme près suit de la propriété universelle élémentaire des quotients. Aucune conclusion n'est tirée pour les représentations qui ajoutent des données non déterminées par les fibres.\qed\end{proof}


\subsection{Conclusion du chapitre}

La représentation, dans le langage des empreintes, n'est pas un agrégat de notions indépendantes. C'est une filtration observationnelle canonique, donc une position dans la géométrie de résolution, dont la généralisation observationnelle constitue le régime intermédiaire. La résolution insuffisante et la résolution terminale décrivent les deux bords de cette géométrie.

Cette lecture montre que les représentations utilisées dans les systèmes modernes peuvent être interprétées comme des coupes d'une géométrie observationnelle canonique, sans introduire de structure supplémentaire. La théorie des empreintes fournit ainsi un cadre unifié où les notions de représentation, de résolution et de raffinement ne sont plus primitives mais dérivées d'une même structure.


%










%
%%%%%%%%%%%%%%%%%%%%%%%%%%%%%%%%%%%%%%%%%%%%%%%%%%%%%%%%%%%%%%%%%%%%%%%%%%%%%%%%%%%%%%%%%%%%%%%%%%%%%%%%%%%%%%%%%%%%%%%%%%%%%%%%%%%%%%%


\section{L'objet observationnel universel et l'unité canonique des faces}

\subsection{Idée directrice}

Le chapitre précédent a montré que les partitions, la fonction de seuil complète, l'arbre, la courbe de résolution et la mesure atomique ne sont pas des objets indépendants.  
La présente section précise cette unité sans introduire de primitive nouvelle.

L'idée n'est pas de postuler un objet caché au-dessus de la théorie, mais de reconnaître que les constructions déjà obtenues sont les \emph{faces} canoniques d'une même classe observationnelle.  
Autrement dit, si l'on veut parler d'un \emph{objet observationnel universel}, il ne s'agit pas ici d'un nouvel élément fondamental : c'est le nom abrégé de la classe canonique des descriptions globales déjà reconstruites par la théorie des empreintes.

\begin{remark}[Formulation retenue]
Le terme \emph{objet observationnel universel} désigne ici la classe d'équivalence canonique des descriptions globales observationnelles déjà produites par la théorie. Il ne désigne ni un tuple supplémentaire, ni une primitive nouvelle, ni un objet situé au-dessus de la logique interne du manuscrit.
\end{remark}

\begin{remark}
Le but de cette section est de montrer que l'on peut lire la théorie entière comme une suite d'oublis contrôlés : à chaque niveau, on oublie une partie de l'information de description, mais ce qui est oublié est reconstruisible canoniquement dans les situations couvertes par les théorèmes de reconstruction de la théorie.
\end{remark}

\subsection{Descriptions observationnelles canoniques}

\begin{definition}[Description observationnelle canonique]
Fixons un contexte \(P\in\mathcal P\), dans les hypothèses de la Partie II, avec \(\mathcal X\) fini, \(\Lambda\) totalement ordonné et \(\iota:\Lambda\to\mathbb R_+\) injective.

On appelle \emph{description observationnelle canonique} de \(P\) toute écriture canonique déjà produite par la théorie dans le contexte \(P\).

On distingue deux familles.

\medskip
\noindent\textbf{(i) Faces globales canoniques.}
Ce sont les descriptions
\[
(\sim_{\Phi,P,\lambda})_{\lambda\in\Lambda},
\qquad
\Theta_P,
\qquad
\mathcal T_P,
\qquad
\widetilde{\mathcal R}_P,
\qquad
\mu_P.
\]

\medskip
\noindent\textbf{(ii) Descriptions partielles.}
Ce sont les images réduites
\[
I_{\Phi,\mathcal G}^P(\mathcal X),
\]
obtenues à partir d'une famille finie de sondes \(\mathcal G\subseteq\mathcal F_\lambda\).

On note
\[
\mathrm{Desc}_{\mathrm{glob}}(P)
:=
\Bigl\{
(\sim_{\Phi,P,\lambda})_{\lambda\in\Lambda},
\Theta_P,
\mathcal T_P,
\widetilde{\mathcal R}_P,
\mu_P
\Bigr\},
\]
\[
\mathrm{Desc}_{\mathrm{part}}(P)
:=
\Bigl\{
I_{\Phi,\mathcal G}^P(\mathcal X)\;:\;\mathcal G\subseteq\mathcal F_\lambda,\ \mathcal G \text{ finie}
\Bigr\},
\]
et
\[
\mathrm{Desc}(P):=\mathrm{Desc}_{\mathrm{glob}}(P)\sqcup \mathrm{Desc}_{\mathrm{part}}(P).
\]
Les deux familles étant de nature différente, on les considère comme disjointes par codage.

On définit sur \(\mathrm{Desc}_{\mathrm{glob}}(P)\) la relation \(\approx_P\) comme la plus petite relation d'équivalence contenant les identifications canoniques de reconstruction fournies par le \emph{Théorème structural des seuils} et par le \emph{Principe universel de descente et de commutation des empreintes}.

On appelle \emph{objet observationnel universel} de \(P\) la classe quotient
\[
\mathbb O_P := \mathrm{Desc}_{\mathrm{glob}}(P)/\!\approx_P.
\]
\end{definition}

\begin{remark}
La classe \(\mathbb O_P\) ne rajoute aucune primitive. Elle nomme le réseau des descriptions globales reliées par les transformations de reconstruction effectivement démontrées ; les résumés scalaires restent distincts des données indexées.
\end{remark}

\begin{remark}[Sens de face]
Une face est ici une description globale canonique appartenant à la classe \(\mathbb O_P\). Elle ne constitue pas un objet indépendant de la théorie.
\end{remark}

\begin{proposition}[Faces canoniques et descriptions partielles — portée corrigée]
La hiérarchie exacte des partitions, la donnée de seuil complétée \(H_P=(\Theta_P,b_P)\) et l'arbre étiqueté construit à partir de cette hiérarchie sont mutuellement reconstructibles sous les hypothèses de bord déclarées. En revanche, \(\widetilde{\mathcal R}_P\) et \(\mu_P\) sont des résumés scalaires : ils conservent la distribution des seuils avec multiplicité mais pas leur affectation aux paires.

Les images réduites \(I_{\Phi,\mathcal G}^P(\mathcal X)\) sont des descriptions partielles obtenues par oubli de coordonnées. Leur relation avec les faces globales est une relation de projection et de factorisation, non une égalité de contenu.
\end{proposition}
\begin{proof}
La première assertion suit de la lecture exacte par \(H_P\) et de la construction de l'arbre à partir de la hiérarchie. La seconde suit de la définition des résumés scalaires comme agrégation des valeurs de \(\Theta_P\).\qed\end{proof}

\begin{theorem}[Unité conditionnelle des faces observationnelles]
Sous les hypothèses finies de la Partie II, la hiérarchie exacte des partitions, la donnée de seuil complétée \(H_P=(\Theta_P,b_P)\) et l'arbre étiqueté construit à partir de cette hiérarchie sont mutuellement reconstructibles. En revanche, la courbe de résolution et la mesure de résolution ne sont que des résumés scalaires : elles ne permettent pas en général de retrouver l'affectation des seuils aux paires.
\end{theorem}
\begin{proof}
La hiérarchie exacte détermine le statut de bord et les seuils par borne supérieure. Réciproquement, \(H_P\) donne la relation exacte à chaque niveau. L'arbre est construit à partir des classes. La mesure et la courbe regroupent les seuils avec multiplicité et perdent l'indexation par paire.\qed\end{proof}

\begin{corollary}[Lecture par oubli de coordonnées]
Soit \(\mathcal G\subseteq\mathcal F_\lambda\) une famille finie de sondes. Alors l'image réduite
\[
I_{\Phi,\mathcal G}^P(\mathcal X)
\]
est une description partielle issue d'un représentant d'une face globale canonique de \(P\), obtenu relativement à la famille \(\mathcal G\). Le passage du total à cette image réduite consiste à oublier toutes les informations qui ne sont pas visibles à la famille \(\mathcal G\).

Les descriptions partielles se lisent donc comme des factorisations locales de la théorie, sans être identifiées aux faces globales canoniques.
\end{corollary}

\begin{proof}
C'est une reformulation du théorème de factorisation et du Principe universel. Toute donnée constante sur les fibres de l'empreinte se factorise de manière unique à travers l'image réduite. Par conséquent, l'image réduite est une description partielle issue d'un représentant d'une face globale canonique de \(P\), obtenu relativement à la famille \(\mathcal G\). La reconstruction canonique de la face globale correspondante relève alors des résultats déjà établis lorsque les hypothèses de reconstruction sont réunies.
\end{proof}

\begin{corollary}[Cas vectoriel comme description partielle finie]
Si un système de représentation fournit une application vectorielle
\[
E=(E_1,\ldots,E_d):\mathcal X\to\mathbb R^d,
\]
et si chaque coordonnée \(E_i\) est utilisée comme une sonde observationnelle, alors la famille
\[
\mathcal G=\{E_1,\ldots,E_d\}
\]
définit une empreinte finie au sens de la théorie.

Dans ce cas, la représentation vectorielle ne constitue pas un objet extérieur à la théorie : elle apparaît comme une description partielle finie issue d'un représentant d'une face globale canonique de \(P\).
\end{corollary}

\begin{proof}
Chaque coordonnée \(E_i\) joue alors exactement le rôle d'une sonde. La famille finie \(\mathcal G\) induit l'empreinte \(I_{\Phi,\mathcal G}^P\), donc une image réduite et une filtration des fibres. La représentation vectorielle n'est donc pas une primitive indépendante : elle est une description partielle finie issue d'un représentant d'une face globale canonique de \(P\), au sens canonique déjà fixé par la théorie.
\end{proof}

\subsection{Conséquences théoriques}

\begin{proposition}[Aucun surcroît d'objet]
Dans la théorie des empreintes, il n'est pas nécessaire d'ajouter un nouvel objet fondamental au-dessus des empreintes, des partitions, des arbres et des mesures. L'objet observationnel universel \(\mathbb O_P\) sert à nommer le réseau des descriptions globales effectivement reliées ; il ne transforme pas les résumés scalaires en invariants complets.
\end{proposition}

\begin{proof}
L'objet \(\mathbb O_P\) est entièrement constitué de données déjà présentes dans la théorie. Il ne crée aucune nouvelle primitive ; il condense seulement en une seule écriture le contenu déjà reconstruit par les théorèmes antérieurs. Aucun objet supplémentaire n'est requis pour décrire les structures canoniques considérées.
\end{proof}

\begin{remark}
Cette manière de présenter les choses modifie la lecture de l'ensemble du manuscrit : les empreintes, les partitions, l'arbre, la résolution et la mesure ne sont plus vues comme des objets séparés, mais comme les faces globales d'une même classe canonique. La théorie gagne ainsi en unité sans rien perdre de sa rigueur.
\end{remark}

\subsection{Conclusion de la section}

L'objet observationnel universel, au sens strict de cette théorie, n'est pas une nouvelle entité cachée.  
C'est la classe canonique déjà reconstruite par la théorie des empreintes : la classe des faces globales canoniques mutuellement équivalentes, tandis que les descriptions partielles sont des factorisations locales associées à des familles finies de sondes.

Ainsi, la théorie ne juxtapose pas des objets hétérogènes. Elle organise plusieurs descriptions canoniquement compatibles d'une même structure observationnelle :
\[
\text{empreinte}
\Longleftrightarrow
\text{partition}
\Longleftrightarrow
\text{fonction de seuil}
\Longleftrightarrow
\text{arbre}
\Longleftrightarrow
\text{mesure}
\Longleftrightarrow
\text{résolution}.
\]
La nouveauté conceptuelle est précisément là : ce n'est pas un nouvel être théorique, mais l'unité canonique de toutes ses faces déjà démontrées.


%









%
%%%%%%%%%%%%%%%%%%%%%%%%%%%%%%%%%%%%%%%%%%%%%%%%%%%%%%%%%%%%%%%%%%%%%%%%%%%%%%%%%%%%%%%%%%%%%%%%%%%%%%%%%%%%%%%%%%%%%%%%%%%%%%%%%%%%%%%%%%%%%%%%%%%%%%%%%%%%%%%%%%%%%%%%%%%%%%%%%%%%%%%%%%%%%%%%%%%%%%%%%%%%%%%%%%%%%%%%%%%%%%%%%%%%%%%%%%%%%%%%%%%%%%%%%%%%%%%%%%%%%%%%%%%%%%%%%%%%%%%%%%%%%%%%%exp 1 puzzel puzzle_experiment.py



\section{Annexe expérimentale : protocole, pipeline, résultats et interprétation}

\subsection{Objectif de l'expérience}

Cette expérience a pour objectif de vérifier, dans un cadre strictement observationnel et sans apprentissage, que la théorie des empreintes permet de reconstruire la structure cachée d'un monde fini à partir des seules réponses des sondes.

L'idée centrale est la suivante : on construit un monde caché composé d'objets munis de propriétés inaccessibles à l'observateur, puis on n'autorise l'observateur qu'à interroger ces objets par une famille croissante de sondes. À partir des seules empreintes observables, l'observateur doit reconstruire successivement
\[
\Pi_k,\qquad \Theta_P,\qquad \mathcal T_P,\qquad \mu_P,\qquad \widetilde{\mathcal R}_P,
\]
dans cet ordre, sans jamais consulter les propriétés cachées pendant la reconstruction.

L'expérience ne cherche donc pas à entraîner un modèle, mais à tester la cohérence interne d'une géométrie de résolution observationnelle.

\subsection{Principe expérimental}

On considère un ensemble fini d'objets
\[
\mathcal X=\{x_1,\dots,x_{100}\},
\]
où chaque objet possède \(10\) propriétés cachées, inaccessibles après la phase de génération.

On définit une famille croissante de sondes
\[
\mathcal G_0\subseteq \mathcal G_1\subseteq \cdots \subseteq \mathcal G_n,
\]
et, pour chaque niveau \(k\), l'empreinte observationnelle
\[
I_{\Phi,\mathcal G_k}^P:\mathcal X\to \mathbb R^{\mathcal G_k}.
\]
La partition associée est
\[
\Pi_k=\mathcal X/{\sim_{\Phi,P,\mathcal G_k}},
\qquad
x\sim_{\Phi,P,\mathcal G_k}y
\Longleftrightarrow
I_{\Phi,\mathcal G_k}^P(x)=I_{\Phi,\mathcal G_k}^P(y).
\]

À partir de ces partitions, on reconstruit la donnée de seuil complétée \(H_P=(\Theta_P,b_P)\) lorsque les niveaux critiques sont observables, puis le spectre critique
\[
\Sigma_P=\{\Theta_P(\{x,y\}) : x\neq y\},
\]
puis l'arbre canonique \(\mathcal T_P\), la mesure atomique \(\mu_P\) et enfin la courbe de résolution
\[
\widetilde{\mathcal R}_P(t)
=
\mu_P((-\infty,t))
=
\frac{1}{\binom{|\mathcal X|}{2}}
\sum_{\{x,y\}\subseteq\mathcal X}
\mathbf 1_{\{t>\Theta_P(\{x,y\})\}}.
\]

\subsection{Pipeline suivi}

Le pipeline a été imposé dans un ordre strict :

\begin{enumerate}
\item génération du monde caché ;
\item génération des sondes ;
\item calcul des empreintes \(I_{\Phi,\mathcal G_k}^P\) ;
\item calcul des partitions \(\Pi_k\) ;
\item reconstruction de \(\Theta_P\) ;
\item reconstruction du spectre \(\Sigma_P\) ;
\item reconstruction de l'arbre \(\mathcal T_P\) ;
\item reconstruction de la mesure atomique \(\mu_P\) ;
\item reconstruction de la courbe \(\widetilde{\mathcal R}_P\) ;
\item vérifications théoriques ;
\item validation finale contre le monde caché, uniquement à la fin.
\end{enumerate}

Aucun accès aux propriétés cachées n'a été autorisé pendant les étapes 3 à 10.

\subsection{Protocole de sécurité observationnelle}

Le protocole impose les contraintes suivantes :

\begin{itemize}[leftmargin=1.4em]
\item séparation stricte entre la couche \emph{monde caché} et la couche \emph{observateur} ;
\item interdiction de lire les propriétés cachées après leur génération ;
\item reconstruction uniquement à partir des réponses de sondes ;
\item absence totale d'apprentissage, d'optimisation, de clustering ou d'heuristique de regroupement ;
\item comparaison avec le monde caché uniquement après la reconstruction complète ;
\item conservation de tous les artefacts intermédiaires pour audit externe ;
\item reproductibilité par graine fixée.
\end{itemize}

Cette séparation est essentielle : sans elle, l'expérience deviendrait une simple illustration et non une reconstruction observationnelle falsifiable.

\subsection{Équations utilisées}

Les équations structurantes de l'expérience sont les suivantes.

\paragraph{Empreintes observationnelles.}
Pour chaque niveau \(k\),
\[
I_{\Phi,\mathcal G_k}^P(x)(g)=\Phi(P,x,g),
\qquad g\in\mathcal G_k.
\]

\paragraph{Partitions observables.}
\[
x\sim_{\Phi,P,\mathcal G_k}y
\Longleftrightarrow
I_{\Phi,\mathcal G_k}^P(x)=I_{\Phi,\mathcal G_k}^P(y).
\]

\paragraph{Fonction de seuil.}
Pour chaque paire \(\{x,y\}\),
\[
x\sim_{\Phi,P,\lambda}y
\Longleftrightarrow
\lambda\in S_{x,y}.
\]

\paragraph{Spectre critique.}
\[
\Sigma_P=\{\Theta_P(\{x,y\}) : x\neq y\}.
\]

\paragraph{Mesure atomique.}
\[
\nu_P
=
\frac{1}{\binom{|\mathcal X|}{2}}
\sum_{\{x,y\}\subseteq\mathcal X}\delta_{\Theta_P(\{x,y\})}.
\]

\paragraph{Courbe de résolution.}
\[
\widetilde{\mathcal R}_P(t)
=
\nu_P((-\infty,t))
=
\frac{1}{\binom{|\mathcal X|}{2}}
\sum_{\{x,y\}\subseteq\mathcal X}
\mathbf 1_{\{t>\Theta_P(\{x,y\})\}}.
\]

\subsection{Résultats obtenus}

L'expérience a confirmé les points suivants :

\begin{itemize}[leftmargin=1.4em]
\item le raffinement des partitions est monotone ;
\item la loi du minimum double est vérifiée pour tous les triplets testés ;
\item les classes observées forment une structure laminaire ;
\item l'arbre du protocole testé est reconstruit à partir de la donnée de seuil et de la convention de bord effectivement présente dans l'expérience ;
\item le plateau terminal est bien détecté ;
\item la validation finale contre le monde caché confirme l'identité entre la structure reconstruite et la structure réelle.
\end{itemize}

Le seul point signalé comme falsifié est
\[
\texttt{curve\_and\_measure : FALSIFIED}.
\]
Ce point ne remet pas en cause la théorie elle-même : il provient du script et de la façon dont le test a été écrit. En particulier, le test vérifie un \emph{jump at \(b\)} alors que la théorie impose un \emph{jump immediately after \(\alpha\)} avec \(\alpha=2.5\). Autrement dit, le problème vient de la convention de détection du saut, non de la structure atomique de la courbe.

La correction conceptuelle est simple : dans la théorie, la courbe de résolution est une fonction en escalier à saut droit au voisinage de chaque atome, donc le test doit vérifier la présence d'un atome dans l'intervalle adéquat autour du saut, et non l'égalité du saut au point de grille suivant.

\subsection{Interprétation}

L'expérience montre que la théorie peut être lue comme un mécanisme de reconstruction purement observationnel :

\[
\text{sondes}
\Longrightarrow
\text{empreintes}
\Longrightarrow
\Pi_k
\Longrightarrow
\Theta_P
\Longrightarrow
\mathcal T_P
\Longrightarrow
\mu_P
\Longrightarrow
\widetilde{\mathcal R}_P.
\]

Elle met donc en évidence que la représentation n'est pas d'abord un vecteur ou un embedding, mais une position dans une filtration de fibres. La généralisation, dans ce cadre, n'est pas une erreur statistique : c'est un régime intermédiaire de résolution. La résolution insuffisante et la résolution terminale en sont les deux bords.

\subsection{Conséquences}

Cette expérience a plusieurs conséquences théoriques.

D'abord, elle fournit une test expérimental de la reconstruction observationnelle dans un monde fini, sans apprentissage ni optimisation.

Ensuite, elle montre que la hiérarchie cachée peut être reconstruite sans accès direct aux propriétés cachées, ce qui renforce l'idée que la donnée fondamentale n'est pas la propriété elle-même, mais l'empreinte qu'elle laisse à travers les sondes.

Enfin, elle suggère que la théorie des empreintes peut servir de cadre unifié pour analyser des systèmes de représentation, des partitions observées, des seuils de séparation et des régimes de résolution dans d'autres domaines où l'observation progressive joue un rôle central.

\subsection{Conclusion expérimentale}

L'expérience ne prétend pas prouver la théorie dans toute sa généralité. Elle établit cependant quelque chose de fort : dans un monde synthétique entièrement contrôlé, les objets cachés, les sondes et les empreintes suffisent à reconstruire la géométrie de résolution annoncée par la théorie.

Le seul échec repéré, \texttt{curve\_and\_measure}, est un artefact de vérification lié à la convention de test du saut. Il doit être corrigé au niveau du script de validation, non interprété comme une réfutation de la théorie.

Ainsi, l'expérience est cohérente avec le cadre observationnel sur le protocole étudié, la fidélité du protocole et la capacité de la théorie à faire apparaître ses propres structures à partir des seules réponses observables.


%















%
%%%%%%%%%%%%%%%%%%%%%%%%%%%%%%%%%%%%%%%%%%%%%%%%%%%%%%%%%%%%%%%%%%%%%%%%%%%%%%%%%%%%%%%%%%%%%%%%%%%%%%%%%%%%%%%%%%%%%%%%%%%%%%%%%%%%%%%%%%%%%%%%%%%%%%%%%%%%%%%%%%%%%%%%%%%%%%%%%%%%%%%%%%%%%%%%
%extension : chat de naszimger :
\section{Exemple géométrique : le chat de Naszimger et la reconstruction d'une expérience de type Schrödinger}

\subsection{Principe de l'exemple}

On construit ici un exemple entièrement formulé dans le langage des empreintes observationnelles. L'objectif n'est pas de postuler une interprétation standard extérieure à la théorie, mais de montrer qu'un objet interne à l'espace des empreintes peut coder, de manière rigoureuse, une expérience à plusieurs niveaux de résolution, où une classe observationnelle est d'abord indiscernable puis se raffine lorsque les sondes deviennent plus discriminantes.

Dans ce cadre, le \emph{chat de Naszimger} ne désigne pas un état physique élémentaire, mais une \emph{classe observationnelle}. Le symbole \(\chi_{\mathrm{ferme}}\) n'est ici qu'un représentant choisi de cette classe à basse résolution, et non un état physique autonome. Cette précision évite toute confusion entre un représentant particulier et la classe géométrique qu'il représente.

Nous ne prétendons pas reconstruire le formalisme quantique de Schrödinger ; nous montrons uniquement que la structure logique de séparation observationnelle admet un représentant canonique dans la théorie des empreintes.

\subsection{Espace géométrique des empreintes}

Fixons un contexte \(P\in\mathcal P\) et une famille finie de sondes
\[
\mathcal F=
\{f_{\mathrm{boite}},f_{\mathrm{detecteur}},f_{\mathrm{vivant}},f_{\mathrm{mort}},f_{\mathrm{memoire}}\}.
\]
On considère un ensemble fini d'objets observationnels
\[
\mathcal X=
\{\chi_0,\chi_{\mathrm{ouvert}},\chi_{\mathrm{ferme}},\chi_{\mathrm{vivant}},\chi_{\mathrm{mort}}\}.
\]

On définit le fonctionnel d'observation \(\Phi:\mathcal P\times\mathcal X\times\mathcal F\to\mathbb R\) par le tableau de réponses suivant :
\[
\begin{array}{c|ccccc}
& f_{\mathrm{boite}} & f_{\mathrm{detecteur}} & f_{\mathrm{vivant}} & f_{\mathrm{mort}} & f_{\mathrm{memoire}} \\ \hline
\chi_0 & 0 & 0 & 0 & 0 & 0 \\
\chi_{\mathrm{ouvert}} & 1 & 1 & 0 & 0 & 0 \\
\chi_{\mathrm{ferme}} & 1 & 0 & 0 & 0 & 0 \\
\chi_{\mathrm{vivant}} & 1 & 1 & 1 & 0 & 1 \\
\chi_{\mathrm{mort}} & 1 & 1 & 0 & 1 & 1
\end{array}
\]
et l'on note
\[
I_\Phi^P(x)=\bigl(\Phi(P,x,f_{\mathrm{boite}}),\Phi(P,x,f_{\mathrm{detecteur}}),\Phi(P,x,f_{\mathrm{vivant}}),\Phi(P,x,f_{\mathrm{mort}}),\Phi(P,x,f_{\mathrm{memoire}})\bigr).
\]

Le niveau géométrique pertinent est le quotient observationnel. On introduit donc, pour tout \(\lambda\in\Lambda\),
\[
\pi_\lambda:\mathcal X\longrightarrow \mathcal X/{\sim_{\Phi,P,\lambda}},
\qquad
\pi_\lambda(x)=[x]_\lambda.
\]
La classe du chat de Naszimger est alors l'élément
\[
[\chi_{\mathrm{ferme}}]_{\lambda_0}\in \mathcal X/{\sim_{\Phi,P,\lambda_0}},
\]
dont \(\chi_{\mathrm{ferme}}\) n'est qu'un représentant parmi d'autres à basse résolution.

À basse résolution, plusieurs objets peuvent donner la même observation. En particulier,
\[
\chi_{\mathrm{ouvert}} \sim_{\Phi,P,\lambda_0} \chi_{\mathrm{ferme}}
\sim_{\Phi,P,\lambda_0} \chi_{\mathrm{vivant}}
\sim_{\Phi,P,\lambda_0} \chi_{\mathrm{mort}},
\qquad
\chi_0 \not\sim_{\Phi,P,\lambda_0} \chi_{\mathrm{ouvert}}.
\]
À l'échelle \(\lambda_1\), la sonde \(f_{\mathrm{detecteur}}\) sépare déjà \(\chi_{\mathrm{ferme}}\) des trois autres objets non triviaux ; à l'échelle \(\lambda_2\), les états \(\chi_{\mathrm{vivant}}\) et \(\chi_{\mathrm{mort}}\) deviennent en outre séparables grâce aux coordonnées finales.

On obtient ainsi une hiérarchie de trois régimes : un régime de fermeture où les objets non triviaux sont confondus, un régime intermédiaire où l'état ferme se sépare, puis un régime final où les issues vivante et morte sont distinguées.

\subsection{Hiérarchie des seuils}

On fixe une échelle
\[
\Lambda=\{\lambda_0\prec\lambda_1\prec\lambda_2\},
\qquad
\iota(\lambda_0)=0.5,\ \iota(\lambda_1)=1.5,\ \iota(\lambda_2)=2.5,
\]
et l'on pose
\[
\mathcal F_{\lambda_0}=\{f_{\mathrm{boite}}\},\qquad
\mathcal F_{\lambda_1}=\{f_{\mathrm{boite}},f_{\mathrm{detecteur}}\},\qquad
\mathcal F_{\lambda_2}=\{f_{\mathrm{boite}},f_{\mathrm{detecteur}},f_{\mathrm{vivant}},f_{\mathrm{mort}},f_{\mathrm{memoire}}\}.
\]

Les classes d'indiscernabilité sont alors déterminées par les égalités des empreintes restreintes. En particulier,
\[
\chi_{\mathrm{ouvert}} \sim_{\Phi,P,\lambda_0} \chi_{\mathrm{ferme}}
\sim_{\Phi,P,\lambda_0} \chi_{\mathrm{vivant}}
\sim_{\Phi,P,\lambda_0} \chi_{\mathrm{mort}},
\]
tandis que
\[
\chi_{\mathrm{ferme}} \not\sim_{\Phi,P,\lambda_1} \chi_{\mathrm{ouvert}},
\qquad
\chi_{\mathrm{vivant}} \not\sim_{\Phi,P,\lambda_2} \chi_{\mathrm{mort}}.
\]

Le point essentiel est le suivant : le \emph{chat de Naszimger} est la classe observationnelle non résolue dont \(\chi_{\mathrm{ferme}}\) est un représentant. Sa fonction de seuil complète vérifie
\[
\Theta_P(\{\chi_{\mathrm{ferme}},\chi_{\mathrm{ouvert}}\})=
\Theta_P(\{\chi_{\mathrm{ferme}},\chi_{\mathrm{vivant}}\})=
\Theta_P(\{\chi_{\mathrm{ferme}},\chi_{\mathrm{mort}}\})
=\alpha_0,
\qquad
\alpha_0\in(\iota(\lambda_0),\iota(\lambda_1)),
\]
tandis que
\[
\Theta_P(\{\chi_{\mathrm{ouvert}},\chi_{\mathrm{vivant}}\})
=
\Theta_P(\{\chi_{\mathrm{ouvert}},\chi_{\mathrm{mort}}\})
=
\Theta_P(\{\chi_{\mathrm{vivant}},\chi_{\mathrm{mort}}\})
=\alpha_1,
\qquad
\alpha_1\in(\iota(\lambda_1),\iota(\lambda_2)).
\]
On obtient ainsi une filtration à trois niveaux : un niveau de fermeture, un niveau de séparation intermédiaire, puis un niveau final où les deux issues résolues deviennent distinguables.

\begin{proposition}[Fibre non résolue du chat de Naszimger]
À l'échelle basse \(\lambda_0\), la fibre de la classe du chat de Naszimger contient exactement les objets qui ne sont pas distingués par la seule sonde \(f_{\mathrm{boite}}\). En particulier,
\[
\pi_{\lambda_0}^{-1}\bigl([\chi_{\mathrm{ferme}}]_{\lambda_0}\bigr)
=
\{\chi_{\mathrm{ouvert}},\chi_{\mathrm{ferme}},\chi_{\mathrm{vivant}},\chi_{\mathrm{mort}}\}.
\]
Autrement dit, plusieurs objets donnent la même observation à basse résolution.
\end{proposition}

\begin{proof}
À l'échelle \(\lambda_0\), la seule coordonnée observée est \(f_{\mathrm{boite}}\). Or
\[
\Phi(P,\chi_{\mathrm{ouvert}},f_{\mathrm{boite}})
=
\Phi(P,\chi_{\mathrm{ferme}},f_{\mathrm{boite}})
=
\Phi(P,\chi_{\mathrm{vivant}},f_{\mathrm{boite}})
=
\Phi(P,\chi_{\mathrm{mort}},f_{\mathrm{boite}})
=1,
\]
tandis que
\[
\Phi(P,\chi_0,f_{\mathrm{boite}})=0.
\]
La fibre de la classe \([\chi_{\mathrm{ferme}}]_{\lambda_0}\) est donc exactement l'ensemble annoncé.
\end{proof}

\begin{proposition}[Reconstruction observationnelle d'une expérience de type Schrödinger]
Dans le cadre ci-dessus, la description complète de l'expérience est reconstruite par la fonction de seuil \(\Theta_P\) et la filtration \(\bigl(\sim_{\Phi,P,\lambda}\bigr)_{\lambda\in\Lambda}\). Plus précisément :
\begin{enumerate}[label=\arabic*.]
\item le représentant \(\chi_{\mathrm{ferme}}\) désigne la classe observationnelle non résolue à basse intensité ;
\item \(\chi_{\mathrm{ouvert}}\) joue le rôle d'état intermédiaire dans la filtration ;
\item les points \(\chi_{\mathrm{vivant}}\) et \(\chi_{\mathrm{mort}}\) représentent les deux issues résolues ;
\item la transition entre les régimes est entièrement codée par les seuils \(\alpha_0\) et \(\alpha_1\) ;
\item aucune donnée extérieure à l'empreinte n'est nécessaire pour reconstruire la structure logique de l'expérience.
\end{enumerate}
\end{proposition}

\begin{proof}
La proposition est une conséquence directe des définitions précédentes. Le rôle de \(\chi_{\mathrm{ferme}}\) est de servir de représentant de la classe non résolue à basse intensité. Les états \(\chi_{\mathrm{vivant}}\) et \(\chi_{\mathrm{mort}}\) sont séparés lorsque la famille de sondes contient \(f_{\mathrm{vivant}}\) et \(f_{\mathrm{mort}}\), tandis que \(\chi_{\mathrm{ouvert}}\) constitue un état intermédiaire distinct à haute résolution. La filtration des niveaux \(\lambda_0\prec\lambda_1\prec\lambda_2\) encode exactement le passage d'une information non résolue à une information progressivement résolue. La reconstruction est donc entièrement portée par l'empreinte et la fonction de seuil.
\end{proof}

\subsection{Lecture géométrique}

L'espace géométrique associé à cet exemple est l'image observationnelle réduite
\[
\mathcal L_{\lambda_2,P}=I_{\Phi,\lambda_2}^P(\mathcal X)
=
\bigl\{
(0,0,0,0,0),
(1,1,0,0,0),
(1,0,0,0,0),
(1,1,1,0,1),
(1,1,0,1,1)
\bigr\}
\subseteq \mathbb R^5.
\]
Dans cet espace, le chat de Naszimger n'est pas un postulat narratif, mais une classe d'orbite observationnelle : il est l'objet qui reste stable sous les sondes de basse résolution et qui se décompose seulement lorsque le seuil critique est franchi.

On peut représenter cette géométrie par le schéma suivant :
\[
[\chi_{\mathrm{ferme}}]_{\lambda_0}
\longrightarrow
\text{empreinte basse résolution}
\longrightarrow
\text{seuil critique}
\longrightarrow
\{\chi_{\mathrm{vivant}},\chi_{\mathrm{mort}}\}.
\]

Cette lecture est fidèle à la théorie des empreintes : l'expérience n'est pas d'abord une histoire physique imposée de l'extérieur, mais une hiérarchie de fibres qui devient visible au fur et à mesure que les sondes se raffinent.

\begin{remark}[Figure source]
Le manuscrit source contient ici une figure illustrative ; sa présence est documentée dans le corpus, mais l'image externe n'est pas embarquée dans cette reconstruction.
\end{remark}

\subsection{Diagramme de commutation}

La compatibilité entre quotient, projection et reconstruction s'écrit par le diagramme commutatif suivant :
\[
\begin{tikzpicture}[baseline=(current bounding box.center),node distance=22mm and 20mm,>=Stealth]
\node (X) {$\mathcal X$};
\node (R2) [right=of X] {$\mathbb R^{\mathcal F_{\lambda_2}}$};
\node (Q) [below=of X] {$\mathcal X/{\sim_{\Phi,P,\lambda_1}}$};
\node (R1) [below=of R2] {$\mathbb R^{\mathcal F_{\lambda_1}}$};

\draw[->] (X) -- node[above] {$I_{\Phi,\lambda_2}^P$} (R2);
\draw[->] (X) -- node[left] {$\pi_{\lambda_1}$} (Q);
\draw[->] (R2) -- node[right] {$r_{\lambda_2\to\lambda_1}$} (R1);
\draw[->] (Q) -- node[below] {$\widetilde I_{\lambda_1}$} (R1);
\end{tikzpicture}
\]
c'est-à-dire
\[
I_{\Phi,\lambda_1}^P
=
r_{\lambda_2\to\lambda_1}\circ I_{\Phi,\lambda_2}^P
=
\widetilde I_{\lambda_1}\circ \pi_{\lambda_1}.
\]

\begin{proposition}[Compatibilité des restrictions]
Le diagramme précédent commute par définition de la restriction des coordonnées. En effet, \(r_{\lambda_2\to\lambda_1}\) est l'application qui oublie les coordonnées correspondant à \(\mathcal F_{\lambda_2}\setminus\mathcal F_{\lambda_1}\), et \(\widetilde I_{\lambda_1}\) est l'application induite sur le quotient par \(\sim_{\Phi,P,\lambda_1}\).
\end{proposition}

\begin{proof}
Pour tout \(x\in\mathcal X\), la restriction de \(I_{\Phi,\lambda_2}^P(x)\) à \(\mathcal F_{\lambda_1}\) est exactement \(I_{\Phi,\lambda_1}^P(x)\) par définition des empreintes à l'échelle. De plus, si \(x\sim_{\Phi,P,\lambda_1}y\), alors \(I_{\Phi,\lambda_1}^P(x)=I_{\Phi,\lambda_1}^P(y)\), si bien que l'application \(\widetilde I_{\lambda_1}([x]_{\lambda_1})=I_{\Phi,\lambda_1}^P(x)\) est bien définie. Les deux compositions coïncident donc avec \(I_{\Phi,\lambda_1}^P\).
\end{proof}

Le lecteur voit ainsi simultanément la projection, le quotient, la filtration et la reconstruction.

\subsection{Interprétation interne à la théorie}

La reconstruction ci-dessus n'impose aucune lecture classique préalable. Elle montre seulement qu'une expérience à plusieurs issues peut être codée par une structure d'empreintes dont la géométrie est déterminée par des seuils critiques.

Autrement dit, dans la théorie des empreintes, l'opposition entre « état non résolu » et « état résolu » n'est pas un axiome extérieur : c'est une propriété de la filtration des fibres. Le chat de Naszimger joue alors le rôle d'objet test minimal pour cette séparation.

\begin{proposition}[Rigidité observationnelle à seuil fixé]
Soient deux expériences finies \((\mathcal X,\Phi,P,\Lambda,\iota)\) et \((\mathcal X',\Phi',P',\Lambda',\iota')\), munies de fonctions de seuil complètes \(\Theta_P\) et \(\Theta'_{P'}\). S'il existe une bijection \(\sigma:\mathcal X\to\mathcal X'\) telle que
\[
\Theta'_{P'}(\{\sigma(x),\sigma(y)\})=\Theta_P(\{x,y\})
\qquad\text{pour tous }x,y\in\mathcal X,
\]
alors les deux expériences sont isomorphes comme structures observationnelles.
\end{proposition}

\begin{proof}
La condition impose la conservation de toutes les valeurs de seuil, donc de tous les niveaux de séparation et de toutes les classes de la filtration. En effet, pour tout \(t\in\mathbb R_+\),
\[
x\equiv_t y
\quad\Longleftrightarrow\quad
\Theta_P(\{x,y\})\ge t
\quad\Longleftrightarrow\quad
\Theta'_{P'}(\{\sigma(x),\sigma(y)\})\ge t
\quad\Longleftrightarrow\quad
\sigma(x)\equiv'_t \sigma(y).
\]
Les partitions de niveau coïncident donc niveau par niveau, ainsi que l'arbre des classes et la mesure de résolution qui en dérivent. Les deux expériences sont donc la même structure observationnelle, transportée par \(\sigma\).
\end{proof}

\begin{corollary}[Principe de résolution par empreinte]
Si un objet observationnel est tel que deux issues deviennent séparables seulement à partir d'une certaine intensité d'observation, alors sa description minimale est donnée par la fonction de seuil associée à ses fibres, indépendamment de toute interprétation classique préalable.
\end{corollary}

\begin{proof}
C'est une reformulation de la proposition précédente en langage de seuils. L'issue non résolue correspond à une classe d'indiscernabilité commune à basse intensité, tandis que les deux issues séparées correspondent aux classes distinctes à intensité suffisante. La fonction de seuil encode exactement le point de transition.
\end{proof}

%














%
%%%%%%%%%%%%%%%%%%%%%%%%%%%%%%%%%%%%%%%%%%%%%%%%%%%%%%%%%%%%%%%%%%%%%%%%%%%%%%%%%%%%%%%%%%%%%%%%%%%%%%%%%%%%%%%%%%%%%%%%%%%%%%%%%%%%%%%%%%%%%%%%%%%%%%%%%%%%%%%%%%%%%%%%%%%%%%%%%%%%%%%%%%%%%%%%%experience de double résolutions

% Références internes au manuscrit global :
%   (filtration canonique, arbre étiqueté, enfants, ancêtre commun)
% - corollaire \emph{Non-recouvrement partiel}
% - théorème \emph{Représentation atomique de la résolution}
% - corollaire \emph{Invariance par reparamétrisation monotone}

\subsection*{Expérience à double résolutions : probabilité émergente et structure de type double fente}

\paragraph{Principe de l'exemple.}
On formalise ici une reconstruction observationnelle d'une structure observationnelle de type double fente dans le langage des empreintes. La donnée primitive n'est pas une probabilité attribuée à un objet isolé, mais une classe observationnelle qui se raffine lorsque le système de sondes devient plus discriminant. Dans ce cadre, la probabilité n'appartient pas à l'objet lui-même : elle appartient au couple formé par une classe d'objets, un système d'observation et un processus de résolution.

\begin{remark}
Dans tout ce qui suit, une superposition observationnelle n'est rien d'autre qu'une fibre non réduite à une échelle donnée. La lecture probabiliste n'apparaît qu'après avoir fixé une mesure de résolution compatible avec l'arbre de résolution. Aucune fonction d'onde n'est postulée ici.
\end{remark}

\subsection*{Dispositif observationnel}

On fixe un contexte \(P\in\mathcal P\) et l'on considère le système d'observation
\[
\mathcal O:=(P,\Phi,\Lambda,\mathcal F),
\]
où \(\Phi\) est le fonctionnel d'observation, \(\Lambda\) l'échelle de résolution, et \(\mathcal F\) la famille de sondes disponibles.

On considère l'ensemble fini
\[
\mathcal X=\{x_0,x_L,x_R,x_L^\sharp,x_R^\sharp\},
\]
où \(x_0\) désigne l'état de préparation, \(x_L\) et \(x_R\) les deux branches élémentaires, et \(x_L^\sharp,x_R^\sharp\) les deux branches après enregistrement mémoriel.

On fixe une famille de sondes
\[
\mathcal F=\{f_{\mathrm{source}},f_{\mathrm{traj}},f_{\mathrm{memoire}}\},
\]
et l'on définit le fonctionnel d'observation \(\Phi:\mathcal P\times\mathcal X\times\mathcal F\to\mathbb R\) par
\[
\begin{array}{c|ccc}
& f_{\mathrm{source}} & f_{\mathrm{traj}} & f_{\mathrm{memoire}}\\ \hline
x_0 & 0 & 0 & 0\\
x_L & 1 & 1 & 0\\
x_R & 1 & 0 & 0\\
x_L^\sharp & 1 & 1 & 1\\
x_R^\sharp & 1 & 0 & 1
\end{array}
\]
et l'on note
\[
I_\Phi^P(x)=\bigl(\Phi(P,x,f_{\mathrm{source}}),\Phi(P,x,f_{\mathrm{traj}}),\Phi(P,x,f_{\mathrm{memoire}})\bigr).
\]

On introduit pour tout \(\lambda\in\Lambda\) la projection quotient
\[
\pi_\lambda:\mathcal X\longrightarrow \mathcal X/{\sim_{\Phi,P,\lambda}},
\qquad
\pi_\lambda(x)=[x]_\lambda.
\]

On fixe trois échelles
\[
\Lambda=\{\lambda_0\prec\lambda_1\prec\lambda_2\},
\qquad
\mathcal F_{\lambda_0}=\{f_{\mathrm{source}}\},\quad
\mathcal F_{\lambda_1}=\{f_{\mathrm{source}},f_{\mathrm{traj}}\},\quad
\mathcal F_{\lambda_2}=\mathcal F.
\]

À basse résolution, la classe cohérente
\[
A=[x_L]_{\lambda_0}=[x_R]_{\lambda_0}=[x_L^\sharp]_{\lambda_0}=[x_R^\sharp]_{\lambda_0}
\]
est une superposition observationnelle au sens de la théorie : elle n'est pas un mélange ontologique, mais une fibre non réduite. On a en particulier
\[
\pi_{\lambda_0}^{-1}(A)=\{x_L,x_R,x_L^\sharp,x_R^\sharp\},
\qquad
\pi_{\lambda_0}^{-1}([x_0]_{\lambda_0})=\{x_0\}.
\]
À la première résolution, cette classe se décompose en deux branches,
\[
A_L=[x_L]_{\lambda_1}=[x_L^\sharp]_{\lambda_1},
\qquad
A_R=[x_R]_{\lambda_1}=[x_R^\sharp]_{\lambda_1}.
\]
À la seconde résolution, chaque branche se raffine encore en feuilles distinctes :
\[
[x_L]_{\lambda_2},\ [x_L^\sharp]_{\lambda_2},\ [x_R]_{\lambda_2},\ [x_R^\sharp]_{\lambda_2}.
\]

\subsection*{Mesure de résolution et probabilité émergente}

La théorie globale fournit une mesure atomique de la distribution des seuils. Cette mesure ne conserve pas, à elle seule, l'affectation des seuils aux branches. Toute probabilité de branche introduite ici doit donc déclarer une règle d'échantillonnage supplémentaire.

\begin{proposition}[Poids empirique des branches — construction conditionnelle]
Soit \(A\) un nœud d'une hiérarchie étiquetée finie et soit \(\operatorname{Ch}(A)\) son ensemble d'enfants. Sous la règle d'échantillonnage uniforme sur les objets de \(A\), on définit
\[
\omega_A(B):=\frac{|B|}{|A|},\qquad B\in\operatorname{Ch}(A),
\]
avec \(|A|>0\). Alors \(\sum_{B\in\operatorname{Ch}(A)}\omega_A(B)=1\). Cette construction est une règle probabiliste du protocole et non une conséquence de \(\mu_P\) seule.
\end{proposition}
\begin{proof}
Les enfants partitionnent \(A\), donc leurs cardinalités somment à \(|A|\). La normalisation donne la somme unitaire. La dernière assertion suit de la perte d'affectation des résumés scalaires.\qed\end{proof}

Sous la règle uniforme déclarée, on utilise \(\omega_A(B)=|B|/|A|\) comme poids de branche. La mesure \(\mu_P\) reste le résumé de la distribution des seuils et n'est pas identifiée à ce poids sans hypothèse supplémentaire.

\begin{definition}[Probabilité émergente de résolution — protocole conditionnel]
Soit \(A\) un nœud non vide d'une hiérarchie finie et \(B\in\operatorname{Ch}(A)\). Sous la règle d'échantillonnage uniforme sur les objets de \(A\), on définit
\[
\mathbb P_P(B\mid A):=\frac{|B|}{|A|}.
\]
Cette probabilité dépend de la règle d'échantillonnage choisie. Elle n'est pas déterminée par la mesure scalaire globale \(\mu_P\) seule.
\end{definition}

\begin{theorem}[La probabilité est une propriété du protocole de résolution]
Sous la règle d'échantillonnage uniforme, \(\mathbb P_P(\cdot\mid A)\) est une probabilité sur \(\operatorname{Ch}(A)\) et
\[
\sum_{B\in\operatorname{Ch}(A)}\mathbb P_P(B\mid A)=1.
\]
Pour une chaîne \(A=A_0\supset A_1\supset\cdots\supset A_n\), la probabilité de la trajectoire vaut
\[
\prod_{k=0}^{n-1}\mathbb P_P(A_{k+1}\mid A_k)=\frac{|A_n|}{|A_0|}.
\]
\end{theorem}
\begin{proof}
Les enfants partitionnent \(A\), donc leurs cardinalités somment à \(|A|\). Le produit télescope : \(\prod_{k=0}^{n-1}|A_{k+1}|/|A_k|=|A_n|/|A_0|\).\qed\end{proof}

\begin{corollary}[Invariance observationnelle du comptage uniforme]
Sous la règle uniforme, les probabilités de branche sont invariantes par isomorphisme de la hiérarchie étiquetée et par toute reparamétrisation strictement croissante qui ne modifie pas les classes aux niveaux considérés.
\end{corollary}
\begin{proof}
La probabilité \(|B|/|A|\) ne dépend que des cardinalités des blocs correspondants. Une isomorphie étiquetée conserve ces cardinalités ; une reparamétrisation strictement croissante qui conserve les coupes conserve les blocs.\qed\end{proof}

\begin{remark}
Le comptage uniforme n'intervient qu'à titre de cas particulier. Dans le modèle symétrique retenu ici, les quatre états de la classe cohérente sont répartis en deux branches de deux objets, puis chaque branche en deux feuilles. Sous la règle uniforme déclarée précédemment,
\[
\mathbb P_P(A_L\mid A)=\mathbb P_P(A_R\mid A)=\frac12,
\]
et, au niveau terminal,
\[
\mathbb P_P([x_L]_{\lambda_2}\mid A_L)=\mathbb P_P([x_L^\sharp]_{\lambda_2}\mid A_L)=\frac12,
\]
avec la même formule pour la branche droite. Ces valeurs proviennent du cardinal des blocs et de la règle d'échantillonnage, non de la mesure scalaire des seuils.
\end{remark}

\begin{proposition}[Reconstruction observationnelle d'une structure de type double fente]
Dans le modèle ci-dessus, la classe cohérente \(A\) se décompose en deux branches de même probabilité sous l'échantillonnage uniforme, puis chaque branche se raffine en deux feuilles de même probabilité. Cette construction reproduit, dans le langage des empreintes, la logique d'une structure de type double fente : une classe cohérente, une première séparation en deux chemins, puis un raffinement secondaire des traces.

Le terme « double fente » reste une analogie structurelle interne au protocole ; aucune conclusion physique sur un phénomène quantique n'est déduite de cette reconstruction.
\end{proposition}

\begin{proof}
La classe \(A\) contient quatre objets et chacun des deux enfants contient deux objets, d'où les probabilités \(2/4=1/2\). Chaque enfant terminal contient deux feuilles d'un objet chacune, d'où \(1/2\) sous l'échantillonnage uniforme local.\qed\end{proof}

\section{Perspective : un modèle de transport observationnel fondé sur les empreintes}

\subsection{Idée directrice}

La théorie des empreintes montre qu'un objet n'est pas saisi d'abord par la totalité de sa microdescription, mais par ce qui reste invariant sous le système d'observation choisi. Autrement dit, la donnée primitive n'est pas l'état complet, mais l'empreinte observationnelle.

Le point mathématique essentiel est alors le suivant :

\begin{quote}
une opération est observationnellement représentable si, et seulement si, elle préserve les fibres observables ; dans ce cas, elle se factorise de manière unique sur l'espace des empreintes.
\end{quote}

Cette section formalise cette descente, puis montre comment elle se prolonge aux échelles, aux seuils et à la géométrie de résolution déjà construits dans le manuscrit global.

\begin{remark}
Le résultat ci-dessous ne prétend pas que toute opération physique ou informatique se transporte automatiquement sur les empreintes. Il caractérise exactement la classe des opérations qui préservent les fibres observables, donc représentables au niveau observationnel. La supériorité matérielle éventuelle d'une architecture concrète n'est pas démontrée ici.
\end{remark}

\subsection{Cadre de représentation}

Fixons un contexte \(P\in\mathcal P\) et une famille finie de sondes \(\mathcal G\subseteq\mathcal F_\lambda\). On considère l'empreinte réduite
\[
I_{\Phi,\mathcal G}^P:\mathcal X\longrightarrow \mathbb R^{\mathcal G},
\qquad
\mathcal L_{\mathcal G,P}:=I_{\Phi,\mathcal G}^P(\mathcal X).
\]
L'ensemble \(\mathcal L_{\mathcal G,P}\) est l'image observationnelle réduite canonique des objets à travers la famille \(\mathcal G\).

Par le \emph{Théorème de réduction canonique des sondes} du manuscrit global, on peut remplacer \(\mathcal G\) par sa fermeture observationnelle
\[
\mathcal G^\sharp:=\operatorname{Cl}_{P,\lambda}(\mathcal G),
\]
sans modifier les fibres observables.

\begin{proposition}[Réduction canonique préalable des sondes]
\label{prop:reduction-prealable}
Pour toute famille finie \(\mathcal G\subseteq\mathcal F_\lambda\), les empreintes
\[
I_{\Phi,\mathcal G}^P
\qquad\text{et}\qquad
I_{\Phi,\operatorname{Cl}_{P,\lambda}(\mathcal G)}^P
\]
ont exactement les mêmes fibres.
Autrement dit, la réduction canonique des sondes ne change pas le contenu classificatoire.
\end{proposition}

\begin{proof}
C'est exactement le \emph{Théorème de réduction canonique des sondes} du manuscrit global, appliqué au contexte \(P\), à l'échelle \(\lambda\) et à la famille \(\mathcal G\).
\end{proof}

\subsection{Caractérisation de la représentation observationnelle}

\begin{definition}[Opération observationnellement représentable]
Une application \(F:\mathcal X\to\mathcal X\) est dite \(\mathcal G\)-compatible si
\[
I_{\Phi,\mathcal G}^P(x)=I_{\Phi,\mathcal G}^P(y)
\quad\Longrightarrow\quad
I_{\Phi,\mathcal G}^P(F(x))=I_{\Phi,\mathcal G}^P(F(y))
\]
pour tous \(x,y\in\mathcal X\).

On dit alors que \(F\) est observationnellement représentable à travers \(\mathcal G\).
\end{definition}

\begin{theorem}[Caractérisation de la représentation observationnelle par la propriété universelle du quotient]
\label{thm:representation-observationnelle}
Fixons \(P\in\mathcal P\) et une famille réduite de sondes \(\mathcal G\subseteq\mathcal F_\lambda\). Soit \(F:\mathcal X\to\mathcal X\) une application. Les assertions suivantes sont équivalentes :
\begin{enumerate}[label=\textbf{(\arabic*)}]
\item \(F\) est \(\mathcal G\)-compatible, c'est-à-dire qu'elle préserve les fibres observables de \(I_{\Phi,\mathcal G}^P\).
\item Il existe une unique application
\[
\mathcal C_F:\mathcal L_{\mathcal G,P}\longrightarrow \mathcal L_{\mathcal G,P}
\]
telle que
\[
I_{\Phi,\mathcal G}^P\circ F=\mathcal C_F\circ I_{\Phi,\mathcal G}^P.
\]
\item Il existe une unique application
\[
\widehat{\mathcal C}_F:\mathcal X/{\sim_{\Phi,P,\mathcal G}}\longrightarrow \mathcal X/{\sim_{\Phi,P,\mathcal G}}
\]
telle que
\[
\pi_{\mathcal G}\circ F=\widehat{\mathcal C}_F\circ \pi_{\mathcal G}.
\]
\end{enumerate}

En particulier, toute opération observationnellement représentable descend de manière unique sur l'espace des empreintes et sur le quotient observationnel.
\end{theorem}

\begin{proof}
Le \emph{Théorème de factorisation de l'empreinte} établit la correspondance entre les fibres de \(I_{\Phi,\mathcal G}^P\) et le quotient observationnel associé. La propriété universelle du quotient donne alors l'équivalence entre la compatibilité sur les fibres et l'existence d'une factorisation unique.

Si \(F\) est \(\mathcal G\)-compatible, on définit
\[
\mathcal C_F\bigl(I_{\Phi,\mathcal G}^P(x)\bigr):=I_{\Phi,\mathcal G}^P(F(x)).
\]
La définition est indépendante du représentant \(x\) précisément parce que \(F\) préserve les fibres observables de \(I_{\Phi,\mathcal G}^P\). Cela donne \((1)\Rightarrow(2)\). L'implication \((2)\Rightarrow(3)\) suit du passage au quotient par les fibres, déjà établi par le théorème de factorisation global.

Réciproquement, si une factorisation existe, alors \(I_{\Phi,\mathcal G}^P(x)=I_{\Phi,\mathcal G}^P(y)\) implique
\[
I_{\Phi,\mathcal G}^P(F(x))
=
\mathcal C_F(I_{\Phi,\mathcal G}^P(x))
=
\mathcal C_F(I_{\Phi,\mathcal G}^P(y))
=
I_{\Phi,\mathcal G}^P(F(y)),
\]
donc \(F\) est \(\mathcal G\)-compatible. L'unicité est immédiate car \(\mathcal L_{\mathcal G,P}\) est l'image de \(I_{\Phi,\mathcal G}^P\).
\end{proof}

\begin{remark}
Le théorème précédent exprime exactement la propriété universelle du quotient observationnel : les applications qui préservent les fibres observables sont précisément celles qui se factorisent de manière unique à travers \(\mathcal X/{\sim_{\Phi,P,\mathcal G}}\) ou, de façon équivalente, à travers \(\mathcal L_{\mathcal G,P}\).
\end{remark}

\begin{corollary}[Critère de non-descente]
Si une application \(F:\mathcal X\to\mathcal X\) sépare au moins une fibre observable de \(I_{\Phi,\mathcal G}^P\), alors elle ne descend ni à \(\mathcal L_{\mathcal G,P}\) ni au quotient \(\mathcal X/{\sim_{\Phi,P,\mathcal G}}\).
\end{corollary}

\begin{proof}
C'est la contraposée immédiate du théorème précédent.
\end{proof}

\begin{proposition}[Structure algébrique des opérations compatibles]
\label{prop:monoid-operations}
L'ensemble
\[
\mathfrak M_{\mathcal G,P}:=\{F:\mathcal X\to\mathcal X\mid F\ \text{est \(\mathcal G\)-compatible}\}
\]
est un monoïde pour la composition, dont l'élément neutre est \(\operatorname{id}_{\mathcal X}\).
De plus, l'application
\[
F\longmapsto \mathcal C_F
\]
est un morphisme de monoïdes de \((\mathfrak M_{\mathcal G,P},\circ,\operatorname{id}_{\mathcal X})\) vers \((\operatorname{End}(\mathcal L_{\mathcal G,P}),\circ,\operatorname{id}_{\mathcal L_{\mathcal G,P}})\).
\end{proposition}

\begin{proof}
Si \(F_1\) et \(F_2\) sont \(\mathcal G\)-compatibles, alors \(F_2\circ F_1\) l'est encore, car la préservation des fibres observables est conservée par composition. Par le théorème précédent,
\[
I_{\Phi,\mathcal G}^P\circ F_i=\mathcal C_{F_i}\circ I_{\Phi,\mathcal G}^P
\qquad (i=1,2).
\]
Donc
\[
I_{\Phi,\mathcal G}^P\circ(F_2\circ F_1)
=
\mathcal C_{F_2}\circ\mathcal C_{F_1}\circ I_{\Phi,\mathcal G}^P.
\]
Par unicité de la factorisation, on obtient
\[
\mathcal C_{F_2\circ F_1}=\mathcal C_{F_2}\circ \mathcal C_{F_1}.
\]
L'identité se factorise trivialement par elle-même. Ainsi, on obtient bien un morphisme de monoïdes.
\end{proof}

\begin{corollary}[Suite finie d'opérations observationnelles]
Si
\[
W=F_n\circ F_{n-1}\circ\cdots\circ F_1
\]
est une suite finie d'opérations \(\mathcal G\)-compatibles, alors il existe une unique application induite
\[
\mathcal C_W=\mathcal C_{F_n}\circ\cdots\circ\mathcal C_{F_1}
\]
telle que
\[
I_{\Phi,\mathcal G}^P\circ W=\mathcal C_W\circ I_{\Phi,\mathcal G}^P.
\]
Ainsi, toute suite finie d'opérations compatibles admet une représentation entièrement définie sur l'espace des empreintes.
\end{corollary}

\begin{proof}
C'est une conséquence immédiate de la proposition précédente par récurrence sur le nombre de compositions.
\end{proof}

\begin{corollary}[Représentation observationnelle d'une fonction de coût]
Soit \(\Theta\) un espace de paramètres et soit
\[
L:\Theta\times\mathcal X\longrightarrow \mathbb R
\]
une fonction de coût telle que, pour tout \(\theta\in\Theta\),
\[
I_{\Phi,\mathcal G}^P(x)=I_{\Phi,\mathcal G}^P(y)
\Longrightarrow
L(\theta,x)=L(\theta,y).
\]
Alors il existe une unique fonction
\[
\widetilde L:\Theta\times\mathcal L_{\mathcal G,P}\longrightarrow \mathbb R
\]
vérifiant
\[
L(\theta,x)=\widetilde L\bigl(\theta,I_{\Phi,\mathcal G}^P(x)\bigr).
\]
En particulier, toute fonction de coût constante sur les fibres admet une représentation unique sur l'espace des empreintes.
\end{corollary}

\begin{proof}
Pour chaque \(\theta\), l'application \(x\mapsto L(\theta,x)\) est constante sur les fibres de \(I_{\Phi,\mathcal G}^P\). Le \emph{Théorème de factorisation de l'empreinte} donne donc une factorisation unique à travers \(\mathcal L_{\mathcal G,P}\). En faisant varier \(\theta\), on obtient \(\widetilde L\).
\end{proof}

\subsection{Compatibilité avec les échelles et la résolution}

On passe maintenant au cadre multi-échelle du manuscrit global. Pour \(\lambda\preceq\mu\), on dispose de la restriction
\[
r_{\mu\to\lambda}:\mathbb R^{\mathcal F_\mu}\longrightarrow \mathbb R^{\mathcal F_\lambda},
\]
qui oublie les coordonnées de \(\mathcal F_\mu\setminus\mathcal F_\lambda\). Par définition,
\[
I_{\Phi,\lambda}^P=r_{\mu\to\lambda}\circ I_{\Phi,\mu}^P.
\]

\begin{proposition}[Compatibilité inter-échelles]
\label{prop:inter-echelles}
Supposons qu'une opération \(F:\mathcal X\to\mathcal X\) soit \(\mathcal F_\lambda\)-compatible et \(\mathcal F_\mu\)-compatible. Notons \(\mathcal C_{F,\lambda}\) et \(\mathcal C_{F,\mu}\) les applications induites respectivement sur \(\mathcal L_{\lambda,P}\) et \(\mathcal L_{\mu,P}\). Alors
\[
r_{\mu\to\lambda}\circ \mathcal C_{F,\mu}
=
\mathcal C_{F,\lambda}\circ r_{\mu\to\lambda}
\]
sur \(\mathcal L_{\mu,P}\).

Autrement dit, la représentation observée à haute résolution est cohérente avec sa restriction à plus basse résolution.
\end{proposition}

\begin{proof}
Pour tout \(x\in\mathcal X\),
\[
r_{\mu\to\lambda}\bigl(\mathcal C_{F,\mu}(I_{\Phi,\mu}^P(x))\bigr)
=
r_{\mu\to\lambda}\bigl(I_{\Phi,\mu}^P(F(x))\bigr)
=
I_{\Phi,\lambda}^P(F(x)).
\]
D'un autre côté,
\[
\mathcal C_{F,\lambda}\bigl(r_{\mu\to\lambda}(I_{\Phi,\mu}^P(x))\bigr)
=
\mathcal C_{F,\lambda}(I_{\Phi,\lambda}^P(x))
=
I_{\Phi,\lambda}^P(F(x)).
\]
Les deux compositions coïncident donc sur l'image, et l'unicité de la factorisation donne l'égalité annoncée.
\end{proof}

\begin{center}
\[
\begin{tikzcd}[column sep=large,row sep=large]
\mathcal L_{\mu,P} \arrow[r,"\mathcal C_{F,\mu}"] \arrow[d,"r_{\mu\to\lambda}"'] &
\mathcal L_{\mu,P} \arrow[d,"r_{\mu\to\lambda}"] \\
\mathcal L_{\lambda,P} \arrow[r,"\mathcal C_{F,\lambda}"'] &
\mathcal L_{\lambda,P}
\end{tikzcd}
\]
\end{center}

\begin{proposition}[Compatibilité avec la géométrie des seuils]
\label{prop:seuils-calcul}
Supposons maintenant que \(F:\mathcal X\to\mathcal X\) soit bijective et qu'il préserve la fonction de seuil complète :
\[
\Theta_P(\{F(x),F(y)\})=\Theta_P(\{x,y\})
\qquad\text{pour tous }x,y\in\mathcal X.
\]
Alors \(F\) préserve toutes les relations de superniveau, donc toutes les partitions \(\equiv_t\), et induit un automorphisme de l'arbre de résolution \(\mathcal T_P\).

De plus, la mesure empirique \(\nu_P\) et la courbe de résolution \(\widetilde{\mathcal R}_P\) sont invariantes sous \(F\).
\end{proposition}

\begin{proof}
Par le \emph{Théorème structural des seuils}, on a
\[
x\equiv_t y
\Longleftrightarrow
\Theta_P(\{x,y\})\ge t.
\]
L'hypothèse impose donc
\[
\Theta_P(\{F(x),F(y)\})\ge t
\Longleftrightarrow
\Theta_P(\{x,y\})\ge t,
\]
soit
\[
x\equiv_t y \Longleftrightarrow F(x)\equiv_t F(y).
\]
Ainsi, chaque partition de niveau est préservée, donc l'arbre complet de classes est préservé.

Comme la mesure \(\nu_P\) est construite à partir des valeurs de \(\Theta_P\) avec multiplicité, et que \(F\) en préserve les valeurs, \(\nu_P\) est invariant. Par le \emph{Théorème de représentation atomique de la résolution} et le \emph{Corollaire d'invariance par reparamétrisation monotone} du manuscrit global, la courbe \(\widetilde{\mathcal R}_P\) est alors également invariante.
\end{proof}

\begin{remark}
Le résultat précédent relie directement la représentation observationnelle à la géométrie des seuils du manuscrit global. Une application qui préserve la fonction de seuil complète (et donc les partitions de résolution) peut alors être vue comme une action sur l'arbre de résolution, et non seulement comme une application sur un ensemble abstrait.
\end{remark}

\subsection{Exemple minimal de descente observationnelle}

\begin{proposition}[Exemple élémentaire]
Soit \(\mathcal X=\mathbb R^2\), \(\mathcal F=\{f_1,f_2,f_3\}\), et définissons
\[
\Phi(P,(u,v),f_1)=u,\qquad
\Phi(P,(u,v),f_2)=v,\qquad
\Phi(P,(u,v),f_3)=u+v.
\]
Considérons la famille \(\mathcal G=\{f_1\}\). L'opération
\[
F(u,v)=(u,\psi(u,v))
\]
est \(\mathcal G\)-compatible pour toute fonction \(\psi\). L'application induite sur l'image observationnelle réduite \(\mathcal L_{\mathcal G,P}\simeq\mathbb R\) est alors l'identité :
\[
\mathcal C_F(u)=u.
\]
Ainsi, une évolution microscopique non triviale peut être invisible au niveau observationnel.
\end{proposition}

\begin{proof}
Avec \(\mathcal G=\{f_1\}\), deux points \((u,v)\) et \((u,v')\) ont la même empreinte dès qu'ils partagent la même première coordonnée. Or
\[
F(u,v)=(u,\psi(u,v))
\]
ne change pas la première coordonnée, donc l'empreinte observée reste la même après application de \(F\). La descente sur \(\mathcal L_{\mathcal G,P}\simeq\mathbb R\) est donc bien définie et vaut l'identité.
\end{proof}

\begin{proposition}[Exemple de non-descente]
Soit \(\mathcal G=\{f_1\}\) dans le cadre précédent et définissons
\[
G(u,v)=(u+v,v).
\]
Alors \(G\) n'est pas \(\mathcal G\)-compatible.
\end{proposition}

\begin{proof}
Les points \((0,0)\) et \((0,1)\) ont la même empreinte relativement à \(\mathcal G\), puisqu'ils ont la même première coordonnée. En revanche,
\[
G(0,0)=(0,0)
\qquad\text{et}\qquad
G(0,1)=(1,1),
\]
donc
\[
I_{\Phi,\mathcal G}^P(G(0,0))\neq I_{\Phi,\mathcal G}^P(G(0,1)).
\]
L'opération \(G\) sépare donc une fibre observationnelle et ne descend pas sur l'espace des empreintes.
\end{proof}

\subsection{Architecture proposée}

Le schéma de transport observationnel obtenu ici peut être résumé par la chaîne
\[
\mathcal X
\;\xrightarrow{\ I_{\Phi,\mathcal G}^P\ }\;
\mathcal L_{\mathcal G,P}
\;\xrightarrow{\ \mathcal C\ }\;
\mathcal L_{\mathcal G,P},
\]
avec reconstruction éventuelle d'un représentant dans \(\mathcal X\) seulement si une interprétation microscopique est nécessaire.

\begin{remark}
Cette architecture ne démontre pas qu'un ordinateur devient automatiquement plus rapide. Elle démontre quelque chose de plus précis : toute opération qui respecte les fibres observables descend de manière unique sur les empreintes, où elle se compose, s'itère et se compare sans reconstruction systématique de l'état complet. C'est un résultat mathématique de descente, pas une promesse matérielle.
\end{remark}

\subsection{Conséquences théoriques}

\begin{proposition}[Stabilité observationnelle de la représentation]
Si deux états \(x,y\in\mathcal X\) ont la même empreinte relativement à \(\mathcal G\), alors toute suite finie d'opérations \(\mathcal G\)-compatibles produit la même empreinte finale à partir de \(x\) ou de \(y\).
\end{proposition}

\begin{proof}
Soit \(W=F_n\circ\cdots\circ F_1\) une suite finie d'opérations \(\mathcal G\)-compatibles. Par le corollaire sur les suites finies, il existe une application induite \(\mathcal C_W\) telle que
\[
I_{\Phi,\mathcal G}^P\circ W=\mathcal C_W\circ I_{\Phi,\mathcal G}^P.
\]
Ainsi, si \(I_{\Phi,\mathcal G}^P(x)=I_{\Phi,\mathcal G}^P(y)\), alors
\[
I_{\Phi,\mathcal G}^P(W(x))
=
\mathcal C_W(I_{\Phi,\mathcal G}^P(x))
=
\mathcal C_W(I_{\Phi,\mathcal G}^P(y))
=
I_{\Phi,\mathcal G}^P(W(y)).
\]
La sortie observationnelle finale est donc identique.
\end{proof}

\begin{corollary}[Compatibilité compositionnelle]
Les opérations \(\mathcal G\)-compatibles forment un monoïde, et leur action sur \(\mathcal L_{\mathcal G,P}\) est compatible avec la composition. Si les opérations sont bijectives et que leurs inverses sont encore compatibles, on obtient un groupe d'automorphismes de l'image observationnelle réduite.
\end{corollary}

\begin{proof}
C'est une reformulation de la proposition the registre des opérations représentables. La partie groupe suit de l'existence d'inverses compatibles.
\end{proof}

\subsection{Statut scientifique}

\chapter{Reprise détaillée des développements anciens et de leurs limites}
\section{Empreinte observationnelle universelle à contexte fixé}
Un développement ancien de \texttt{D\_Dev} définit, sans ajouter de primitive externe, l'empreinte complète
\[
I_{\Phi,\mathcal F}^{P}(x)(f):=\Phi(P,x,f),\qquad f\in\mathcal F,
\]
et, si l'on veut agréger simultanément les contextes déjà présents dans la théorie,
\[
I_{\Phi}^{\mathrm{univ}}(x)(P,f):=\Phi(P,x,f).
\]
Ces constructions sont des extensions directes de la définition de l'empreinte lorsque les familles considérées sont autorisées à être non finies. Elles restent des objets observationnels internes et ne sont pas identifiées à une sémantique externe.

\begin{proposition}[Indiscernabilité universelle]
La relation
\[
x\equiv_{\Phi}^{\mathrm{univ}}y
\Longleftrightarrow
\Phi(P,x,f)=\Phi(P,y,f)\quad\text{pour tous }P\in\mathcal P,\ f\in\mathcal F
\]
est une relation d'équivalence et implique toute indiscernabilité obtenue par restriction à un contexte et à une sous-famille de sondes.
\end{proposition}
\begin{proof}
Elle est l'égalité de deux applications. La restriction à une sous-famille ne peut qu'oublier des coordonnées, donc l'égalité universelle implique l'égalité restreinte.\qed\end{proof}

\begin{theorem}[Factorisation universelle d'une donnée observationnelle]
Pour une application \(T:\mathcal X\to Y\), les trois propriétés suivantes sont équivalentes : \(T\) est constante sur les fibres de \(I_{\Phi}^{\mathrm{univ}}\) ; \(T\) se factorise de manière unique par l'image de cette empreinte ; \(T\) se factorise de manière unique par son quotient universel.
\end{theorem}
\begin{proof}
C'est la factorisation élémentaire par les fibres : la valeur sur une fibre est bien définie précisément lorsque \(T\) y est constante. L'unicité vient du fait que toute valeur de l'image possède un représentant.\qed\end{proof}

\section{Correction d'un énoncé de maximalité}
Un texte source qualifiait l'indiscernabilité universelle de « plus fine ». Cette formulation est inversée si l'ordre des relations est l'inclusion : la relation universelle est la \emph{plus grande} relation d'équivalence dont chaque paire liée est indiscernable pour toutes les observations. Une relation plus fine au sens de moins d'identifications peut toujours être choisie comme sous-relation, mais elle ne représente pas davantage d'information observationnelle.

\begin{proposition}[Maximalité observationnelle correctement orientée]
Si \(\approx\) est une relation telle que
\[
x\approx y\Longrightarrow x\equiv_{\Phi}^{\mathrm{univ}}y,
\]
alors \(\approx\subseteq\equiv_{\Phi}^{\mathrm{univ}}\). Ainsi \(\equiv_{\Phi}^{\mathrm{univ}}\) est maximale par inclusion parmi les relations dont les paires sont rendues indiscernables par la totalité des observations.
\end{proposition}
\begin{proof}
L'inclusion est exactement l'implication de l'hypothèse.\qed\end{proof}

\section{Signatures, profils et stabilisation observationnelle}
Les \texttt{D\_Base} anciens sont récupérés sous les hypothèses explicites de séparation des profils. Une signature de chaîne \(\Sigma(\Gamma)=(\mathcal A([\rho_i]))_i\) détermine les classes de la chaîne dès que l'application de profil \(\mathcal A\) est injective. Pour une chaîne infinie, tout invariant qui est constant sur les classes de stabilisation observationnelle se factorise de manière unique par cette stabilisation. Enfin, un morphisme observationnel qui préserve l'ordre envoie les chaînes sur des chaînes et les profils par inclusion.

Ces résultats sont conservés comme développements internes conditionnels : ils ne supposent ni métrique externe, ni temps physique, ni sémantique ajoutée.

\chapter{Origine documentaire, audit et traçabilité}

La théorie consolidée repose sur le corpus source et non sur les productions des ordres antérieurs. Les registres fournis avec ce volume relient les unités de source aux destinations logiques du livre. Les chapitres historiques des ordres antérieurs ont été déplacés hors du corps principal afin de préserver l'information sans imposer une organisation chronologique au lecteur.

\section{Corpus}
Le fichier source fondamental possède l'empreinte SHA-256 :
\[
\texttt{c88425d2d46c2fb27cc448b070c850909ebf896865ede3d560b88459d71530b0}.
\]
Le corpus contient 255 fichiers source, répartis dans $67$ fichiers v0, $26$ fichiers v1 et $162$ fichiers v2, auxquels s'ajoute $RG.txt$ comme entrée documentaire de référence. Le registre formel O4 recense 6262 environnements structurés.

\section{Statuts de preuve}
Les registres distinguent les propositions démontrées, les résultats conditionnels, les conjectures, les observations empiriques et les diagnostics documentaires. Les résultats O49--O51 qui dépendent de code absent de l'archive courante sont explicitement marqués comme résultats rapportés, non comme reproductions indépendantes.


\chapter{Historique scientifique condensé}

L'historique des ordres reste consultable comme chronologie documentaire. Il n'est pas utilisé comme principe d'organisation du livre. Les étapes anciennes sont scientifiquement retenues lorsqu'elles introduisent une construction, une démonstration, un contre-exemple ou une limite qui demeure pertinente après correction. Les suppressions sont traitées dans le registre des suppressions ; aucune branche n'est supprimée dans le seul but de raccourcir le livre.


\backmatter

\chapter*{Conclusion générale}

La théorie des empreintes se présente ici comme une architecture où l'observation produit une empreinte, l'empreinte produit des fibres et des quotients, les familles de niveaux produisent des seuils et une géométrie hiérarchique, et les extensions recherchent des transformations, des contraintes et des réalisations dont la provenance doit rester explicite. La correction des seuils n'a pas supprimé la géométrie : elle a séparé le squelette $\Theta_P$ de son bord $b_P$. La correction de la mesure n'a pas supprimé les probabilités : elle les a replacées sous une règle d'échantillonnage. La correction de la canonicité n'a pas supprimé la représentation : elle a limité la canonicité aux quotients et factorisations effectivement définis. Enfin, l'imagination et les contraintes sont conservées comme constructions propres de la théorie, avec leurs frontières de réalisation.

\end{document}